% Les marchés publics (corruption, contrebande, falsification des produits) — Allemand Tle A — Cours signé OnBuch+
% Programme officiel MINESEC. Compiler avec : tectonic AL08-les-marches-publics-corruption-contrebande-falsification-des.tex
\newcommand{\DOCMATIERE}{Allemand}
\newcommand{\DOCNIVEAU}{Tle A}
\newcommand{\DOCMODULE}{Chapitre 2 : Vie économique}
\newcommand{\DOCLECON}{AL08}
\newcommand{\DOCTITRE}{Les marchés publics (corruption, contrebande, falsification des produits)}
% OnBuch+ — préambule « cours signés » ENRICHI (compilé avec tectonic / XeLaTeX)
% Système visuel « Pop » d'OnBuch+ : fond crème (jamais blanc), bordures encre
% épaisses, ombres « dures » décalées, titres Archivo Black en capitales.
% Version MATHÉMATIQUES : graphes (pgfplots/tikz), boîtes pédagogiques
% (définition, propriété, méthode, attention, le savais-tu, exercice résolu,
% exercice, TP, bilan), page de garde et signature OnBuch+ ancrée en bas.
%
% Macros attendues dans main.tex AVANT \input{preamble} :
%   \DOCMATIERE  \DOCNIVEAU  \DOCMODULE  \DOCLECON (numéro)  \DOCTITRE
\documentclass[11pt]{article}

\usepackage{fontspec}
\usepackage[ngerman,french]{babel} % français = langue principale ; l'allemand via \all{..} / textebox
\usepackage[a4paper,margin=2cm,top=3.4cm,bottom=2.8cm,headheight=1.4cm,footskip=1.3cm]{geometry}
\usepackage{amsmath,amssymb,amsfonts}
\usepackage{mathtools} % \xmapsto, \coloneqq... souvent utilisés par les modèles
\usepackage{cancel} % simplifications barrées (bilans de forces)
% \abs{z} (module, valeur absolue) et \norm{u} (norme), utilisés par les modèles
\providecommand{\abs}{}\let\abs\relax\DeclarePairedDelimiter{\abs}{\lvert}{\rvert}
\providecommand{\norm}{}\let\norm\relax\DeclarePairedDelimiter{\norm}{\lVert}{\rVert}
\usepackage[table]{xcolor} % [table] : fournit aussi \rowcolors (alternance de lignes)
\usepackage{pagecolor}
\usepackage{tikz}
\usetikzlibrary{arrows.meta,positioning,calc,shapes.geometric,shapes.misc,shapes.arrows,decorations.pathmorphing,decorations.pathreplacing,decorations.markings,patterns,shadows,mindmap,trees,fit,backgrounds,matrix,intersections,angles,quotes}
\usepackage{pgfplots}
\usepackage[version=4]{mhchem} % équations nucléaires \ce{^{235}_{92}U}
\usepackage[european,siunitx]{circuitikz} % schémas électriques
\pgfplotsset{compat=1.18}
\usepgfplotslibrary{fillbetween,groupplots} % aires entre courbes (calcul intégral)
\usepackage{enumitem}
\usepackage{fancyhdr}
% [most] charge toutes les bibliothèques tcolorbox (listings, minted, raster,
% documentation...) et gonfle inutilement la mémoire TeX sur un document long
% (~300 boîtes) : on ne charge que ce qui est réellement utilisé.
\usepackage[skins,breakable]{tcolorbox}
\usepackage{graphicx}
\usepackage{colortbl}
\usepackage{booktabs}
\usepackage{tabularx}
\usepackage{diagbox} % case d'en-tête diagonale des tableaux à double entrée
% Colonnes extensibles centrée / gauche / droite (C, L, R), souvent utilisées par les modèles
\newcolumntype{C}{>{\centering\arraybackslash}X}
\newcolumntype{L}{>{\raggedright\arraybackslash}X}
\newcolumntype{R}{>{\raggedleft\arraybackslash}X}
\usepackage{array}
\usepackage{multicol}
\usepackage{multirow}
\usepackage{siunitx}
% parse-numbers=false : accepte aussi \SI{2,5 \times 10^{-3}}{...}
\sisetup{locale=FR,output-decimal-marker={,},group-minimum-digits=4,parse-numbers=false}
\DeclareSIUnit\ohm{\ensuremath{\symup{\Omega}}} % U+2126 absent de la police
\DeclareSIUnit\division{div} % réglages d'oscilloscope (ms/div, V/div)
\DeclareSIUnit\tr{tr} % tours (vitesse de rotation, ex. \SI{1500}{\tr\per\minute})
\DeclareSIUnit\molar{M} % concentration molaire (ex. \SI{3}{\molar}, \SI{1}{\micro\molar})
\DeclareSIUnit\an{an} % année (le modèle confond parfois avec \year, un compteur TeX) — ex. \SI{5}{\kilo\meter\per\an}
\DeclareSIUnit\FCFA{FCFA} % franc CFA (ex. \SI{500}{\FCFA\per\kilo\gram})
\DeclareSIUnit\degreeBrix{\textdegree Brix} % teneur en sucre d'un sirop/jus (réfractométrie)
\DeclareSIUnit\month{mois} % le modèle utilise parfois \month, un compteur TeX, comme unité de durée
\usepackage{qrcode}
% \degree : commande fréquemment utilisée par les modèles pour un angle ou une
% température en dehors de \SI{}{} (non fournie par les paquets chargés).
\providecommand{\degree}{^{\circ}}
% \qed (fin de démonstration), utilisé par les modèles sans amsthm
\providecommand{\qed}{\hfill$\square$}
% \barre : double barre des valeurs interdites dans un tableau de variations (tkz-tab)
\providecommand{\barre}{\ensuremath{\|}}
% environnement proof (amsthm non chargé) utilisé par les modèles
\ifdefined\proof\else\newenvironment{proof}[1][Démonstration]{\par\noindent\textit{#1.}\ }{\qed\par}\fi
\usepackage{lastpage}
\usepackage{hyperref}
\hypersetup{hidelinks}

% ============================================================
% Palette « Pop » OnBuch+ — mêmes hex que la classe P (plus_ui.dart)
% ============================================================
\definecolor{poporange}{HTML}{F59321}
\definecolor{popdark}{HTML}{E07A0C}
\definecolor{popcream}{HTML}{FFF6E8}
\definecolor{popink}{HTML}{1C1714}
\definecolor{poppurple}{HTML}{7A5AE0}
\definecolor{popgreen}{HTML}{2E9E6A}
\definecolor{poppink}{HTML}{E0457B}
\definecolor{popblue}{HTML}{2F7DE1}
\definecolor{popgold}{HTML}{C98A12}
\definecolor{popmuted}{HTML}{8A7F76}
\definecolor{popline}{HTML}{EADFD2}
\definecolor{popgray}{HTML}{8A7F76}
\colorlet{popgrey}{popgray}
% Alias tolérants (le modèle nomme parfois les couleurs différemment)
\colorlet{popwhite}{white}
\colorlet{popblack}{popink}
\colorlet{popred}{poppink}
\colorlet{popyellow}{popgold}
% Teintes claires (fonds de tableaux/encadrés) — le modèle nomme parfois
% les couleurs avec un suffixe L (ex. popblueL) sans les avoir définies.
\colorlet{poporangeL}{poporange!20}
\colorlet{popdarkL}{popdark!20}
\colorlet{poppurpleL}{poppurple!20}
\colorlet{popgreenL}{popgreen!20}
\colorlet{poppinkL}{poppink!20}
\colorlet{popblueL}{popblue!20}
\colorlet{popgoldL}{popgold!20}
\colorlet{popredL}{popred!20}
\colorlet{popgrayL}{popgray!20}
% Teintes claires pour fonds de tableaux / zones
\colorlet{popblueL}{popblue!10!white}
\colorlet{poporangeL}{poporange!14!white}
\colorlet{popgreenL}{popgreen!12!white}
\colorlet{poppurpleL}{poppurple!12!white}
\colorlet{poppinkL}{poppink!12!white}
\colorlet{popgoldL}{popgold!14!white}

\pagecolor{popcream}

% ============================================================
% Typographie
% ============================================================
\defaultfontfeatures{Ligatures=TeX}
\setmainfont{PlusJakartaSans-Medium.ttf}[Path=../../../fonts/,BoldFont=PlusJakartaSans-Bold.ttf]
% Maths en sans-serif (Fira Math) avec chiffres et lettres droites de Plus
% Jakarta, pour que les formules se fondent dans le texte.
\usepackage[math-style=ISO,bold-style=ISO]{unicode-math}
\setmathfont{FiraMath-Regular.otf}[Path=../../../fonts/]
\setmathfont{PlusJakartaSans-Medium.ttf}[Path=../../../fonts/,range={up/num,up/{Latin,latin},\mathup}]
\usepackage{icomma} % virgule décimale sans espace en mode math
% Caractères mathématiques Unicode absents des polices de texte (Plus Jakarta,
% Archivo Black) : rendus par leur équivalent mathématique, en texte comme en
% formule — sinon ils s'affichent en carré vide (ex. « Arithmétique dans ℤ »).
\usepackage{newunicodechar}
\newunicodechar{ℝ}{\ensuremath{\mathbb{R}}}
\newunicodechar{ℤ}{\ensuremath{\mathbb{Z}}}
\newunicodechar{ℂ}{\ensuremath{\mathbb{C}}}
\newunicodechar{ℕ}{\ensuremath{\mathbb{N}}}
\newunicodechar{ℚ}{\ensuremath{\mathbb{Q}}}
\newunicodechar{𝔻}{\ensuremath{\mathbb{D}}}
\newunicodechar{α}{\ensuremath{\alpha}}
\newunicodechar{β}{\ensuremath{\beta}}
\newunicodechar{γ}{\ensuremath{\gamma}}
\newunicodechar{δ}{\ensuremath{\delta}}
\newunicodechar{ε}{\ensuremath{\varepsilon}}
\newunicodechar{θ}{\ensuremath{\theta}}
\newunicodechar{λ}{\ensuremath{\lambda}}
\newunicodechar{μ}{\ensuremath{\mu}}
\newunicodechar{σ}{\ensuremath{\sigma}}
\newunicodechar{τ}{\ensuremath{\tau}}
\newunicodechar{φ}{\ensuremath{\varphi}}
\newunicodechar{ψ}{\ensuremath{\psi}}
\newunicodechar{ω}{\ensuremath{\omega}}
\newunicodechar{Σ}{\ensuremath{\Sigma}}
% majuscules produites par \MakeUppercase dans les titres (α-aminés -> Α)
\newunicodechar{Α}{\ensuremath{\alpha}}
\newunicodechar{Β}{\ensuremath{\beta}}
\newunicodechar{Γ}{\ensuremath{\Gamma}}
\newunicodechar{Δ}{\ensuremath{\Delta}}
\newunicodechar{Ε}{\ensuremath{\varepsilon}}
\newunicodechar{Θ}{\ensuremath{\Theta}}
\newunicodechar{Λ}{\ensuremath{\Lambda}}
\newunicodechar{Μ}{\ensuremath{\mu}}
\newunicodechar{Τ}{\ensuremath{\tau}}
\newunicodechar{Φ}{\ensuremath{\Phi}}
\newunicodechar{Ψ}{\ensuremath{\Psi}}
\newunicodechar{Ω}{\ensuremath{\Omega}}
\newunicodechar{↦}{\ensuremath{\mapsto}}
\newunicodechar{∈}{\ensuremath{\in}}
\newunicodechar{∉}{\ensuremath{\notin}}
\newunicodechar{∧}{\ensuremath{\wedge}}
\newunicodechar{⇔}{\ensuremath{\Leftrightarrow}}
\newunicodechar{⟺}{\ensuremath{\Longleftrightarrow}}
\newunicodechar{⇒}{\ensuremath{\Rightarrow}}
\newunicodechar{⟹}{\ensuremath{\Longrightarrow}}
\newunicodechar{∘}{\ensuremath{\circ}}
\newunicodechar{∪}{\ensuremath{\cup}}
\newunicodechar{∩}{\ensuremath{\cap}}
\newunicodechar{≡}{\ensuremath{\equiv}}
\newunicodechar{⊂}{\ensuremath{\subset}}
\newunicodechar{⊥}{\ensuremath{\perp}}
\newunicodechar{∃}{\ensuremath{\exists}}
\newunicodechar{∀}{\ensuremath{\forall}}
\newunicodechar{∛}{\ensuremath{\sqrt[3]{\,}}}
\newunicodechar{∜}{\ensuremath{\sqrt[4]{\,}}}
\newunicodechar{⃗}{\ensuremath{{}^{\to}}}
\newunicodechar{ⁿ}{\ifmmode^{n}\else\textsuperscript{\ensuremath{n}}\fi}
\newunicodechar{ˣ}{\ifmmode^{x}\else\textsuperscript{\ensuremath{x}}\fi}
\newunicodechar{ᵇ}{\ifmmode^{b}\else\textsuperscript{\ensuremath{b}}\fi}
\newunicodechar{ʳ}{\ifmmode^{r}\else\textsuperscript{\ensuremath{r}}\fi}
\newunicodechar{ᵘ}{\ifmmode^{u}\else\textsuperscript{\ensuremath{u}}\fi}
\newunicodechar{ʸ}{\ifmmode^{y}\else\textsuperscript{\ensuremath{y}}\fi}
\newunicodechar{ᵃ}{\ifmmode^{a}\else\textsuperscript{\ensuremath{a}}\fi}
\newunicodechar{ᵉ}{\ifmmode^{e}\else\textsuperscript{\ensuremath{e}}\fi}
\newunicodechar{ᵏ}{\ifmmode^{k}\else\textsuperscript{\ensuremath{k}}\fi}
\newunicodechar{ᵖ}{\ifmmode^{p}\else\textsuperscript{\ensuremath{p}}\fi}
\newunicodechar{ᴺ}{\ifmmode^{N}\else\textsuperscript{\ensuremath{N}}\fi}
\newunicodechar{ᶜ}{\ifmmode^{c}\else\textsuperscript{\ensuremath{c}}\fi}
\newunicodechar{ʰ}{\ifmmode^{h}\else\textsuperscript{\ensuremath{h}}\fi}
\newunicodechar{ᵠ}{\ifmmode^{q}\else\textsuperscript{\ensuremath{q}}\fi}
\newunicodechar{ₙ}{\ifmmode_{n}\else\textsubscript{\ensuremath{n}}\fi}
\newunicodechar{ₐ}{\ifmmode_{a}\else\textsubscript{\ensuremath{a}}\fi}
\newunicodechar{ᵢ}{\ifmmode_{i}\else\textsubscript{\ensuremath{i}}\fi}
\newunicodechar{ᵣ}{\ifmmode_{r}\else\textsubscript{\ensuremath{r}}\fi}
\newunicodechar{ₖ}{\ifmmode_{k}\else\textsubscript{\ensuremath{k}}\fi}
\newunicodechar{ᵦ}{\ifmmode_{\beta}\else\textsubscript{\ensuremath{\beta}}\fi}
\newunicodechar{ₓ}{\ifmmode_{x}\else\textsubscript{\ensuremath{x}}\fi}
\newfontfamily\popdisplay{ArchivoBlack-Regular.ttf}[Path=../../../fonts/]
\newfontfamily\poplogo{SpaceGrotesk-Bold.ttf}[Path=../../../fonts/]
\newfontfamily\popsemi{PlusJakartaSans-SemiBold.ttf}[Path=../../../fonts/]
\newfontfamily\popxbold{PlusJakartaSans-ExtraBold.ttf}[Path=../../../fonts/]
\newfontfamily\popxboldls{PlusJakartaSans-ExtraBold.ttf}[Path=../../../fonts/,LetterSpace=3.0]

\setlist[enumerate]{leftmargin=1.7em,itemsep=5pt,topsep=4pt}
\setlist[itemize]{leftmargin=1.7em,itemsep=4pt,topsep=4pt,label={\color{poporange}\textbullet}}
\setlength{\parindent}{0pt}
\setlength{\parskip}{5pt}
\renewcommand{\arraystretch}{1.25}

% pgfplots : style Pop par défaut
\pgfplotsset{
  every axis/.append style={axis line style={popink,line width=1pt},tick style={popink},
    grid style={popline,line width=0.5pt},label style={font=\small},tick label style={font=\footnotesize}},
  popcurve/.style={poporange,line width=1.8pt,no markers,smooth},
}
% Même style disponible dans un \draw TikZ simple (hors pgfplots)
\tikzset{popcurve/.style={poporange,line width=1.8pt}}
\tikzset{popgrid/.style={popline,very thin}}
\colorlet{popcurve}{poporange} % le nom du style est parfois utilisé comme couleur

% ============================================================
% Pill / squiggle
% ============================================================
\newcommand{\poppill}[3]{%
  \tikz[baseline=-0.55ex]\node[draw=popink,line width=1.2pt,rounded corners=2.6pt,%
    fill=#2,inner xsep=6.5pt,inner ysep=3pt]{%
    \popxboldls\fontsize{7}{7}\selectfont\color{#3}\MakeUppercase{#1}};%
}
\newcommand{\popsquiggle}[1]{%
  \tikz[baseline=-0.4ex]{
    \draw[#1,line width=2.6pt,line cap=round]
      plot [smooth] coordinates {(0,0.09) (0.34,0.01) (0.68,0.21) (1.02,0.01) (1.36,0.10)};
  }%
}
% Mot-clé surligné (marqueur orange sous le texte)
\newcommand{\cle}[1]{{\popxbold\color{popdark}#1}}

% ============================================================
% En-tête / pied
% ============================================================
\pagestyle{fancy}
\fancyhf{}
\renewcommand{\headrulewidth}{0pt}
\renewcommand{\footrulewidth}{0pt}
\lhead{\raisebox{-0.15ex}{\poplogo\fontsize{13}{13}\selectfont\color{popink}OnBuch\color{poporange}+}}
\rhead{\poppill{\DOCMATIERE\ \textbullet\ \DOCNIVEAU\ \textbullet\ Leçon \DOCLECON}{white}{popink}}
\lfoot{\popxbold\fontsize{7.6}{7.6}\selectfont\color{popmuted}\MakeUppercase{Cours signé OnBuch+}\ \textbullet\ {\popsemi\DOCTITRE}}
\rfoot{\tikz[baseline=-0.6ex]\node[rounded corners=8.5pt,fill=popink,minimum height=17pt,minimum width=17pt,inner xsep=5pt,inner sep=0]{\popxbold\fontsize{8}{8}\selectfont\color{popcream}\thepage\,/\,\pageref*{LastPage}};}

% ============================================================
% Titres
% ============================================================
\newcommand{\chaptitre}[2]{%
  \begin{tcolorbox}[enhanced,colback=poporange,colframe=popink,boxrule=2.6pt,arc=11pt,
      drop shadow=popink,left=15pt,right=15pt,top=12pt,bottom=11pt,before skip=3pt,after skip=18pt]
    \poppill{#2}{popink}{popcream}\par\vspace{9pt}
    {\raggedright\hyphenpenalty=10000\exhyphenpenalty=10000\popdisplay\fontsize{22}{24}\selectfont\color{popcream}\MakeUppercase{#1}\par}
  \end{tcolorbox}
}
% Section numérotée : pastille orange avec le numéro + titre Archivo
\newcounter{popsec}
\newcommand{\coursec}[1]{%
  \stepcounter{popsec}\setcounter{popsub}{0}%
  \par\vspace{16pt}\noindent
  \begin{minipage}{\linewidth}
  \tikz[baseline=-0.7ex]\node[draw=popink,line width=1.6pt,fill=poporange,rounded corners=4pt,minimum size=22pt,inner sep=0]{\popdisplay\fontsize{12}{12}\selectfont\color{popcream}\thepopsec};%
  \hspace{8pt}{\popdisplay\fontsize{15}{17}\selectfont\color{popink}\MakeUppercase{#1}}\par
  \vspace{4pt}\noindent{\color{poporange}\rule{36pt}{3.6pt}}
  \end{minipage}\par\nopagebreak\vspace{8pt}
}
\newcounter{popsub}
\newcommand{\courssub}[1]{%
  \stepcounter{popsub}%
  \par\vspace{9pt}\noindent{\popxbold\fontsize{11.5}{13}\selectfont\color{popink}{\color{poporange}\thepopsec.\thepopsub}\hspace{6pt}#1}\par\nopagebreak\vspace{4pt}
}

% ============================================================
% Boîtes pédagogiques — même recette : fond blanc, bordure épaisse colorée,
% ombre dure encre, badge plein. Titre optionnel [#1].
% ============================================================
\tcbset{popbase/.style={breakable,enhanced,colback=white,boxrule=2.3pt,arc=9pt,
  shadow={3pt}{-3pt}{0pt}{popink},left=14pt,right=14pt,top=11pt,bottom=11pt,before skip=11pt,after skip=15pt,
  fontupper=\popsemi\fontsize{10.6}{14.5}\selectfont\color{popink},
  fontlower=\popsemi\fontsize{10.6}{14.5}\selectfont\color{popink}}}
% Coche dessinée (le glyphe ✓ n'existe pas dans les polices du document)
\newcommand{\popcheck}[1]{\tikz[baseline=-0.5ex]\draw[#1,line width=1.6pt,line cap=round,line join=round](-0.13,0.0)--(-0.04,-0.09)--(0.13,0.1);}
\providecommand{\coche}{\popcheck{popgreen}}
\providecommand{\resultat}[1]{\par\smallskip\noindent\fcolorbox{popgreen}{popgreenL}{\parbox{\dimexpr\linewidth-2\fboxsep-2\fboxrule\relax}{\textbf{Résultat.} #1}}\par\smallskip}
\newcommand{\popboxhead}[3]{\poppill{#1}{#2}{white}\ifx\relax#3\relax\else\hspace{6pt}{\popxbold\fontsize{10.6}{12}\selectfont\color{popink}#3}\fi\par\vspace{7pt}}

\newtcolorbox{aretenir}[1][]{popbase,colframe=popgreen,before upper={\popboxhead{À retenir}{popgreen}{#1}}}
% Texte en allemand (document de lecture, dialogue, extrait) : cadre
% d'étude. Un titre optionnel (source, auteur) s'affiche dans l'en-tête.
\newtcolorbox{textebox}[1][]{popbase,colframe=popblue,before upper={\popboxhead{Text}{popblue}{#1}\begin{otherlanguage}{ngerman}},after upper={\end{otherlanguage}}}
% Marques ✓ / ✗ (forme correcte / fautive) souvent utilisées par les modèles
\providecommand{\cross}{\textcolor{poppink}{\ensuremath{\times}}}
\providecommand{\xmark}{\cross}\providecommand{\crossmark}{\cross}\providecommand{\cmark}{\ensuremath{\checkmark}}
% Ligne de réponse / justification à compléter (inventée par les modèles)
\providecommand{\justification}[1]{\par\noindent\textit{Justification~:} \dotfill\par}
% \textipa (package tipa non chargé) : la police ne couvre pas l'API -> simple passage
\providecommand{\textipa}[1]{#1}
% Soulignement (ulem non chargé) : \uline, \ul
\providecommand{\uline}[1]{\underline{#1}}\providecommand{\ul}[1]{\underline{#1}}
% \alert (beamer) inventée par les modèles : texte mis en évidence
\providecommand{\alert}[1]{\textbf{\textcolor{poppink}{#1}}}
% variantes ASCII d'ordinaux écrites en macro
\providecommand{\iere}{\textsuperscript{re}}\providecommand{\ieme}{\textsuperscript{e}}\providecommand{\ere}{\textsuperscript{re}}\providecommand{\eme}{\textsuperscript{e}}\providecommand{\er}{\textsuperscript{er}}\providecommand{\ie}{\textsuperscript{e}}
% Ordinaux français écrits en macro par les modèles : 1\ière, 3\ième
\newcommand{\ière}{\textsuperscript{re}}\newcommand{\ième}{\textsuperscript{e}}
% \exemple (inventée par les modèles) : étiquette « Exemple : »
\providecommand{\exemple}{\textbf{Exemple~:}\ }
% \square utilisable aussi hors mode math (cases à cocher)
\AtBeginDocument{\let\squareMath\square\renewcommand{\square}{\ensuremath{\squareMath}}}
% Mot ou phrase en allemand à l'intérieur d'un paragraphe français
\newcommand{\all}[1]{\ifmmode\text{\textit{#1}}\else\begingroup\selectlanguage{ngerman}\itshape #1\endgroup\fi}
\let\esp\all % alias toléré
\newtcolorbox{exemplebox}[1][]{popbase,colframe=poporange,before upper={\popboxhead{Exemple}{poporange}{#1}}}
\newtcolorbox{definition}[1][]{popbase,colframe=popblue,before upper={\popboxhead{Définition}{popblue}{#1}}}
\newtcolorbox{propriete}[1][]{popbase,colframe=poppurple,before upper={\popboxhead{Propriété}{poppurple}{#1}}}
\newtcolorbox{methode}[1][]{popbase,colframe=popgold,before upper={\popboxhead{Méthode}{popgold}{#1}}}
\newtcolorbox{attention}[1][]{popbase,colframe=poppink,before upper={\popboxhead{Attention}{poppink}{#1}}}
\newtcolorbox{experience}[1][]{popbase,colframe=popblue,colback=popblueL,before upper={\popboxhead{Expérience}{popblue}{#1}}}
% « Le savais-tu ? » : carte violette pleine (contexte, culture, Cameroun)
\newtcolorbox{savaistu}[1][]{popbase,colback=poppurple,colframe=popink,
  fontupper=\popsemi\fontsize{10.4}{14.2}\selectfont\color{white},
  before upper={\poppill{Le savais-tu\,?}{white}{poppurple}\ifx\relax#1\relax\else\hspace{6pt}{\popxbold\color{white}#1}\fi\par\vspace{7pt}}}
% Exercice résolu : énoncé en haut, \tcblower puis la solution sur fond crème
\newcounter{popexo}\newcounter{popexr}
\newtcolorbox{exoresolu}[1][]{popbase,colframe=poporange,colbacklower=popcream,
  segmentation style={popink,line width=1.2pt,dashed},
  before upper={\stepcounter{popexr}\popboxhead{Exercice résolu \thepopexr}{poporange}{#1}},
  before lower={\poppill{Solution}{popink}{popcream}\par\vspace{6pt}}}
% Exercice (non résolu en place) — la correction est dans la partie Corrigés
\newtcolorbox{exercice}[1][]{popbase,colframe=popink,
  before upper={\stepcounter{popexo}\popboxhead{Exercice \thepopexo}{popink}{#1}}}
% Corrigé
\newtcolorbox{corrige}[1][]{popbase,colframe=popgreen,colback=popgreenL,
  before upper={\popboxhead{Corrigé}{popgreen}{#1}}}
% Objectifs / prérequis / bilan
\newtcolorbox{objectifs}{popbase,colframe=popink,colback=poporangeL,
  before upper={\popboxhead{Objectifs de la leçon}{popink}{}}}
\newtcolorbox{prerequis}{popbase,colframe=popink,colback=popblueL,
  before upper={\popboxhead{Prérequis}{popblue}{}}}
\newtcolorbox{bilan}{popbase,colframe=popink,colback=white,boxrule=2.8pt,
  before upper={\popboxhead{Fiche bilan}{poporange}{L'essentiel à connaître pour l'examen}}}
% Figure encadrée avec légende
\newtcolorbox{popfigure}[1][]{enhanced,breakable=false,colback=white,colframe=popline,boxrule=1.4pt,arc=8pt,
  left=10pt,right=10pt,top=10pt,bottom=8pt,before skip=10pt,after skip=12pt,halign=center,
  lower separated=false,halign lower=center,
  fontlower=\popsemi\fontsize{9}{11.5}\selectfont\color{popmuted},
  #1}
\newcommand{\legende}[1]{\par\vspace{6pt}{\popsemi\fontsize{9}{11.5}\selectfont\color{popmuted}#1}}

% ============================================================
% Signature OnBuch+
% ============================================================
\newcommand{\poplogoinline}[1]{{\poplogo\fontsize{#1}{#1}\selectfont\color{popink}OnBuch\color{poporange}+}}
\newcommand{\signatureonbuch}{%
  \begin{tcolorbox}[enhanced,colback=popink,colframe=popink,boxrule=2.2pt,arc=10pt,
      drop shadow=poporange,left=14pt,right=14pt,top=10pt,bottom=10pt,before skip=0pt,after skip=0pt]
    \tikz[baseline=-0.7ex]\node[circle,fill=popgreen,draw=popcream,line width=1.3pt,minimum size=22pt,inner sep=0]{\popcheck{white}};%
    \hspace{9pt}\begin{minipage}[c]{0.62\linewidth}
      {\popxbold\fontsize{12}{13}\selectfont\color{popcream}Signé\ \textbullet\ {\poplogo OnBuch\color{poporange}+}}\par\vspace{2pt}
      {\raggedright\popsemi\fontsize{8}{10.5}\selectfont\color{popline}\MakeUppercase{Équipe pédagogique OnBuch+}\\\MakeUppercase{Conforme au programme officiel MINESEC}\par}
    \end{minipage}\hfill
    {\poplogo\fontsize{22}{22}\selectfont\color{popcream}OnBuch\color{poporange}+}
  \end{tcolorbox}%
}

% ============================================================
% Page de garde
% #1 titre · #2 matière · #3 niveau · #4 module · #5 n° leçon · #6 durée ·
% #7 description / objectifs (texte ou itemize)
% ============================================================
\newcommand{\pagegarde}[7]{%
  \thispagestyle{empty}
  \newgeometry{a4paper,margin=1.7cm,top=1.6cm,bottom=1.4cm}%
  \begin{tikzpicture}[remember picture,overlay]
    \fill[poporange!18] ([xshift=-3.2cm,yshift=-3.6cm]current page.north east) circle (4.6cm);
    \fill[poppurple!14] ([xshift=2.4cm,yshift=7.6cm]current page.south west) circle (3.8cm);
  \end{tikzpicture}%
  {\poplogo\fontsize{30}{30}\selectfont\color{popink}OnBuch\color{poporange}+}\hfill
  \raisebox{0.6ex}{\poppill{Cours signé \textbullet\ Premium}{popink}{popcream}}\par
  \vspace{1pt}\popsquiggle{poppurple}\par
  \vspace{8pt}
  \begin{tcolorbox}[enhanced,colback=poporange,colframe=popink,boxrule=2.8pt,arc=13pt,
      drop shadow=popink,left=18pt,right=18pt,top=12pt,bottom=13pt,after skip=10pt]
    \poppill{#2 \textbullet\ #3}{popink}{popcream}\hspace{5pt}\poppill{#4}{white}{popink}\par\vspace{8pt}
    {\popdisplay\fontsize{14}{14}\selectfont\color{popink}LEÇON #5}\par\vspace{4pt}
    {\raggedright\hyphenpenalty=10000\exhyphenpenalty=10000\popdisplay\fontsize{25}{27}\selectfont\color{popcream}\MakeUppercase{#1}\par}
    \vspace{8pt}
    {\popxbold\fontsize{9.5}{9.5}\selectfont\color{popink}\MakeUppercase{Durée conseillée : #6}}
  \end{tcolorbox}
  \begin{tcolorbox}[enhanced,colback=white,colframe=popink,boxrule=2.2pt,arc=10pt,
      drop shadow=popink,left=16pt,right=16pt,top=10pt,bottom=10pt,after skip=8pt]
    \poppill{Dans cette leçon}{poppurple}{white}\par\vspace{5pt}
    \popsemi\fontsize{9.3}{12.6}\selectfont\color{popink}#7
  \end{tcolorbox}
  \noindent\begin{minipage}[t]{0.487\linewidth}
    \begin{tcolorbox}[enhanced,colback=popgreen,colframe=popink,boxrule=2.2pt,arc=10pt,
        drop shadow=popink,left=11pt,right=11pt,top=10pt,bottom=10pt,height=3.1cm,valign=center]
      \tcbox[colback=white,colframe=popink,boxrule=1.3pt,arc=5pt,boxsep=0pt,nobeforeafter,
        left=3pt,right=3pt,top=3pt,bottom=3pt,valign=center]{\qrcode[height=1.9cm]{https://play.google.com/store/apps/details?id=com.luvvix.onbuch&hl=fr}}%
      \hspace{8pt}\begin{minipage}[c]{0.5\linewidth}
        {\popdisplay\fontsize{11}{12.5}\selectfont\color{white}\MakeUppercase{Télécharge OnBuch}}\par\vspace{4pt}
        {\raggedright\popsemi\fontsize{8.4}{11}\selectfont\color{white}Gratuit \textbullet\ Google Play\\Tuteur IA, annales, concours, résultats\par}
      \end{minipage}
    \end{tcolorbox}
  \end{minipage}\hfill
  \begin{minipage}[t]{0.487\linewidth}
    \begin{tcolorbox}[enhanced,colback=poppurple,colframe=popink,boxrule=2.2pt,arc=10pt,
        drop shadow=popink,left=11pt,right=11pt,top=10pt,bottom=10pt,height=3.1cm,valign=center]
      \tikz[baseline=-0.7ex]\node[circle,fill=white,draw=popink,line width=1.3pt,minimum size=1.6cm,inner sep=0]{\popdisplay\fontsize{24}{24}\selectfont\color{poppurple}+};%
      \hspace{8pt}\begin{minipage}[c]{0.6\linewidth}
        {\popdisplay\fontsize{11}{12.5}\selectfont\color{white}\MakeUppercase{Passe à OnBuch+}}\par\vspace{4pt}
        {\raggedright\popsemi\fontsize{8.4}{11}\selectfont\color{white}Tous les cours signés, Tuteur IA illimité, fiches PDF — onglet OnBuch+ de l'app\par}
      \end{minipage}
    \end{tcolorbox}
  \end{minipage}
  \vspace{8pt}
  \popcontenu
  \vfill
  \signatureonbuch
  \restoregeometry
  \newpage
}

% Rangée « Ce cours contient » (page de garde)
\newcommand{\poptile}[3]{% #1 couleur #2 gros texte #3 légende
  \begin{tcolorbox}[enhanced,colback=#1,colframe=popink,boxrule=2pt,arc=9pt,shadow={2.5pt}{-2.5pt}{0pt}{popink},
      left=6pt,right=6pt,top=9pt,bottom=9pt,halign=center,valign=center,height=1.75cm,width=\linewidth,nobeforeafter]
    {\popdisplay\fontsize{12.5}{13}\selectfont\color{white}\MakeUppercase{#2}}\par\vspace{3pt}
    {\centering\popsemi\fontsize{7.8}{9.5}\selectfont\color{white}#3\par}
  \end{tcolorbox}}
\newcommand{\popcontenu}{%
  \par\noindent{\popxboldls\fontsize{7.5}{7.5}\selectfont\color{popmuted}CE COURS SIGNÉ CONTIENT}\par\vspace{7pt}
  \noindent\begin{minipage}[t]{0.235\linewidth}\poptile{popblue}{Cours}{complet, illustré, pas à pas}\end{minipage}\hfill
  \begin{minipage}[t]{0.235\linewidth}\poptile{poporange}{Méthodes}{et exercices résolus}\end{minipage}\hfill
  \begin{minipage}[t]{0.235\linewidth}\poptile{poppink}{Exercices}{type examen + corrigés}\end{minipage}\hfill
  \begin{minipage}[t]{0.235\linewidth}\poptile{popgreen}{Fiche bilan}{l'essentiel à retenir}\end{minipage}\par}

% Sommaire
\newtcolorbox{sommaire}{enhanced,colback=white,colframe=popink,boxrule=2.2pt,arc=10pt,
  drop shadow=popink,left=18pt,right=18pt,top=13pt,bottom=13pt,before skip=6pt,after skip=20pt,
  before upper={\poppill{Sommaire}{poppurple}{white}\par\vspace{9pt}\popsemi\fontsize{10.6}{14.5}\selectfont\color{popink}}}

% Signature de fin de cours — poussée tout en bas de la dernière page
\newcommand{\findecours}{%
  \par\vfill\nopagebreak
  \begin{center}\popsquiggle{poporange}\end{center}\vspace{4pt}
  \signatureonbuch
}
\AtBeginDocument{\def\llbracket{\mathopen{[\mkern-3.5mu[}}\def\rrbracket{\mathclose{]\mkern-3.5mu]}}} % intervalles d'entiers (glyphe ⟦ absent de la police)
% Glyphes absents de Fira Math : redéfinis à partir de glyphes disponibles
\AtBeginDocument{%
  \def\setminus{\mathbin{\backslash}}\let\smallsetminus\setminus
  \def\vdots{\mathord{\vbox{\baselineskip3.2pt\lineskiplimit0pt\kern4pt\hbox{.}\hbox{.}\hbox{.}}}}%
  \def\ddots{\mathinner{\mkern1mu\raise6.5pt\hbox{.}\mkern2mu\raise3.6pt\hbox{.}\mkern2mu\raise0.7pt\hbox{.}\mkern1mu}}%
  \def\triangle{\mathord{\scalebox{0.9}{$\symup{\Delta}$}}}%
  \def\nabla{\mathord{\raisebox{\depth}{\scalebox{1}[-1]{$\symup{\Delta}$}}}}%
  \def\checkmark{\popcheck{popdark}}%
  \def\star{\mathbin{\ast}}\def\bigstar{\mathord{\ast}}%
  \def\diamond{\mathbin{\scalebox{0.6}{$\Diamond$}}}%
  \def\bigcup{\mathop{\mathchoice{\raisebox{-0.3em}{\scalebox{1.7}{$\cup$}}}{\scalebox{1.25}{$\cup$}}{\cup}{\cup}}\displaylimits}%
  \def\bigcap{\mathop{\mathchoice{\raisebox{-0.3em}{\scalebox{1.7}{$\cap$}}}{\scalebox{1.25}{$\cap$}}{\cap}{\cap}}\displaylimits}%
  \def\lesseqqgtr{\mathrel{\lessgtr}}\def\bigoplus{\mathop{\oplus}\displaylimits}%
}
\providecommand{\ssi}{si et seulement si} % abréviation fréquente
\AtBeginDocument{\providecommand{\wideparen}{\overparen}} % arc de cercle

%% FIGFIX-BEGIN v5 — correctifs globaux des figures (docs/onbuch-generation/tools/figfix)
\usepackage{adjustbox}\usepackage{etoolbox}\tcbuselibrary{raster}
% toute figure TikZ/pgfplots plus large que la ligne est réduite à la largeur disponible
\BeforeBeginEnvironment{tikzpicture}{\begin{adjustbox}{max width=\linewidth}}
\AfterEndEnvironment{tikzpicture}{\end{adjustbox}}
% légendes pgfplots lisibles par-dessus les courbes
\pgfplotsset{every axis/.append style={legend style={fill=white,fill opacity=0.92,text opacity=1,draw=black!15,inner sep=3pt}}}
% étiquettes d'arêtes (syntaxe edge["..."]) avec fond blanc
\tikzset{every edge quotes/.append style={fill=white,inner sep=1.2pt}}
%% FIGFIX-END
\begin{document}

\pagegarde{Les marchés publics (corruption, contrebande, falsification des produits)}{Allemand}{Tle A}{Chapitre 2 : Vie économique}{AL08}{2 h}{Cette leçon de type texte propose la lecture et le commentaire d'un article de presse allemand sur les marchés publics, la corruption, la contrebande et la falsification des produits. Elle consolide le lexique économique et judiciaire et introduit deux outils grammaticaux essentiels du journalisme allemand : le passif d'état (Zustandspassiv) et le discours indirect au subjonctif I.
\vspace{4pt}
\begin{multicols}{2}\small
\begin{itemize}
\item À la fin de cette leçon, je sais comprendre globalement puis en détail un texte de presse allemand sur les marchés publics, la corruption, la contrebande et la falsification des produits.
\item À la fin de cette leçon, je sais employer correctement le vocabulaire thématique (der Markt, die Korruption, der Schmuggel, die Fälschung, die Ausschreibung, der Betrug) avec articles, pluriels et familles de mots.
\item À la fin de cette leçon, je sais reconnaître et former le passif d'état (Zustandspassiv) et le distinguer du passif d'action (Vorgangspassiv).
\item À la fin de cette leçon, je sais former le subjonctif I à tous les temps utiles et l'employer pour rapporter des paroles au discours indirect.
\item À la fin de cette leçon, je sais conjuguer les verbes forts et faibles au subjonctif I et au subjonctif II et choisir le bon auxiliaire.
\item À la fin de cette leçon, je sais répondre par écrit et à l'oral à des questions de compréhension sur un texte économique.
\end{itemize}\end{multicols}}

\begin{sommaire}
\begin{multicols}{2}
\begin{enumerate}
\item Texte support : lecture et première compréhension
\item Lexique thématique : marchés publics et criminalité économique
\item Grammaire 1 : le passif d'état (Zustandspassiv)
\item Grammaire 2 : le discours indirect au subjonctif I
\item Conjugaison : subjonctif I et II, tableaux de synthèse
\item Méthode du commentaire de texte et commentaire modèle
\item Production guidée : écrire et parler sur la criminalité économique
\item Activité d'intégration
\item Méthodes et exercices résolus
\item Exercices
\item Corrigés des exercices
\item Fiche bilan
\end{enumerate}
\end{multicols}
\end{sommaire}

\begin{objectifs}
\begin{itemize}
\item À la fin de cette leçon, je sais comprendre globalement puis en détail un texte de presse allemand sur les marchés publics, la corruption, la contrebande et la falsification des produits.
\item À la fin de cette leçon, je sais employer correctement le vocabulaire thématique (der Markt, die Korruption, der Schmuggel, die Fälschung, die Ausschreibung, der Betrug) avec articles, pluriels et familles de mots.
\item À la fin de cette leçon, je sais reconnaître et former le passif d'état (Zustandspassiv) et le distinguer du passif d'action (Vorgangspassiv).
\item À la fin de cette leçon, je sais former le subjonctif I à tous les temps utiles et l'employer pour rapporter des paroles au discours indirect.
\item À la fin de cette leçon, je sais conjuguer les verbes forts et faibles au subjonctif I et au subjonctif II et choisir le bon auxiliaire.
\item À la fin de cette leçon, je sais répondre par écrit et à l'oral à des questions de compréhension sur un texte économique.
\item À la fin de cette leçon, je sais rédiger un résumé structuré et un court commentaire de texte en allemand.
\item À la fin de cette leçon, je sais produire un court paragraphe d'opinion guidé sur la lutte contre la corruption au Cameroun ou en Allemagne.
\end{itemize}
\end{objectifs}

\begin{prerequis}
\begin{itemize}
\item Le passif d'action (Vorgangspassiv) au présent et au passé composé (classe de Première)
\item Le subjonctif II et ses emplois (souhait, hypothèse, conseil)
\item Les temps principaux de l'indicatif : présent, Präteritum, Perfekt, Plusquamperfekt
\item Le vocabulaire de base de l'économie (Geld, Wirtschaft, Unternehmen, Preis)
\item La construction de la phrase allemande : place du verbe, subordonnées avec dass, weil, obwohl
\end{itemize}
\end{prerequis}

\newpage

\coursec{Texte support : lecture et première compréhension}

À Yaoundé, un chantier de route est attribué à une entreprise après un appel d'offres : mais que se passe-t-il quand des pots-de-vin faussent le jeu ? Quand les médicaments vendus au marché sont des contrefaçons dangereuses ? En Allemagne aussi, les journaux dénoncent régulièrement la \all{Korruption} (corruption), le \all{Schmuggel} (contrebande) et la \all{Fälschung} de produits, des médicaments aux pièces automobiles. Ces fléaux ont un nom collectif : la \all{Wirtschaftskriminalität}, la criminalité économique.

\textbf{Problématique :} comment la presse allemande rapporte-t-elle ces faits sans accuser directement des personnes qui ne sont pas encore condamnées ? Grâce à deux outils linguistiques précis : le \cle{discours indirect au subjonctif I} (\all{der Minister erklärte, das Verfahren sei korrekt gewesen}) et le \cle{passif} (\all{gefälschte Medikamente wurden beschlagnahmt}). Cette leçon vous donne les clés pour lire, comprendre et commenter un article économique allemand sur ces sujets d'actualité — au Cameroun comme dans les pays germanophones.

\courssub{Situation d'entrée et objectifs}

\begin{aretenir}[Objectifs de la section]
À la fin de cette section, vous serez capable de :
\begin{itemize}
\item identifier le thème, le genre et le ton d'un article de presse allemand ;
\item employer le vocabulaire des marchés publics et de la criminalité économique ;
\item repérer les verbes de parole et les faits rapportés (préparation au discours indirect) ;
\item répondre à des questions de compréhension et résumer chaque paragraphe en allemand.
\end{itemize}
\end{aretenir}

\courssub{Texte support : Korruption und Fälschung – Wenn der Markt krankt}

\begin{methode}[Lire un article de presse allemand en quatre étapes]
\begin{enumerate}
\item \textbf{Lire le titre et le chapeau} : ils annoncent le thème et les faits essentiels (qui ? quoi ? où ?).
\item \textbf{Repérer les verbes de parole} : \all{sagen} (dire), \all{behaupten} (prétendre, affirmer), \all{erklären} (déclarer), \all{berichten} (rapporter). Ils signalent le discours rapporté.
\item \textbf{Distinguer faits et opinions} : un fait est vérifiable (\all{wurden beschlagnahmt}) ; une opinion est une affirmation de quelqu'un (\all{behauptet die Opposition}).
\item \textbf{Souligner le vocabulaire du thème} : ici, les mots de la criminalité économique. Notez chaque nom avec son article et son pluriel.
\end{enumerate}
\end{methode}

\begin{textebox}[Korruption und Fälschung – Wenn der Markt krankt]
München/Douala. Bei einer öffentlichen Ausschreibung für den Bau einer Schule in Bayern sei ein Unternehmen bevorzugt worden, behauptet die Opposition im Landtag. Ein Mitarbeiter des Ministeriums habe Geld angenommen, damit die Firma den Zuschlag erhalte. Der zuständige Minister erklärte jedoch gestern vor Journalisten, das Verfahren sei korrekt gewesen; alle Angebote seien fair geprüft worden. Die Staatsanwaltschaft ermittle nun, sagte ein Sprecher, aber es gebe noch keine Beweise.

Unterdessen wurden am Flughafen München mehrere Kisten mit gefälschten Medikamenten beschlagnahmt. Die Waren waren in einem Container versteckt und sollten über Hamburg nach Westafrika transportiert werden. Ein Zollbeamter berichtete, die Packungen sähen täuschend echt aus, aber der Inhalt sei wirkungslos oder sogar gefährlich. Besonders Antibiotika und Malariamittel würden häufig gefälscht.

Ein Experte für Wirtschaftskriminalität sagte in einem Interview, der Schmuggel habe in den letzten Jahren deutlich zugenommen. Die Täter seien gut organisiert und nutzten neue Routen. Er forderte härtere Strafen und eine bessere Zusammenarbeit zwischen den Ländern. „Wer gefälschte Medikamente verkauft, riskiert Menschenleben“, warnte er.

Auch in Kamerun kämpfen die Behörden gegen Fälschung und Schmuggel, besonders an den Grenzen und auf den Märkten von Douala und Yaoundé. Verbraucher sollen verdächtige Produkte melden, erklärte ein Beamter. Denn nur gemeinsam könne man den Markt schützen.
\end{textebox}

\textbf{Glossaire allemand-français}

\begin{center}
\begin{tabularx}{\linewidth}{|l|X|}
\hline
\rowcolor{popblueL} \textbf{Deutsch} & \textbf{Français} \\
\hline
die Ausschreibung, -en & l'appel d'offres \\
\hline
der Zuschlag, \dots Zuschläge & l'adjudication (attribution du marché) \\
\hline
bevorzugen & favoriser, avantager \\
\hline
die Bestechung, -en & la corruption (active), les pots-de-vin \\
\hline
das Angebot, -e & l'offre, la soumission \\
\hline
die Staatsanwaltschaft, -en & le parquet, le ministère public \\
\hline
ermitteln & enquêter \\
\hline
der Beweis, -e & la preuve \\
\hline
beschlagnahmen & saisir, confisquer \\
\hline
schmuggeln / der Schmuggel & faire passer en contrebande / la contrebande \\
\hline
fälschen / die Fälschung, -en & falsifier, contrefaire / la falsification \\
\hline
der Zollbeamte, -n & le douanier \\
\hline
täuschend echt & à s'y méprendre, d'un réalisme trompeur \\
\hline
die Strafe, -n & la peine, la sanction \\
\hline
der Betrug, \dots Betrüge & l'escroquerie, la fraude \\
\hline
melden & signaler, dénoncer \\
\hline
\end{tabularx}
\end{center}

\begin{definition}[die Korruption]
\all{Die Korruption} (féminin, sans pluriel courant) désigne \all{die Bestechung von Amtsträgern, um Vorteile zu erlangen} : la corruption d'agents publics pour obtenir des avantages (un marché, un document, un service). Famille de mots : \all{bestechen} (corrompre, verbe fort : \emph{besticht, bestach, hat bestochen}), \all{bestechlich} (corruptible), \all{der Bestecher} (le corrupteur). Attention au faux ami : \all{die Korruption} se dit bien « la corruption », mais on dit « un pot-de-vin » pour \all{das Schmiergeld}.
\end{definition}

\courssub{Première lecture : compréhension orale et globale}

\textbf{Déroulement en classe.} 1) Lecture silencieuse individuelle (3 minutes) : chaque élève souligne les mots du thème. 2) Lecture à voix haute par l'enseignant, paragraphe par paragraphe : les élèves écoutent la prononciation (\all{Ausschreibung} se prononce « aous-chraï-boung », \all{beschlagnahmt} « bé-chlak-nâmt »). 3) Compréhension globale collective.

\begin{exemplebox}[Compréhension globale — réponses attendues]
\begin{itemize}
\item \textbf{Thème :} la criminalité économique (marchés publics truqués, contrefaçon, contrebande).
\item \textbf{Genre du texte :} un article de presse (reportage d'actualité avec déclarations).
\item \textbf{Source :} un journal allemand, daté de Munich et Douala (correspondance).
\item \textbf{Ton :} la dénonciation — le journaliste rapporte des accusations graves mais reste prudent : il ne condamne personne lui-même, il \emph{rapporte} les paroles au subjonctif I.
\end{itemize}
\end{exemplebox}

\begin{attention}[Pourquoi le journaliste n'accuse-t-il pas directement ?]
En droit de la presse, tant qu'il n'y a pas de condamnation, on ne peut pas écrire \all{Der Mitarbeiter hat Geld angenommen} (indicatif = fait présenté comme certain). Le journaliste écrit donc \all{Ein Mitarbeiter \emph{habe} Geld angenommen} : le subjonctif I signale « c'est ce qu'on affirme, je ne fais que rapporter ». C'est la prudence juridique du discours indirect — nous l'étudierons en détail dans la section grammaire.
\end{attention}

\courssub{Repérage : qui parle ? quels faits sont rapportés ?}

Ce repérage prépare le travail sur le discours indirect. Complétez mentalement le tableau :

\begin{center}
\begin{tabularx}{\linewidth}{|X|X|X|}
\hline
\rowcolor{popblueL} \textbf{Personne} & \textbf{Verbe de parole} & \textbf{Contenu rapporté} \\
\hline
die Opposition & behauptet & ein Unternehmen sei bevorzugt worden \\
\hline
der Minister & erklärte & das Verfahren sei korrekt gewesen \\
\hline
ein Sprecher & sagte & die Staatsanwaltschaft ermittle \\
\hline
ein Zollbeamter & berichtete & die Packungen sähen echt aus \\
\hline
der Experte & sagte / warnte / forderte & der Schmuggel habe zugenommen \\
\hline
ein Beamter (Kamerun) & erklärte & Verbraucher sollen Produkte melden \\
\hline
\end{tabularx}
\end{center}

\begin{popfigure}
\begin{tikzpicture}[mindmap, grow cyclic, every node/.style={concept}, root concept/.append style={concept color=popblue, font=\bfseries}, level 1/.append style={level distance=3.4cm, sibling angle=90}, level 2/.append style={level distance=2.2cm, sibling angle=45, concept color=popblueL, font=\small}]
\node {Wirtschafts\-kriminalität}
  child[concept color=poporange] { node {Korruption} child { node[align=center] {Geld für den\\Zuschlag} } }
  child[concept color=popgreen] { node {Schmuggel} child { node[align=center] {Waren im\\Container} } }
  child[concept color=poppurple] { node {Fälschung} child { node[align=center] {gefälschte\\Medikamente} } }
  child[concept color=poppink] { node {Betrug} child { node[align=center] {falsche\\Angebote} } };
\end{tikzpicture}
\legende{Carte mentale : les quatre visages de la criminalité économique, avec un exemple tiré du texte pour chacun.}
\end{popfigure}

\courssub{Compréhension détaillée : questions}

\textbf{Questions fermées (vrai/faux) — Richtig oder falsch ?}

\begin{exercice}[Richtig oder falsch ? Korrigieren Sie die falschen Aussagen.]
\begin{enumerate}
\item Die Ausschreibung betrifft den Bau eines Krankenhauses.
\item Der Minister sagt, das Verfahren sei korrekt gewesen.
\item Die gefälschten Medikamente wurden am Bahnhof gefunden.
\item Der Experte fordert härtere Strafen.
\item In Kamerun gibt es keine Probleme mit Fälschungen.
\end{enumerate}
\end{exercice}

\begin{exoresolu}[Questions ouvertes — Antworten Sie auf Deutsch]
\textbf{Énoncé.} Répondez par des phrases complètes :
\begin{enumerate}
\item Wer behauptet, dass ein Unternehmen bevorzugt wurde?
\item Was wurde am Flughafen München beschlagnahmt?
\item Wohin sollten die Waren transportiert werden?
\item Warum sind gefälschte Medikamente gefährlich?
\item Was sollen die Verbraucher in Kamerun tun?
\end{enumerate}
\tcblower
\textbf{Solution détaillée.}
\begin{enumerate}
\item \all{Die Opposition im Landtag behauptet, dass ein Unternehmen bevorzugt wurde.} (C'est l'opposition qui affirme cela — et non le journaliste !)
\item \all{Am Flughafen München wurden mehrere Kisten mit gefälschten Medikamenten beschlagnahmt.} (Notez le passif au Präteritum : \all{wurden beschlagnahmt}.)
\item \all{Die Waren sollten über Hamburg nach Westafrika transportiert werden.}
\item \all{Sie sind gefährlich, weil der Inhalt wirkungslos oder sogar giftig sein kann; der Experte warnt: „Wer gefälschte Medikamente verkauft, riskiert Menschenleben.“}
\item \all{Die Verbraucher sollen verdächtige Produkte melden.}
\end{enumerate}
\end{exoresolu}

\begin{corrige}[Exercice « Richtig oder falsch ? »]
\begin{enumerate}
\item \textbf{Falsch.} \all{Die Ausschreibung betrifft den Bau einer \emph{Schule}.} (une école, pas un hôpital)
\item \textbf{Richtig.} \all{Der Minister erklärte, das Verfahren sei korrekt gewesen.}
\item \textbf{Falsch.} \all{Die Medikamente wurden am \emph{Flughafen} München beschlagnahmt.} (à l'aéroport, pas à la gare)
\item \textbf{Richtig.} \all{Er forderte härtere Strafen und eine bessere Zusammenarbeit.}
\item \textbf{Falsch.} \all{Auch in Kamerun kämpfen die Behörden gegen Fälschung und Schmuggel.}
\end{enumerate}
\end{corrige}

\courssub{Reformulation : résumer chaque paragraphe en une phrase}

\begin{methode}[Résumer un paragraphe en allemand]
\begin{enumerate}
\item Soulignez le sujet principal et le verbe central du paragraphe.
\item Supprimez les détails (chiffres, citations, exemples secondaires).
\item Rédigez une phrase simple : sujet + verbe + complément essentiel. Utilisez le discours indirect si vous rapportez une parole : \all{Der Minister erklärte, \dots}
\end{enumerate}
\end{methode}

\begin{exemplebox}[Résumés modèles — une phrase par paragraphe]
\begin{itemize}
\item \textbf{§ 1 :} \all{Die Opposition behauptet, bei einer Ausschreibung in Bayern sei ein Unternehmen bevorzugt worden; der Minister bestreitet das.} (L'opposition affirme qu'une entreprise a été favorisée ; le ministre le nie.)
\item \textbf{§ 2 :} \all{Am Flughafen München wurden gefälschte Medikamente beschlagnahmt, die nach Westafrika transportiert werden sollten.}
\item \textbf{§ 3 :} \all{Ein Experte sagt, der Schmuggel habe zugenommen, und fordert härtere Strafen.}
\item \textbf{§ 4 :} \all{Auch Kamerun kämpft gegen Fälschung und Schmuggel und ruft die Verbraucher zur Mitarbeit auf.}
\end{itemize}
\end{exemplebox}

\begin{attention}[Erreurs fréquentes des francophones]
\begin{itemize}
\item \textbf{Majuscules :} tous les noms prennent une majuscule : \all{die Korruption, der Schmuggel, die Strafe} — jamais \all{die korruption}.
\item \textbf{Genre des noms :} apprenez chaque mot avec son article : \emph{die} Ausschreibung, \emph{der} Zuschlag, \emph{die} Fälschung, \emph{der} Betrug, \emph{der} Schmuggel.
\item \textbf{Faux ami :} \all{die Strafe} = la peine/sanction (et non « la souffrance ») ; « une amende » se dit \all{die Geldstrafe}.
\item \textbf{Place du verbe :} dans une subordonnée avec \all{dass}, le verbe va à la fin : \all{Die Opposition behauptet, dass ein Unternehmen bevorzugt \emph{wurde}.}
\item \textbf{Ne traduisez pas} « faire passer en contrebande » mot à mot : le verbe allemand est simplement \all{schmuggeln}.
\end{itemize}
\end{attention}

\begin{savaistu}[Transparency International : une ONG berlinoise]
L'Allemagne publie chaque année un rapport sur la corruption, et la police criminelle fédérale (\all{das Bundeskriminalamt}) consacre un chapitre entier à la \all{Wirtschaftskriminalität}. Saviez-vous que \textbf{Transparency International}, la plus célèbre ONG de lutte contre la corruption, a son siège à Berlin ? Elle publie chaque année un indice mondial de perception de la corruption, dans lequel tous les pays — y compris le Cameroun — sont classés. Le mot \all{Schmiergeld} (pot-de-vin) vient de \all{schmieren}, « graisser » : on « graisse la patte » en allemand comme en français !
\end{savaistu}

\begin{experience}[À vous de parler : Nachrichtensprecher]
Par groupes de trois : l'un joue le journaliste qui présente le fait divers, l'autre le ministre qui se défend (\all{Das Verfahren war korrekt!}), le troisième l'expert qui commente. Puis le journaliste rapporte les paroles au discours indirect : \all{Der Minister erklärte, das Verfahren sei korrekt gewesen.} Chaque groupe présente sa scène à la classe.
\end{experience}

\coursec{Lexique thématique : marchés publics et criminalité économique}

Dans la section précédente, nous avons lu et compris globalement le texte support sur les marchés publics et la criminalité économique. Vous avez remarqué que ce texte contient beaucoup de mots nouveaux : \all{die Ausschreibung}, \all{der Zuschlag}, \all{der Schmuggel}, \all{die Fälschung}... Cette deuxième section est consacrée à l'étude systématique de ce vocabulaire. Un bon lexique se construit toujours de la même manière : on observe le mot dans une phrase, on apprend son \cle{article} et son \cle{pluriel}, on le rattache à sa \cle{famille de mots}, puis on le réemploie. C'est exactement la démarche que nous allons suivre, champ lexical par champ lexical.

\courssub{Observation : un court texte pour rencontrer le vocabulaire en contexte}

\begin{textebox}[Korruption im Wirtschaftsleben — ein Zeitungsbericht]
In der Stadt Douala wurde gestern ein \textbf{großer} Skandal bekannt. Eine Firma hatte bei einer öffentlichen Ausschreibung den Zuschlag für den Bau einer Schule erhalten. Jetzt ermittelt die Polizei, weil der Auftragnehmer angeblich Bestechungsgelder an einen Beamten gezahlt hat. „Wir kämpfen entschlossen gegen die Korruption“, erklärte ein Vertreter der Regierung. „Wer bestecht, steht unter Strafe.“

Auch an der Grenze gibt es Probleme. Der Zoll hat letzte Woche mehrere Lastwagen kontrolliert und gefälschte Medikamente beschlagnahmt. Die Schmuggler wollten die Waren illegal einführen. Gefälschte Produkte und Plagiate sind eine Gefahr für die Gesundheit der Bevölkerung. Ein Händler wurde angeklagt und später zu zwei Jahren Gefängnis und einer hohen Geldstrafe verurteilt. Das Gesetz ist in diesem Punkt sehr streng. Experten warnen: Betrug, Schmuggel und Fälschungen schaden der ganzen Wirtschaft des Landes. Deshalb müssen Bürger, Firmen und Staat zusammenarbeiten, damit öffentliche Aufträge fair vergeben werden und der Handel ehrlich bleibt.
\end{textebox}

\textbf{Glossaire} : \all{der Zuschlag} : l'adjudication ; \all{der Auftragnehmer} : le soumissionnaire, l'entreprise attributaire ; \all{die Bestechungsgelder} (pl.) : les pots-de-vin ; \all{beschlagnahmen} : saisir, confisquer ; \all{illegal einführen} : importer illégalement ; \all{das Plagiat} : le produit contrefait ; \all{unter Strafe stehen} : être passible d'une peine.

\textbf{Questions de compréhension rapide} :
\begin{enumerate}
\item \all{Wer hat den Zuschlag erhalten?} (Qui a obtenu l'adjudication ?)
\item \all{Was hat der Zoll beschlagnahmt?} (Qu'est-ce que la douane a saisie ?)
\item \all{Wozu wurde der Händler verurteilt?} (À quoi le commerçant a-t-il été condamné ?)
\end{enumerate}

\courssub{Champ 1 — Les marchés publics}

\begin{definition}[Vocabulaire des marchés publics]
\begin{itemize}
\item \all{der Markt, die Märkte} : le marché (au sens économique et concret)
\item \all{der öffentliche Auftrag, die öffentlichen Aufträge} : le marché public, la commande publique
\item \all{die Ausschreibung, -en} : l'appel d'offres
\item \all{ausschreiben — schrieb aus — hat ausgeschrieben} (verbe fort, \textbf{séparable}) : lancer un appel d'offres
\item \all{das Angebot, -e} : l'offre, la soumission
\item \all{der Zuschlag, die Zuschläge} : l'adjudication (l'attribution du marché)
\item \all{zuschlagen — schlug zu — hat zugeschlagen} (verbe fort, \textbf{séparable}) : ici, \all{den Zuschlag erhalten} (obtenir l'adjudication) / \all{den Zuschlag erteilen} (attribuer l'adjudication)
\item \all{der Auftragnehmer, -} : l'attributaire, le soumissionnaire retenu
\item \all{die Firma, die Firmen} / \all{das Unternehmen, -} : l'entreprise
\item \all{vergeben — vergab — hat vergeben} (verbe fort, \textbf{inséparable}) : attribuer, octroyer (un marché)
\end{itemize}
\end{definition}

\begin{exemplebox}[Le vocabulaire en phrases]
\all{Die Regierung schreibt den Bau einer Brücke öffentlich aus.} « L'État lance un appel d'offres public pour la construction d'un pont. »

\all{Zehn Firmen geben ein Angebot ab.} « Dix entreprises déposent une offre. »

\all{Die Firma hat den Zuschlag erhalten.} « L'entreprise a obtenu l'adjudication. »

\all{Der Staat vergibt einen Auftrag an das beste Unternehmen.} « L'État attribue un marché à la meilleure entreprise. »
\end{exemplebox}

\begin{aretenir}[Expression clé]
\all{einen Auftrag vergeben} = attribuer un marché. Le verbe \all{vergeben} est un verbe fort irrégulier : \all{vergibt — vergab — hat vergeben}. Exemple : \all{Der Minister hat den Auftrag an eine deutsche Firma vergeben.} « Le ministre a attribué le marché à une entreprise allemande. »
\end{aretenir}

\courssub{Champ 2 — La corruption}

\begin{definition}[Vocabulaire de la corruption]
\begin{itemize}
\item \all{die Korruption} (sans pluriel, indénombrable) : la corruption
\item \all{die Bestechung, -en} : la corruption, l'acte de corrompre
\item \all{die Bestechungsgelder} (pluriel uniquement) : les pots-de-vin
\item \all{bestechen — bestach — hat bestochen} (verbe fort) : corrompre, soudoyer
\item \all{korrupt} (adjectif) : corrompu ; \all{bestechlich} : corruptible
\item \all{der Schmiergeld-Skandal, -e} : le scandale des pots-de-vin (\all{das Schmiergeld} : le pot-de-vin, littéralement « l'argent graisse »)
\end{itemize}
\end{definition}

\begin{exemplebox}[Exemples traduits]
\all{Der Beamte wurde bestochen.} « Le fonctionnaire a été corrompu. »

\all{Korruption zerstört das Vertrauen der Bürger.} « La corruption détruit la confiance des citoyens. »

\all{Wir müssen gegen die Korruption kämpfen.} « Nous devons lutter contre la corruption. »
\end{exemplebox}

\courssub{Champ 3 — La contrebande}

\begin{definition}[Vocabulaire de la contrebande]
\begin{itemize}
\item \all{der Schmuggel} (indénombrable) : la contrebande
\item \all{schmuggeln — schmuggelte — hat geschmuggelt} (verbe faible) : faire passer en contrebande
\item \all{der Schmuggler, -} : le contrebandier
\item \all{der Zoll, die Zölle} : la douane (attention : \all{die Zölle} = les droits de douane)
\item \all{die Grenze, -n} : la frontière
\item \all{beschlagnahmen — beschlagnahmte — hat beschlagnahmt} (verbe faible, \textbf{inséparable} [préfixe \all{be-}]) : saisir, confisquer
\item \all{illegal einführen — führte illegal ein — hat illegal eingeführt} (verbe fort, \textbf{séparable}) : importer illégalement
\end{itemize}
\end{definition}

\begin{exemplebox}[Exemples traduits]
\all{Die gefälschten Medikamente wurden beschlagnahmt.} « Les médicaments falsifiés ont été saisis. »

\all{Der Zoll kontrolliert die Lastwagen an der Grenze.} « La douane contrôle les camions à la frontière. »

\all{Die Schmuggler haben Zigaretten illegal eingeführt.} « Les contrebandiers ont importé des cigarettes illégalement. »
\end{exemplebox}

\courssub{Champ 4 — La falsification et la fraude}

\begin{definition}[Vocabulaire de la falsification]
\begin{itemize}
\item \all{die Fälschung, -en} : la falsification, le faux (le produit falsifié)
\item \all{fälschen — fälschte — hat gefälscht} (verbe faible) : falsifier, contrefaire
\item \all{gefälscht} (adjectif/participe) : falsifié, contrefait ; \all{das gefälschte Produkt} : le produit falsifié
\item \all{die Plagiate} (pl.) : les contrefaçons, les produits copiés (sing. \all{das Plagiat})
\item \all{der Betrug} (indénombrable) : l'escroquerie, la fraude
\item \all{betrügen — betrog — hat betrogen} (verbe fort) : tromper, escroquer
\item \all{der Betrüger, -} : l'escroc, le fraudeur
\end{itemize}
\end{definition}

\begin{attention}[Faux amis : \all{der Betrug} et \all{die Fälschung}]
\all{der Betrug} ne signifie pas « le bruit » mais « l'escroquerie, la fraude ». Le bruit se dit \all{der Lärm}. De même, \all{die Fälschung} désigne la falsification ou le faux (objet contrefait), et non « la fausseté » au sens moral (qui se dirait plutôt \all{die Falschheit}). Enfin, ne confondez pas \all{fälschen} (falsifier) et \all{fallen} (tomber) : \all{Er fälscht Dokumente} « Il falsifie des documents » / \all{Er fällt} « Il tombe ».
\end{attention}

\courssub{Champ 5 — Justice et sanctions}

\begin{definition}[Vocabulaire de la justice]
\begin{itemize}
\item \all{das Gesetz, -e} : la loi
\item \all{die Strafe, -n} : la peine, la sanction ; \all{die Geldstrafe, -n} : l'amende
\item \all{das Gefängnis, -se} : la prison
\item \all{ermitteln — ermittelte — hat ermittelt} (verbe faible, intransitif, souvent avec \all{gegen} + Akk.) : enquêter, mener une enquête (\all{die Polizei ermittelt} : la police enquête)
\item \all{anklagen — klagte an — hat angeklagt} (verbe faible, \textbf{séparable}) : accuser, inculper
\item \all{verurteilen (zu + Dativ) — verurteilte — hat verurteilt} (verbe faible, \textbf{inséparable} [préfixe \all{ver-}]) : condamner (à)
\end{itemize}
\end{definition}

\begin{exemplebox}[Exemples traduits]
\all{Der Händler wurde zu zwei Jahren Gefängnis verurteilt.} « Le commerçant a été condamné à deux ans de prison. »

\all{Wer schmuggelt, steht unter Strafe.} « Qui fait de la contrebande est passible d'une peine. »

\all{Die Staatsanwaltschaft klagt den Betrüger an.} « Le parquet inculpe l'escroc. »
\end{exemplebox}

\courssub{Familles de mots : construire son vocabulaire en réseau}

En allemand, les mots d'une même famille partagent une racine. Apprendre la famille entière multiplie votre vocabulaire sans effort supplémentaire.

\begin{center}
\begin{tabularx}{\linewidth}{|l|l|l|X|}
\hline
\rowcolor{popblueL} Verbe & Nom (action) & Nom (personne) & Adjectif \\
\hline
fälschen (falsifier) & die Fälschung & der Fälscher & gefälscht \\
\hline
schmuggeln (contrebande) & der Schmuggel & der Schmuggler & geschmuggelt \\
\hline
bestechen (corrompre) & die Bestechung & der Bestecher & bestechlich \\
\hline
betrügen (escroquer) & der Betrug & der Betrüger & betrügerisch \\
\hline
\end{tabularx}
\end{center}

\begin{propriete}[Les noms en \cle{-ung} sont féminins]
Tous les noms formés avec le suffixe \cle{-ung} sont \textbf{féminins} (\all{die}) et forment leur pluriel en \cle{-en} : \all{die Ausschreibung, -en} ; \all{die Bestechung, -en} ; \all{die Fälschung, -en} ; \all{die Ermittlung, -en} (l'enquête) ; \all{die Verurteilung, -en} (la condamnation). Ce suffixe transforme un verbe en nom d'action, comme le français « -ation » : \all{fälschen} → \all{die Fälschung} (falsifier → la falsification).
\end{propriete}

\begin{popfigure}
\begin{tikzpicture}[mindmap, grow cyclic, every node/.style={concept}, root concept/.append style={concept color=popblue, font=\small\bfseries}, level 1/.append style={level distance=3.4cm, sibling angle=90, concept color=poporangeL}, level 2/.append style={level distance=2.2cm, sibling angle=45, concept color=popgreenL, font=\scriptsize}]
\node[font=\small\bfseries]{Wirtschaftskriminalität}
  child { node {Markt} child { node {die Ausschreibung} } child { node {der Zuschlag} } }
  child { node {Korruption} child { node {bestechen} } child { node {die Bestechung} } }
  child { node {Schmuggel} child { node {der Zoll} } child { node {beschlagnahmen} } }
  child { node {Fälschung} child { node {der Betrug} } child { node {die Strafe} } };
\end{tikzpicture}
\legende{Carte mentale du champ lexical : la criminalité économique et ses quatre domaines.}
\end{popfigure}

\courssub{Tableau récapitulatif : articles et pluriels à compléter}

\begin{center}
\begin{tabularx}{\linewidth}{|l|l|X|}
\hline
\rowcolor{popblueL} Deutsch (artikel + Plural) & Français & Famille \\
\hline
der Markt, die Märkte & le marché & — \\
\hline
die Ausschreibung, -en & l'appel d'offres & ausschreiben \\
\hline
der Zuschlag, die Zuschläge & l'adjudication & zuschlagen \\
\hline
die Korruption (unz.) & la corruption & korrupt \\
\hline
die Bestechung, -en & la corruption (acte) & bestechen \\
\hline
der Schmuggel (unz.) & la contrebande & schmuggeln \\
\hline
der Zoll, die Zölle & la douane & — \\
\hline
die Grenze, -n & la frontière & — \\
\hline
die Fälschung, -en & la falsification & fälschen \\
\hline
der Betrug (unz.) & l'escroquerie & betrügen \\
\hline
das Gesetz, -e & la loi & gesetzlich \\
\hline
die Geldstrafe, -n & l'amende & die Strafe \\
\hline
das Gefängnis, -se & la prison & — \\
\hline
\end{tabularx}
\end{center}

\begin{attention}[Pièges fréquents des francophones]
\begin{itemize}
\item N'oubliez jamais la majuscule des noms : on écrit \all{die Grenze}, \all{der Zoll}, jamais \all{die grenze}.
\item \all{der Zoll} au pluriel \all{die Zölle} prend un Umlaut : mémorisez \all{der Zoll, die Zölle}.
\item \all{die Bestechungsgelder} n'existe qu'au pluriel, comme « les pots-de-vin » ; on ne dit pas un singulier courant.
\item Les noms indénombrables (\all{die Korruption}, \all{der Schmuggel}, \all{der Betrug}) ne prennent pas de pluriel : ne dites jamais \all{die Korruptionen}.
\item Verbes à particule : \all{ausschreiben}, \all{anklagen}, \all{illegal einführen} sont \textbf{séparables} (particule en fin de proposition principale) ; \all{beschlagnahmen}, \all{vergeben}, \all{verurteilen} sont \textbf{inséparables} (pas de \all{ge-} au Partizip II).
\end{itemize}
\end{attention}

\begin{savaistu}[Le saviez-tu ?]
En Suisse, l'appel d'offres se dit aussi \all{die Submission}, un mot hérité du latin que l'on rencontre dans les textes administratifs suisses ; en Autriche, on parle volontiers de \all{die Ausschreibung} comme en Allemagne. Quant au mot \all{Schmuggel}, il vient d'un vieux verbe germanique qui signifiait « se glisser, se faufiler » : le contrebandier est donc, étymologiquement, celui qui « se glisse » à travers la frontière ! Au Cameroun, les douanes (\all{der Zoll}) luttent activement contre la contrebande aux frontières, par exemple sur les axes vers le Nigeria et le Tchad.
\end{savaistu}

\begin{exoresolu}[Complétez avec le mot qui convient]
Énoncé. Complétez les phrases avec : \all{Zuschlag}, \all{beschlagnahmt}, \all{Bestechungsgelder}, \all{Geldstrafe}, \all{gefälschten}.

1. \all{Die Firma Müller hat den \ldots\ für den Straßenbau erhalten.}
2. \all{Der Zoll hat die \ldots\ Waren an der Grenze \ldots .}
3. \all{Der Beamte hat \ldots\ angenommen.}
4. \all{Der Schmuggler muss eine hohe \ldots\ zahlen.}
\tcblower
Solution détaillée.

1. \all{Die Firma Müller hat den \textbf{Zuschlag} für den Straßenbau erhalten.} « L'entreprise Müller a obtenu l'adjudication pour la construction de la route. » (accusatif masculin : \all{den Zuschlag}).

2. \all{Der Zoll hat die \textbf{gefälschten} Waren an der Grenze \textbf{beschlagnahmt}.} « La douane a saisi les marchandises falsifiées à la frontière. » (adjectif épithète au pluriel avec article défini : terminaison \all{-en} ; participe II de \all{beschlagnahmen} en fin de phrase au Perfekt, sans \all{ge-} car inséparable).

3. \all{Der Beamte hat \textbf{Bestechungsgelder} angenommen.} « Le fonctionnaire a accepté des pots-de-vin. » (nom qui n'existe qu'au pluriel).

4. \all{Der Schmuggler muss eine hohe \textbf{Geldstrafe} zahlen.} « Le contrebandier doit payer une amende élevée. » (accusatif féminin : \all{eine hohe Geldstrafe}).
\end{exoresolu}

\begin{exercice}[À vous de jouer]
1. Traduisez en allemand : a) L'État attribue le marché à la meilleure entreprise. b) Les médicaments falsifiés ont été saisis par la douane. c) Il a été condamné à une amende.

2. Complétez le tableau : donnez l'article et le pluriel de \all{Grenze}, \all{Gesetz}, \all{Ausschreibung}, \all{Strafe}, \all{Zoll}.

3. Formez la famille de mots complète (verbe, nom d'action, nom de personne, adjectif) de : \all{bestechen} et \all{fälschen}.
\end{exercice}

\begin{corrige}[Exercice « À vous de jouer »]
1. a) \all{Der Staat vergibt den Auftrag an die beste Firma.} b) \all{Die gefälschten Medikamente wurden vom Zoll beschlagnahmt.} c) \all{Er wurde zu einer Geldstrafe verurteilt.}

2. \all{die Grenze, -n} ; \all{das Gesetz, -e} ; \all{die Ausschreibung, -en} ; \all{die Strafe, -n} ; \all{der Zoll, die Zölle}.

3. \all{bestechen → die Bestechung, der Bestecher, bestechlich} ; \all{fälschen → die Fälschung, der Fälscher, gefälscht}.
\end{corrige}

\begin{aretenir}[L'essentiel du lexique]
Retenez les cinq verbes forts ou clés du thème : \all{vergeben — vergab — vergeben} (attribuer), \all{bestechen — bestach — bestochen} (corrompre), \all{betrügen — betrog — betrogen} (escroquer), \all{fälschen} et \all{schmuggeln} (faibles), ainsi que les trois expressions : \all{einen Auftrag vergeben} (attribuer un marché), \all{gegen die Korruption kämpfen} (lutter contre la corruption), \all{unter Strafe stehen} (être passible d'une peine). Ce vocabulaire vous servira dans les sections de grammaire (\all{Zustandspassiv}, \all{Konjunktiv I}) et de production qui suivent.
\end{aretenir}

\coursec{Grammaire 1 : le passif d'état (Zustandspassiv)}

Dans la section précédente, nous avons constitué le lexique des marchés publics et de la criminalité économique : \all{die Ausschreibung}, \all{die Korruption}, \all{der Schmuggel}, \all{die Fälschung}, \all{der Betrug}. En lisant les articles de presse et les rapports officiels qui emploient ce vocabulaire, on remarque une tournure qui revient constamment : \all{Die Medikamente sind gefälscht.} Cette construction, qui ressemble au passif mais n'en est pas tout à fait un, s'appelle le \cle{Zustandspassiv}, le passif d'état. C'est l'une des structures les plus fréquentes de la langue administrative et journalistique allemande — et l'une des plus mal maîtrisées par les francophones, qui la confondent avec le passif d'action.

\courssub{Observation : deux phrases, deux réalités}

Observons attentivement les exemples tirés de notre texte support et de la presse germanophone :

\begin{exemplebox}[Phrases à observer]
\begin{enumerate}
\item \all{Die Medikamente sind gefälscht.} « Les médicaments sont falsifiés. »
\item \all{Das Verfahren ist abgeschlossen.} « La procédure est close. »
\item \all{Die Grenze war gesichert.} « La frontière était sécurisée. »
\item \all{Die Akte ist geschlossen.} « Le dossier est clos. »
\item \all{Der Schmuggler ist verurteilt.} « Le contrebandier est condamné. »
\end{enumerate}
\end{exemplebox}

Questions d'observation : quel est l'auxiliaire employé ? S'agit-il de \all{werden} ou de \all{sein} ? Ces phrases décrivent-elles une action qui se déroule sous nos yeux, ou le résultat d'une action déjà terminée ?

Réponse : l'auxiliaire est \all{sein} (et non \all{werden}), et chaque phrase décrit un \textbf{état}, un constat : la falsification a eu lieu, elle est achevée, et ce qui nous intéresse, c'est le résultat visible maintenant. Comparons avec le passif que vous connaissez déjà :

\begin{center}
\begin{tabularx}{\linewidth}{|X|X|X|}
\hline
\rowcolor{popblueL} Vorgangspassiv (action) & Zustandspassiv (état) & Remarque \\
\hline
\all{Die Medikamente werden gefälscht.} & \all{Die Medikamente sind gefälscht.} & action en cours vs résultat \\
\hline
\all{Das Verfahren wird abgeschlossen.} & \all{Das Verfahren ist abgeschlossen.} & on clôt vs c'est clos \\
\hline
\all{Die Grenze wurde gesichert.} & \all{Die Grenze war gesichert.} & on a sécurisé vs c'était sécurisé \\
\hline
\end{tabularx}
\end{center}

\courssub{Formation du Zustandspassiv}

\begin{propriete}[Règle de formation]
Le \cle{Zustandspassiv} (passif d'état) se forme avec :
\begin{center}
\textbf{sein (conjugué au temps voulu) + Partizip II}
\end{center}
Il exprime l'\textbf{état résultant d'une action terminée} : \all{Die Akte ist geschlossen.} « Le dossier est clos. »

Le \cle{Vorgangspassiv} (passif d'action), lui, se forme avec \textbf{werden + Partizip II} et décrit l'action elle-même : \all{Die Akte wird geschlossen.} « On clôt le dossier. »
\end{propriete}

La distinction est donc entièrement portée par l'auxiliaire : \all{werden} met le projecteur sur le \textbf{processus}, \all{sein} sur le \textbf{résultat}. Le participe passé, dans le Zustandspassiv, fonctionne presque comme un adjectif : \all{gefälscht} dans \all{Die Ware ist gefälscht} décrit la qualité de la marchandise, exactement comme \all{Die Ware ist teuer}. D'ailleurs, placé devant le nom, ce participe se décline comme un adjectif : \all{die gefälschte Ware} « la marchandise falsifiée », \all{das abgeschlossene Verfahren} « la procédure close », \all{die beschlagnahmten Waren} « les marchandises saisies ».

\begin{exemplebox}[Paires action / état]
\begin{itemize}
\item \all{Das Paket wird geöffnet.} « On ouvre le paquet » (action). $\rightarrow$ \all{Das Paket ist geöffnet.} « Le paquet est ouvert » (état).
\item \all{Der Schmuggler wurde verurteilt.} « Le contrebandier a été condamné » (action). $\rightarrow$ \all{Er ist verurteilt.} « Il est condamné » (état).
\item \all{Die Ausschreibung wird veröffentlicht.} « L'appel d'offres est publié » (on le publie maintenant). $\rightarrow$ \all{Die Ausschreibung ist veröffentlicht.} « L'appel d'offres est publié » (il figure déjà au journal officiel).
\item \all{Die Waren werden beschlagnahmt.} « Les marchandises sont saisies » (opération en cours). $\rightarrow$ \all{Die Waren sind beschlagnahmt.} « Les marchandises sont saisies » (elles sont sous scellés).
\end{itemize}
\end{exemplebox}

\begin{popfigure}
\begin{tikzpicture}[
  node distance=0.6cm and 2.2cm,
  box/.style={draw=popline, rounded corners=3pt, align=center, inner sep=6pt, font=\small},
  ]
\node[box, fill=popcream, font=\bfseries] (titre) {Le passif en allemand : deux constructions};
\node[box, fill=poporangeL, below left=0.8cm and -1.5cm of titre, text width=5.6cm] (vorgang)
  {\textbf{Vorgangspassiv} (action)\\[2pt]
   \all{werden + Partizip II}\\[2pt]
   \all{Die Ware \textbf{wird} gefälscht.}\\
   « On falsifie la marchandise. »\\[2pt]
   \emph{flèche dynamique : le processus}};
\node[box, fill=popblueL, below right=0.8cm and -1.5cm of titre, text width=5.6cm] (zustand)
  {\textbf{Zustandspassiv} (état)\\[2pt]
   \all{sein + Partizip II}\\[2pt]
   \all{Die Ware \textbf{ist} gefälscht.}\\
   « La marchandise est falsifiée. »\\[2pt]
   \emph{cadre statique : le résultat}};
\draw[-{Stealth[length=3mm]}, thick, poporange] (titre.south west) ++(0.3,0) -- (vorgang.north);
\draw[-{Stealth[length=3mm]}, thick, popblue] (titre.south east) ++(-0.3,0) -- (zustand.north);
\draw[-{Stealth[length=3mm]}, thick, poporange, decorate, decoration={snake, amplitude=0.4mm}] ($(vorgang.south west)+(0.4,-0.15)$) -- ($(vorgang.south east)+(-0.4,-0.15)$);
\draw[thick, popblue] ($(zustand.south west)+(0.2,-0.35)$) rectangle ($(zustand.south east)+(-0.2,-0.05)$);
\end{tikzpicture}
\legende{Action (werden, flèche ondulée) contre état (sein, cadre fermé)}
\end{popfigure}

\courssub{Les temps du Zustandspassiv}

Le Zustandspassiv se conjugue aux temps de \all{sein}. Le tableau suivant donne les quatre temps courants, avec le modèle \all{fälschen} :

\begin{center}
\begin{tabular}{|l|l|l|}
\hline
\rowcolor{popblueL} Temps & Forme & Traduction \\
\hline
Präsens & \all{Die Ware ist gefälscht.} & La marchandise est falsifiée. \\
\hline
Präteritum & \all{Die Ware war gefälscht.} & La marchandise était falsifiée. \\
\hline
Perfekt & \all{Die Ware ist gefälscht gewesen.} & La marchandise a été falsifiée (état passé). \\
\hline
Plusquamperfekt & \all{Die Ware war gefälscht gewesen.} & La marchandise avait été falsifiée (état antérieur). \\
\hline
\end{tabular}
\end{center}

\begin{attention}[Le Perfekt et le Plusquamperfekt]
Aux temps composés, le participe de \all{sein} est \textbf{toujours} \all{gewesen} (et jamais \all{geworden}, qui appartient au Vorgangspassiv). Comparez :
\begin{itemize}
\item \all{Die Ware ist gefälscht \textbf{gewesen}.} (Zustandspassiv : elle était falsifiée.)
\item \all{Die Ware ist gefälscht \textbf{worden}.} (Vorgangspassiv : elle a été falsifiée — l'action.)
\end{itemize}
Une seule lettre distingue les deux sens : retenez le couple \all{gewesen} = état, \all{worden} = action.
\end{attention}

Conjugaison complète au présent, avec le modèle \all{abschließen} (forme impersonnelle à valeur d'exemple) :

\begin{center}
\begin{tabular}{|l|l|l|}
\hline
\rowcolor{popblueL} Personne & Zustandspassiv & Vorgangspassiv \\
\hline
ich & \all{bin abgeschlossen} & \all{werde abgeschlossen} \\
\hline
du & \all{bist abgeschlossen} & \all{wirst abgeschlossen} \\
\hline
er/sie/es & \all{ist abgeschlossen} & \all{wird abgeschlossen} \\
\hline
wir & \all{sind abgeschlossen} & \all{werden abgeschlossen} \\
\hline
ihr & \all{seid abgeschlossen} & \all{werdet abgeschlossen} \\
\hline
sie/Sie & \all{sind abgeschlossen} & \all{werden abgeschlossen} \\
\hline
\end{tabular}
\end{center}

\courssub{Emplois : la langue du constat}

Le Zustandspassiv est la forme privilégiée de trois types de discours :

\begin{enumerate}
\item \textbf{Les constats administratifs} : \all{Der Auftrag ist vergeben.} « Le marché est attribué. » \all{Die Frist ist abgelaufen.} « Le délai est expiré. » \all{Das Konto ist gesperrt.} « Le compte est bloqué. »
\item \textbf{Les comptes rendus journalistiques} : \all{Nach der Razzia sind zwanzig Verdächtige festgenommen.} « Après la descente de police, vingt suspects sont arrêtés (sous les verrous). »
\item \textbf{Les descriptions de résultats} : \all{Die Grenze ist gesichert, die Lager sind versiegelt.} « La frontière est sécurisée, les entrepôts sont scellés. »
\end{enumerate}

Dans les rapports officiels — audits des marchés publics, rapports de la police douanière, communiqués des ministères —, le Zustandspassiv domine, car ce qui compte juridiquement, c'est l'état de fait, pas le déroulement de l'action.

\begin{methode}[Comment choisir entre werden et sein ?]
\begin{enumerate}
\item Posez-vous la question : la phrase décrit-elle une \textbf{action en cours} ou un \textbf{résultat acquis} ?
\item Testez avec un adverbe : si l'on peut ajouter \all{gerade} (« en train de »), c'est le \textbf{Vorgangspassiv} : \all{Die Ware wird gerade kontrolliert.} « La marchandise est en train d'être contrôlée. »
\item Si l'on peut ajouter \all{schon} ou \all{bereits} (« déjà »), c'est le \textbf{Zustandspassiv} : \all{Die Ware ist bereits kontrolliert.} « La marchandise est déjà contrôlée. »
\item En cas de doute au passé, demandez-vous si le français dirait « a été » (action $\rightarrow$ \all{wurde}) ou « était » (état $\rightarrow$ \all{war}).
\end{enumerate}
\end{methode}

\courssub{Comparaison avec le français : le piège de « est »}

Le français utilise la même forme « est + participe » pour l'action et l'état : « les produits sont falsifiés » peut signifier « on les falsifie » ou « ils sont dans un état falsifié ». L'allemand, lui, \textbf{oblige à choisir} :

\begin{center}
\begin{tabularx}{\linewidth}{|X|X|X|}
\hline
\rowcolor{popblueL} Français & Deutsch & Sens \\
\hline
Les produits sont falsifiés (constat). & \all{Die Produkte sind gefälscht.} & état \\
\hline
Les produits sont falsifiés (par des bandits). & \all{Die Produkte werden (von Banditen) gefälscht.} & action \\
\hline
Le dossier est clos. & \all{Die Akte ist geschlossen.} & état \\
\hline
Le dossier est clos chaque soir à 17 h. & \all{Die Akte wird jeden Abend geschlossen.} & action répétée \\
\hline
\end{tabularx}
\end{center}

\begin{attention}[Erreur fréquente des francophones]
Ne traduisez jamais automatiquement « est falsifié » par \all{wird gefälscht}. Le français « est » cache les deux constructions allemandes. Demandez-vous toujours : \textbf{action en cours} (\all{werden}) ou \textbf{résultat} (\all{sein}) ? Ainsi « Le marché est attribué » (décision prise, constat) se dit \all{Der Auftrag ist vergeben}, et non \all{wird vergeben}, qui signifierait « on est en train de l'attribuer ».
\end{attention}

\begin{propriete}[Limites du Zustandspassiv]
\begin{itemize}
\item \textbf{Pas de futur courant} : on évite \all{Die Akte wird geschlossen sein} dans la langue courante ; on préfère une périphrase : \all{Dann ist die Akte geschlossen.} ou \all{Bis morgen ist die Akte geschlossen.}
\item \textbf{Verbes intransitifs sans résultat visible} : les verbes qui ne produisent pas d'état durable n'admettent pas le Zustandspassiv. On dit \all{Es wird getanzt} (« on danse »), mais jamais \all{Es ist getanzt} pour un état.
\item Le verbe doit pouvoir former un Vorgangspassiv : seuls les verbes transitifs (ou à passif impersonnel) entrent dans le système.
\end{itemize}
\end{propriete}

\courssub{Texte d'application : un rapport de la douane}

\begin{textebox}[Zollbericht aus Douala]
Bei einer Kontrolle im Hafen von Douala sind gestern drei Container geöffnet worden. Das Ergebnis ist alarmierend: Ein großer Teil der Medikamente ist gefälscht. Die Packungen sind professionell bedruckt, aber der Inhalt ist wirkungslos. Zwei Verdächtige sind festgenommen, die Waren sind beschlagnahmt und die Lagerhalle ist versiegelt. Nach Angaben der Zollbeamten war die Grenze im Norden monatelang schlecht gesichert; deshalb war der Schmuggel so einfach gewesen. Jetzt ist das Ermittlungsverfahren eröffnet. Ein Sprecher erklärte: „Die Akte ist noch nicht geschlossen. Weitere Kontrollen sind geplant.“ Experten warnen: Solange gefälschte Produkte auf dem Markt sind, ist die Gesundheit der Bevölkerung gefährdet. Der Fall zeigt, wie wichtig transparente öffentliche Aufträge sind: Wenn die Ausschreibung manipuliert ist, ist der Betrug schon programmiert.
\end{textebox}

\textbf{Glossaire} : der Container, die Container (conteneur) ; die Packung, die Packungen (emballage) ; bedruckt (imprimé) ; wirkungslos (inefficace) ; festnehmen (arrêter) ; beschlagnahmen (saisir) ; versiegeln (sceller) ; der Zollbeamte, die Zollbeamten (douanier) ; das Ermittlungsverfahren, die Ermittlungsverfahren (enquête) ; gefährden (mettre en danger) ; der öffentliche Auftrag, die öffentlichen Aufträge (marché public) ; manipulieren (manipuler).

\textbf{Questions de compréhension} :
\begin{enumerate}
\item Quels éléments du texte sont au Zustandspassiv ? Relevez-en cinq.
\item Pourquoi le journaliste écrit-il \all{sind geöffnet worden} (Vorgangspassiv) pour les containers, mais \all{sind versiegelt} (Zustandspassiv) pour l'entrepôt ?
\item Traduisez : \all{Wenn die Ausschreibung manipuliert ist, ist der Betrug schon programmiert.}
\end{enumerate}

\begin{corrige}[Questions de compréhension sur le texte]
\begin{enumerate}
\item Exemples de Zustandspassiv : \all{ist gefälscht} ; \all{sind bedruckt} ; \all{sind festgenommen} ; \all{sind beschlagnahmt} ; \all{ist versiegelt} ; \all{ist eröffnet} ; \all{ist geschlossen} ; \all{sind geplant} ; \all{ist gefährdet} ; \all{ist manipuliert}.
\item \all{Sind geöffnet worden} est un Vorgangspassiv au Perfekt : le journaliste rapporte l'\textbf{action} d'ouverture menée hier par les douaniers. \all{Ist versiegelt} est un Zustandspassiv au présent : il décrit l'\textbf{état actuel} de l'entrepôt, désormais sous scellés.
\item « Si l'appel d'offres est manipulé, l'escroquerie est déjà programmée. » (Les deux propositions sont au Zustandspassiv : états constatés.)
\end{enumerate}
\end{corrige}

\courssub{Exercices d'application}

\begin{exoresolu}[Transformation : Vorgangspassiv $\rightarrow$ Zustandspassiv]
Énoncé : transformez les phrases suivantes en passif d'état au présent. 1) \all{Die Grenze wird gesichert.} 2) \all{Der Schmuggler wird verurteilt.} 3) \all{Die gefälschten Waren werden beschlagnahmt.}
\tcblower
Solution détaillée : on remplace l'auxiliaire \all{werden} par \all{sein} conjugué au présent, le participe reste inchangé. 1) \all{Die Grenze \textbf{ist} gesichert.} « La frontière est sécurisée. » 2) \all{Der Schmuggler \textbf{ist} verurteilt.} « Le contrebandier est condamné. » 3) \all{Die gefälschten Waren \textbf{sind} beschlagnahmt.} « Les marchandises falsifiées sont saisies. » Attention à l'accord du verbe \all{sein} avec le sujet pluriel à la phrase 3.
\end{exoresolu}

\begin{exoresolu}[Choix entre sein et werden]
Énoncé : complétez avec la bonne forme de \all{sein} ou \all{werden}. 1) \all{Die Medikamente \dots\ gerade kontrolliert.} 2) \all{Das Verfahren \dots\ seit gestern abgeschlossen.} 3) \all{Jede Woche \dots\ neue Fälschungen entdeckt.}
\tcblower
Solution détaillée : 1) \all{werden} : l'adverbe \all{gerade} (« en train de ») signale l'action en cours $\rightarrow$ Vorgangspassiv. 2) \all{ist} : \all{seit gestern} décrit un état qui dure $\rightarrow$ Zustandspassiv. 3) \all{werden} : \all{jede Woche} exprime une action répétée $\rightarrow$ Vorgangspassiv. Traductions : 1) « Les médicaments sont en train d'être contrôlés. » 2) « La procédure est close depuis hier. » 3) « Chaque semaine, de nouvelles contrefaçons sont découvertes. »
\end{exoresolu}

\begin{exercice}[Traduction (thème)]
Traduisez en allemand en choisissant le bon passif :
\begin{enumerate}
\item Les médicaments sont falsifiés. (constat d'un expert)
\item Les médicaments sont falsifiés dans un atelier clandestin. (activité régulière)
\item Le dossier était déjà clos quand le journaliste est arrivé.
\item Le marché public a été attribué hier. (l'action)
\item Depuis le scandale, tous les comptes sont bloqués.
\end{enumerate}
\end{exercice}

\begin{corrige}[Exercice : traduction]
\begin{enumerate}
\item \all{Die Medikamente sind gefälscht.} (état $\rightarrow$ \all{sein})
\item \all{Die Medikamente werden in einer illegalen Werkstatt gefälscht.} (action répétée $\rightarrow$ \all{werden})
\item \all{Die Akte war schon geschlossen, als der Journalist ankam.} (état au Präteritum ; \all{als} pour un événement unique passé)
\item \all{Der öffentliche Auftrag wurde gestern vergeben.} (action au Präteritum $\rightarrow$ \all{wurde})
\item \all{Seit dem Skandal sind alle Konten gesperrt.} (état actuel $\rightarrow$ \all{sind})
\end{enumerate}
\end{corrige}

\begin{aretenir}[L'essentiel du Zustandspassiv]
\begin{itemize}
\item Formation : \textbf{sein + Partizip II} ; temps composés avec \all{gewesen} (jamais \all{worden}).
\item Sens : l'\textbf{état résultant} d'une action achevée ; le Vorgangspassiv (\all{werden + Partizip II}) décrit l'\textbf{action}.
\item Test pratique : \all{gerade} possible $\rightarrow$ \all{werden} ; \all{schon/bereits} possible $\rightarrow$ \all{sein}.
\item Le français « est + participe » est ambigu : l'allemand impose un choix. C'est la langue du constat administratif et journalistique : \all{Der Auftrag ist vergeben, die Akte ist geschlossen, die Ware ist beschlagnahmt.}
\end{itemize}
\end{aretenir}

\begin{savaistu}[Pourquoi l'administration allemande aime tant cette forme]
Dans les textes officiels allemands — lois, rapports d'audit, décisions de justice —, le Zustandspassiv permet de constater un fait sans nommer d'auteur : \all{Der Fehler ist korrigiert.} « L'erreur est corrigée. » Cette tournure impersonnelle est si typique qu'on parle parfois de \all{Behördendeutsch}, « l'allemand des bureaux ». Savoir la décoder, c'est pouvoir lire un communiqué du ministère des Marchés publics ou un article de la \all{Süddeutsche Zeitung} sur un scandale de corruption sans dictionnaire de style.
\end{savaistu}

\coursec{Grammaire 2 : le discours indirect au subjonctif I}

Dans la section précédente, nous avons vu comment le passif d'état (\all{Zustandspassiv}) permet de décrire le résultat d'une action : \all{Die Ware ist gefälscht.} Nous restons dans l'univers de la presse et du reportage, mais nous changeons de perspective : comment un journaliste allemand rapporte-t-il les \emph{paroles} d'un ministre, d'un expert ou d'un accusé sans les prendre à son compte ? C'est toute la fonction du \cle{subjonctif I} (\all{Konjunktiv I}), temps-roi du discours indirect (\all{indirekte Rede}).

\courssub{Observation : deux phrases du texte}

Reprenons deux phrases tirées de notre texte support sur les marchés publics :

\begin{exemplebox}[Observation dans le texte]
\all{Der Minister erklärte, das Verfahren sei korrekt gewesen.} « Le ministre a déclaré que la procédure avait été correcte. »

\all{Ein Experte sagte, der Schmuggel habe zugenommen.} « Un expert a dit que la contrebande avait augmenté. »
\end{exemplebox}

Observons attentivement :

\begin{itemize}
\item Les verbes \all{sei} et \all{habe} ne sont pas des formes de l'indicatif : à l'indicatif, on aurait \all{ist} et \all{hat}.
\item Le journaliste ne dit pas « la procédure était correcte » (ce serait un fait présenté comme avéré) ; il dit « le ministre a déclaré que... » : il \emph{rapporte} une parole sans la garantir.
\item En français, la presse utilise souvent le conditionnel dans ce cas : « la procédure \emph{aurait été} correcte », « la contrebande \emph{aurait} augmenté ». L'allemand, lui, possède un outil grammatical spécialisé : le \cle{subjonctif I}.
\end{itemize}

\begin{definition}[Le discours indirect (indirekte Rede)]
Le \cle{discours indirect} rapporte les paroles ou les pensées de quelqu'un sans les citer mot à mot. En allemand écrit, surtout dans la presse, il se construit au \cle{subjonctif I} (\all{Konjunktiv I}), qui signale au lecteur : « ceci est une parole rapportée, je n'en garantis pas la vérité. »
\end{definition}

\courssub{Formation du subjonctif I présent}

\begin{propriete}[Formation du Konjunktiv I présent]
Le subjonctif I présent se forme sur le \textbf{radical de l'infinitif} + les désinences :

\begin{center}
\begin{tabular}{|l|l|}
\hline
\rowcolor{popblueL} Personne & Désinence \\
\hline
ich & -e \\
\hline
du & -est \\
\hline
er/sie/es & -e \\
\hline
wir & -en \\
\hline
ihr & -et \\
\hline
sie/Sie & -en \\
\hline
\end{tabular}
\end{center}

Exemple avec \all{sagen} : ich \all{sage}, du \all{sagest}, er \all{sage}, wir \all{sagen}, ihr \all{saget}, sie \all{sagen}.

\textbf{Seule exception} : le verbe \all{sein}, qui perd le -e à la première et à la troisième personne du singulier : ich \all{sei}, du \all{seiest}, er \all{sei}, wir \all{seien}, ihr \all{seiet}, sie \all{seien}.
\end{propriete}

\begin{attention}[Pas de voyelle modifiée au subjonctif I]
Contrairement au subjonctif II, le subjonctif I ne modifie \textbf{jamais} la voyelle du radical : \all{kommen} donne er \all{komme} (et non *\all{kämme} ; la forme \all{er käme} appartient au subjonctif II). De même : er \all{habe}, er \all{gebe}, er \all{nehme}.
\end{attention}

\begin{popfigure}
\begin{center}
\begin{tabular}{|l|l|l|l|l|l|}
\hline
\rowcolor{popblueL} Personne & sein & haben & werden & sagen & kommen \\
\hline
ich & sei & habe & werde & sage & komme \\
\hline
du & seiest & habest & werdest & sagest & kommest \\
\hline
er/sie/es & sei & habe & werde & sage & komme \\
\hline
wir & seien & haben & werden & sagen & kommen \\
\hline
ihr & seiet & habet & werdet & saget & kommet \\
\hline
sie/Sie & seien & haben & werden & sagen & kommen \\
\hline
\end{tabular}
\end{center}
\legende{Conjugaison complète au subjonctif I présent : sein, haben, werden, un verbe faible (sagen) et un verbe fort (kommen).}
\end{popfigure}

\courssub{Les trois temps du discours indirect}

Le discours indirect ne connaît que \textbf{trois temps} : présent, passé, futur. La règle de correspondance (concordance des temps) est simple :

\begin{propriete}[Concordance des temps au discours indirect]
\begin{itemize}
\item Discours direct au \textbf{présent} $\rightarrow$ KI \textbf{présent} : \all{„Ich komme.“} $\rightarrow$ Er sagte, er \all{komme}.
\item Discours direct à \textbf{n'importe quel temps passé} (Perfekt, Präteritum ou Plusquamperfekt) $\rightarrow$ un \textbf{seul} KI \textbf{passé} : \all{sei} / \all{habe} + participe II. \all{„Ich bin gekommen.“ / „Ich kam.“ / „Ich war gekommen.“} $\rightarrow$ Er sagte, er \all{sei gekommen}.
\item Discours direct au \textbf{futur} $\rightarrow$ KI \textbf{futur} : \all{werde} + infinitif. \all{„Ich werde kommen.“} $\rightarrow$ Er sagte, er \all{werde kommen}.
\end{itemize}
Le choix entre \all{habe} et \all{sei} au passé suit les mêmes règles qu'au Perfekt : verbes de mouvement ou de changement d'état avec \all{sein} (\all{er sei gekommen}), les autres avec \all{haben} (\all{er habe gesagt}).
\end{propriete}

\begin{popfigure}
\begin{center}
\begin{tikzpicture}[font=\small]
\node[draw=popblue, fill=popblueL, rounded corners, align=center, minimum width=3.2cm] (p1) at (0,3) {Présent\\ \all{„Ich komme.“}};
\node[draw=popgreen, fill=popgreenL, rounded corners, align=center, minimum width=3.2cm] (p2) at (7,3) {KI présent\\ \all{er komme}};
\draw[-{Stealth}, very thick, popdark] (p1) -- (p2);
\node[draw=poporange, fill=poporangeL, rounded corners, align=center, minimum width=3.2cm] (v1) at (0,1) {Perfekt / Präteritum /\\ Plusquamperfekt\\ \all{„Ich kam / bin gekommen.“}};
\node[draw=popgreen, fill=popgreenL, rounded corners, align=center, minimum width=3.2cm] (v2) at (7,1) {KI passé (unique)\\ \all{er sei gekommen}\\ \all{er habe gesagt}};
\draw[-{Stealth}, very thick, popdark] (v1) -- (v2);
\node[draw=poppurple, fill=poppurpleL, rounded corners, align=center, minimum width=3.2cm] (f1) at (0,-1.3) {Futur\\ \all{„Ich werde kommen.“}};
\node[draw=popgreen, fill=popgreenL, rounded corners, align=center, minimum width=3.2cm] (f2) at (7,-1.3) {KI futur\\ \all{er werde kommen}};
\draw[-{Stealth}, very thick, popdark] (f1) -- (f2);
\end{tikzpicture}
\end{center}
\legende{Frise des correspondances : trois temps passés de l'indicatif convergent vers un seul passé du subjonctif I.}
\end{popfigure}

\begin{exemplebox}[Exemples traduits : les trois temps]
\textbf{Présent :} Direkt : \all{„Ich bin unschuldig.“} $\rightarrow$ Indirekt : \all{Er sagte, er sei unschuldig.} « Il a dit qu'il était innocent. »

\textbf{Passé :} Direkt : \all{„Wir haben die Ware kontrolliert.“} $\rightarrow$ Indirekt : \all{Sie erklärten, sie hätten die Ware kontrolliert.} « Ils ont déclaré avoir contrôlé la marchandise. » (subjonctif II de substitution, voir plus bas.)

\textbf{Futur :} Direkt : \all{„Die Regierung wird den Schmuggel bekämpfen.“} $\rightarrow$ Indirekt : \all{Der Minister sagte, die Regierung werde den Schmuggel bekämpfen.} « Le ministre a dit que le gouvernement combattrait la contrebande. »
\end{exemplebox}

\courssub{La règle de substitution : quand le KI ressemble à l'indicatif}

Problème : à certaines personnes, le subjonctif I est \emph{identique} à l'indicatif présent : \all{wir sagen}, \all{sie kommen}, \all{ihr kommet}. Le lecteur ne verrait alors plus la marque du discours indirect.

\begin{propriete}[Règle de substitution (Ersatzregel)]
Quand la forme du subjonctif I est identique à celle de l'indicatif, on la remplace par le \textbf{subjonctif II} pour garder la marque de l'indirect :
\begin{itemize}
\item \all{sie kommen} (KI = indicatif) $\rightarrow$ \all{sie kämen} : \all{Sie sagten, sie kämen morgen.} « Elles ont dit qu'elles viendraient demain. »
\item \all{sie haben kontrolliert} $\rightarrow$ \all{sie hätten kontrolliert} : \all{Sie erklärten, sie hätten alles kontrolliert.} « Elles ont déclaré avoir tout contrôlé. »
\end{itemize}
Attention : ce subjonctif II n'a ici \textbf{aucune valeur d'irréel} ; il sert uniquement de signal du discours indirect.
\end{propriete}

\begin{attention}[Ne pas confondre KI (parole rapportée) et KII (irréel)]
C'est le piège majeur pour les francophones :
\begin{itemize}
\item \all{Er komme.} = « il vient, \emph{paraît-il} » (parole rapportée, subjonctif I).
\item \all{Er käme.} = « il viendrait » (hypothèse, irréel, subjonctif II).
\end{itemize}
Dans un article de presse, \all{Der Minister sei korrupt} signifie « le ministre serait corrompu \emph{selon les déclarations rapportées} », et non « si le ministre était corrompu ». En version, traduisez le KI par le conditionnel journalistique français ou par une tournure avec « selon » : « le ministre aurait déclaré... », « selon l'expert, la contrebande aurait augmenté ».
\end{attention}

\courssub{Discours indirect libre : avec ou sans dass}

\begin{propriete}[Deux constructions possibles]
\begin{itemize}
\item \textbf{Avec \all{dass}} : le verbe conjugué va à la fin : \all{Er sagte, dass er nichts gewusst habe.}
\item \textbf{Sans \all{dass}} (discours indirect libre) : le verbe conjugué prend la \textbf{deuxième place}, comme dans une principale : \all{Er sagte, er habe nichts gewusst.} « Il a dit qu'il n'avait rien su. »
\end{itemize}
Les deux formes sont correctes ; la seconde est très fréquente dans la presse.
\end{propriete}

\begin{exemplebox}[Verbes introducteurs fréquents]
\all{sagen} (dire), \all{behaupten} (prétendre, affirmer), \all{erklären} (déclarer), \all{berichten} (rapporter), \all{meinen} (penser, estimer), \all{zugeben} (admettre, avouer), \all{dementieren} (démentir), \all{betonen} (souligner).

\all{Der Angeklagte behauptete, er habe die Dokumente nie gesehen.} « L'accusé a prétendu n'avoir jamais vu les documents. »

\all{Die Regierung dementierte, Gelder unterschlagen zu haben.} « Le gouvernement a démenti avoir détourné des fonds. »
\end{exemplebox}

\courssub{Comparaison français / allemand}

\begin{center}
\begin{tabularx}{\linewidth}{|X|X|X|}
\hline
\rowcolor{popblueL} Français & Deutsch & Remarque \\
\hline
Le ministre a déclaré que la procédure était correcte. & Der Minister erklärte, das Verfahren sei korrekt gewesen. & L'allemand marque l'indirect par le KI ; le français garde l'indicatif. \\
\hline
La contrebande aurait augmenté (conditionnel journalistique). & Der Schmuggel habe zugenommen, sagte ein Experte. & Le conditionnel français = le KI allemand. \\
\hline
Il prétend être innocent. & Er behauptet, unschuldig zu sein. / Er behauptet, er sei unschuldig. & Deux tournures possibles en allemand. \\
\hline
Selon les témoins, les marchandises auraient été falsifiées. & Den Zeugen zufolge seien die Waren gefälscht worden. & Passif + KI : très fréquent dans la presse. \\
\hline
\end{tabularx}
\end{center}

\begin{methode}[Transformer une phrase du discours direct au discours indirect]
\begin{enumerate}
\item \textbf{Repérer le temps} du discours direct (présent, passé, futur).
\item \textbf{Choisir le temps KI correspondant} : présent $\rightarrow$ KI présent ; tout temps passé $\rightarrow$ KI passé (\all{habe/sei} + participe II) ; futur $\rightarrow$ \all{werde} + infinitif.
\item \textbf{Adapter les pronoms} : \all{ich} $\rightarrow$ \all{er/sie}, \all{mein} $\rightarrow$ \all{sein/ihr}, selon le contexte.
\item \textbf{Supprimer les guillemets} et introduire la parole rapportée avec \all{dass} + verbe final, ou sans \all{dass} avec le verbe en deuxième position.
\item \textbf{Vérifier la substitution} : si la forme KI est identique à l'indicatif (\all{sie kommen}), passer au KII (\all{sie kämen}).
\end{enumerate}
\end{methode}

\courssub{Application sur un texte de presse}

\begin{textebox}[Pressemitteilung : Ermittlungen im Hafen von Douala]
\all{Bei einer Pressekonferenz in Douala erklärte ein Regierungssprecher, die Zollbehörden hätten im Hafen mehrere Container mit gefälschten Medikamenten beschlagnahmt. Die Waren seien als Import aus Europa deklariert worden, stammten aber in Wirklichkeit aus einer illegalen Fabrik. Ein Zollbeamter, der anonym bleiben wolle, sagte, der Schmuggel habe in den letzten Monaten deutlich zugenommen. Die Täter seien gut organisiert und verfügten über gefälschte Dokumente. Ein Vertreter der Staatsanwaltschaft berichtete, mehrere Verdächtige seien bereits verhaftet worden. Sie hätten zugegeben, Bestechungsgelder an Hafenmitarbeiter gezahlt zu haben. Der Sprecher betonte, die Regierung werde den Kampf gegen die Fälschung verstärken. Neue Kontrollsysteme würden bald installiert. Experten meinen jedoch, das Problem sei tiefer: Solange Beamte schlecht bezahlt würden, bleibe die Korruption eine Versuchung. Die Bevölkerung, so der Sprecher, müsse gefälschte Produkte melden, denn sie gefährdeten die Gesundheit aller.}

\textbf{Glossaire :} der Zoll, die Zölle (la douane) ; beschlagnahmen (saisir) ; deklarieren (déclarer en douane) ; der Zollbeamte, -n (le douanier) ; die Staatsanwaltschaft, -en (le parquet) ; der Verdächtige, -n (le suspect) ; das Bestechungsgeld, -er (le pot-de-vin) ; betonen (souligner) ; die Versuchung, -en (la tentation) ; melden (signaler) ; gefährden (mettre en danger).
\end{textebox}

\textbf{Questions de compréhension :}
\begin{enumerate}
\item \all{Was wurde im Hafen von Douala beschlagnahmt?}
\item \all{Woher stammten die Waren wirklich?}
\item \all{Was haben die Verdächtigen zugegeben?}
\item Relevez trois formes de subjonctif I dans le texte et identifiez leur temps (présent, passé, futur).
\end{enumerate}

\begin{exoresolu}[Du discours direct au discours indirect]
Mettez au discours indirect (introduit par \all{Der Minister erklärte,...}) : \all{„Wir haben alle Firmen kontrolliert. Das Verfahren war korrekt. Wir werden die Korruption bekämpfen.“}
\tcblower
Étape 1 : les trois phrases sont respectivement au Perfekt, au Präteritum et au Futur I. Étape 2 : Perfekt et Präteritum $\rightarrow$ KI passé ; Futur $\rightarrow$ KI futur. Étape 3 : \all{wir} $\rightarrow$ \all{sie} (le ministre et son équipe). Or \all{sie haben} et \all{sie werden} seraient identiques à l'indicatif : la règle de substitution impose donc le KII \all{hätten} et \all{würden}.

Solution : \all{Der Minister erklärte, sie hätten alle Firmen kontrolliert. Das Verfahren sei korrekt gewesen. Sie würden die Korruption bekämpfen.}

Traduction : « Le ministre a déclaré qu'ils avaient contrôlé toutes les entreprises, que la procédure avait été correcte et qu'ils combattraient la corruption. »
\end{exoresolu}

\begin{exercice}[Entraînement]
\begin{enumerate}
\item Version : traduisez en français : \all{Ein Experte sagte, der Schmuggel habe zugenommen, weil die Kontrollen zu schwach seien.}
\item Mettez au discours indirect : \all{„Ich habe keine Bestechungsgelder erhalten“, sagte der Beamte.}
\item Thème : traduisez en allemand : « Les témoins ont déclaré que les marchandises avaient été falsifiées. »
\item Thème : traduisez : « Selon le journal, le ministre aurait signé le contrat. »
\end{enumerate}
\end{exercice}

\begin{corrige}[Exercice : Entraînement]
\begin{enumerate}
\item « Un expert a dit que la contrebande avait augmenté parce que les contrôles étaient trop faibles » (ou : « ...aurait augmenté, les contrôles étant, selon lui, trop faibles »). \all{habe zugenommen} = KI passé ; \all{seien} = KI présent.
\item \all{Der Beamte sagte, er habe keine Bestechungsgelder erhalten.} (Perfekt $\rightarrow$ KI passé avec \all{habe} + participe II.)
\item \all{Die Zeugen erklärten, die Waren seien gefälscht worden.} (passif au KI passé : \all{seien} + participe II + \all{worden}.)
\item \all{Laut dem Bericht habe der Minister den Vertrag unterschrieben.} ou \all{Der Zeitung zufolge habe der Minister den Vertrag unterschrieben.}
\end{enumerate}
\end{corrige}

\begin{savaistu}[Le subjonctif I, bouclier du journaliste]
En Allemagne, le subjonctif I est la marque du journalisme sérieux. En écrivant \all{Der Manager soll Millionen unterschlagen haben} ou \all{Er habe Gelder unterschlagen}, le journaliste rapporte une accusation \emph{sans l'affirmer} : juridiquement, il ne fait que citer une source. Cette nuance protège la presse contre les procès en diffamation, tout en informant le public. C'est pourquoi les articles des grands quotidiens germanophones (\all{Süddeutsche Zeitung}, \all{Neue Zürcher Zeitung}) sont remplis de \all{sei}, \all{habe}, \all{werde} : chaque forme est un petit panneau « parole rapportée ».
\end{savaistu}

\begin{aretenir}[L'essentiel du discours indirect]
\begin{itemize}
\item Le discours indirect allemand s'exprime au \textbf{subjonctif I} : radical + \textbf{-e, -est, -e, -en, -et, -en} ; seul \all{sein} est irrégulier (\all{sei, seiest, sei, seien, seiet, seien}).
\item Trois temps seulement : KI présent (\all{er komme}), KI passé unique pour tous les passés (\all{er sei gekommen / habe gesagt}), KI futur (\all{er werde kommen}).
\item Si le KI est identique à l'indicatif (\all{sie kommen}), on le remplace par le KII (\all{sie kämen}) — sans valeur d'irréel.
\item Avec \all{dass}, le verbe va en fin de phrase ; sans \all{dass}, il reste en deuxième position.
\item En version, rendez le KI par le conditionnel journalistique français ou par « selon... ».
\end{itemize}
\end{aretenir}

\coursec{Conjugaison : subjonctif I et II, tableaux de synthèse}

Dans la section précédente, nous avons vu que le discours indirect journalistique s'exprime au \cle{subjonctif I} : \all{Der Zollbeamte sagte, er sei bestochen worden.} Mais notre texte sur les marchés publics contient aussi des hypothèses : \all{Wenn die Kontrollen strenger wären, gäbe es weniger Schmuggel.} Ces deux formes, \all{sei} et \all{wären}, appartiennent à deux modes différents qu'il faut maintenant conjuguer sans erreur. Cette section est la boîte à outils complète : tableaux, règles de formation, emplois contrastés et entraînement.

\courssub{Observer avant de conjuguer}

\begin{textebox}[Extraits du dossier de presse]
Der Staatsanwalt erklärte, die Firma habe bei der Ausschreibung betrogen. Sie habe gefälschte Dokumente vorgelegt und Beamte bestochen. Ein Insider berichtete, der Minister wisse seit Monaten von dem Betrug, aber er schweige. Hätte die Kontrollkommission früher reagiert, wäre der Schmuggel gestoppt worden. Wenn es keine Korruption gäbe, wären die Preise niedriger und die Straßen besser. Ein Händler am Markt von Douala meinte: „Würde der Zoll ehrlich arbeiten, könnten wir alle legal importieren.“
\end{textebox}

\textbf{Observation.} Soulignons les formes : \all{habe betrogen}, \all{habe vorgelegt}, \all{wisse}, \all{schweige} (subjonctif I : on \emph{rapporte} des paroles) ; \all{hätte reagiert}, \all{wäre gestoppt worden}, \all{gäbe}, \all{wären}, \all{würde arbeiten}, \all{könnten} (subjonctif II : on \emph{imagine} l'irréel). Même phrase, deux mondes : le KI rapporte, le KII rêve, suppose ou conseille.

\courssub{Formation : les deux recettes de base}

\begin{propriete}[Formation du subjonctif I (Konjunktiv I)]
On prend le \cle{radical de l'infinitif} et on ajoute les désinences : -e, -est, -e, -en, -et, -en. Seul le verbe \all{sein} est irrégulier : \all{ich sei, du seist, er sei, wir seien, ihr seiet, sie seien}. En pratique, on utilise surtout la 3e personne du singulier (\all{er habe, er sei, er komme}) et du pluriel (\all{sie hätten, sie seien}).
\end{propriete}

\begin{propriete}[Formation du subjonctif II (Konjunktiv II)]
On prend le \cle{radical du Präteritum} ; pour les verbes forts, on ajoute un \cle{Umlaut} sur a, o, u (a $\rightarrow$ ä, o $\rightarrow$ ö, u $\rightarrow$ ü) puis les désinences -e, -est, -e, -en, -et, -en : \all{kam $\rightarrow$ käme}, \all{gab $\rightarrow$ gäbe}, \all{wusste $\rightarrow$ wüsste}. Les verbes faibles ont un KII \emph{identique au Präteritum} : \all{sagte, machte}.
\end{propriete}

\courssub{Tableau de synthèse : les auxiliaires et les modaux}

\begin{center}
\begin{tabularx}{\linewidth}{|X|X|X|X|}
\hline
\rowcolor{popblueL} \textbf{Infinitif} & \textbf{KI présent (er)} & \textbf{KII présent (er)} & \textbf{Français (KII)} \\
\hline
sein & sei & wäre & il serait \\
\hline
haben & habe & hätte & il aurait \\
\hline
werden & werde & würde & il deviendrait \\
\hline
können & könne & könnte & il pourrait \\
\hline
müssen & müsse & müsste & il devrait (obligation) \\
\hline
sollen & solle & sollte & il devrait (conseil) \\
\hline
dürfen & dürfe & dürfte & il serait permis \\
\hline
wollen & wolle & wollte & il voudrait \\
\hline
\end{tabularx}
\end{center}

\begin{exemplebox}[Les auxiliaires dans notre thème]
\all{Der Beamte sagte, er habe nichts unterschrieben.} « Le fonctionnaire a dit qu'il n'avait rien signé. » (KI)\par
\all{An seiner Stelle hätte ich den Vertrag nicht unterschrieben.} « À sa place, je n'aurais pas signé le contrat. » (KII)\par
\all{Ohne Bestechung wäre alles schneller.} « Sans corruption, tout serait plus rapide. »
\end{exemplebox}

\courssub{Verbes faibles, verbes forts, verbes irréguliers}

\begin{center}
\begin{tabularx}{\linewidth}{|X|X|X|X|}
\hline
\rowcolor{popgreenL} \textbf{Infinitif} & \textbf{KI (er)} & \textbf{KII (er)} & \textbf{Remarque} \\
\hline
sagen (faible) & sage & sagte & KII = Präteritum \\
\hline
machen (faible) & mache & machte & KII = Präteritum \\
\hline
kommen (fort) & komme & käme & Umlaut \\
\hline
geben (fort) & gebe & gäbe & Umlaut \\
\hline
nehmen (fort) & nehme & nähme & Umlaut \\
\hline
wissen (irrégulier) & wisse & wüsste & Umlaut \\
\hline
es gibt & es gebe & es gäbe & forme impersonnelle \\
\hline
\end{tabularx}
\end{center}

\begin{attention}[Le piège des verbes faibles]
Pour les verbes faibles, le KII est identique au Präteritum : \all{er sagte} peut vouloir dire « il dit » (passé réel) ou « il dirait » (irréel). Conséquences pratiques : à l'\emph{écrit journalistique}, on préfère le KI pour le discours indirect (\all{er sage}) afin d'éviter toute ambiguïté ; à l'\emph{oral}, on remplace le KII des verbes faibles par \all{würde} + infinitif : \all{Ich würde das nicht sagen.} « Je ne dirais pas cela. » Ne dites jamais \all{ich sagte} pour « je dirais » : c'est fautif à l'oral soutenu.
\end{attention}

\courssub{Les temps passés : habe/sei contre hätte/wäre}

\begin{propriete}[Subjonctif passé : la règle d'or]
\textbf{Subjonctif I passé} (discours indirect au passé) : \all{habe} ou \all{sei} + Partizip II. On choisit \all{sei} avec les verbes de mouvement et de changement d'état (comme l'auxiliaire \all{sein} du Perfekt).\par
\all{Er sagte, er habe das Geld genommen.} « Il a dit qu'il avait pris l'argent. »\par
\all{Sie behauptete, die Ware sei gestern angekommen.} « Elle a affirmé que la marchandise était arrivée hier. »\par
\textbf{Subjonctif II passé} (irréel du passé) : \all{hätte} ou \all{wäre} + Partizip II.\par
\all{Er hätte das Geld genommen.} « Il aurait pris l'argent. »\par
\all{Wäre die Kontrolle gekommen, hätten sie die Fälschung entdeckt.} « Si le contrôle était venu, ils auraient découvert la falsification. »
\end{propriete}

\begin{aretenir}[Le contraste en une paire de phrases]
\all{Der Händler sagte, er habe die gefälschten Medikamente nicht verkauft.} = le commerçant \emph{affirme} (on rapporte, sans juger).\par
\all{Der Händler hätte die gefälschten Medikamente nicht verkaufen sollen.} = il \emph{n'aurait pas dû} les vendre (reproche, irréel du passé).\par
Le français traduit le KI par l'indicatif (il a dit que\ldots) et le KII par le conditionnel (il aurait\ldots).
\end{aretenir}

\courssub{Tableau TikZ à double entrée : les quatre cases du subjonctif}

\begin{popfigure}
\begin{center}
\begin{tikzpicture}
\matrix (m) [matrix of nodes, nodes={draw=popline, minimum height=0.85cm, align=center, anchor=center, inner sep=2pt}, column sep=-\pgflinewidth, row sep=-\pgflinewidth,
    column 1/.style={nodes={fill=popgreenL, minimum width=2.6cm, font=\bfseries, text width=2.4cm}},
    column 2/.style={nodes={fill=popblueL, minimum width=3.4cm, font=\bfseries, text width=3.2cm}},
    column 3/.style={nodes={fill=popblueL, minimum width=3.4cm, text width=3.2cm}},
    column 4/.style={nodes={fill=poporangeL, minimum width=3.4cm, font=\bfseries, text width=3.2cm}},
    column 5/.style={nodes={fill=poporangeL, minimum width=3.4cm, text width=3.2cm}},
    row 1/.style={nodes={font=\bfseries, minimum height=0.85cm}},
    row 1 column 1/.style={nodes={fill=popblueL, minimum width=2.6cm}},
    row 1 column 4/.style={nodes={fill=poporangeL, font=\bfseries}},
    row 1 column 5/.style={nodes={fill=poporangeL, font=\bfseries}},
]
{
& KI présent & KI passé & KII présent & KII passé \\
sein & er sei & er sei gewesen & er wäre & er wäre gewesen \\
haben & er habe & er habe gehabt & er hätte & er hätte gehabt \\
werden & er werde & er sei geworden & er würde & er wäre geworden \\
können & er könne & er habe gekonnt & er könnte & er hätte gekonnt \\
müssen & er müsse & er habe gemusst & er müsste & er hätte gemusst \\
};
\end{tikzpicture}
\end{center}
\legende{Tableau à double entrée : subjonctif I et II, présent et passé, pour les auxiliaires et les verbes modaux.}
\end{popfigure}

\courssub{Emplois contrastés : rapporter ou imaginer}

\begin{definition}[Les deux missions des subjonctifs]
Le \cle{subjonctif I} est le mode du \emph{discours rapporté} : le journaliste transmet sans prendre position. Le \cle{subjonctif II} est le mode de l'\emph{irréel} : souhait irréalisable, hypothèse, conseil, politesse atténuée.
\end{definition}

\begin{center}
\begin{tabularx}{\linewidth}{|X|X|}
\hline
\rowcolor{popblueL} \textbf{Subjonctif I : on rapporte} & \textbf{Subjonctif II : on imagine} \\
\hline
Der Minister erklärte, er kenne die Firma nicht. (Le ministre a déclaré qu'il ne connaissait pas l'entreprise.) & Wenn der Minister die Firma kenne, würde er es zugeben. (Si le ministre connaissait l'entreprise, il l'admettrait.) \\
\hline
Sie sagte, die Waren seien beschlagnahmt worden. (Elle a dit que les marchandises avaient été saisies.) & Die Waren wären nicht beschlagnahmt worden, wenn die Papiere echt gewesen wären. (Les marchandises n'auraient pas été saisies si les papiers avaient été authentiques.) \\
\hline
Er behauptete, er habe nie geschmuggelt. (Il prétendait n'avoir jamais fait de contrebande.) & Hätte er nie geschmuggelt, wäre er jetzt frei. (S'il n'avait jamais fait de contrebande, il serait libre maintenant.) \\
\hline
\end{tabularx}
\end{center}

\begin{exemplebox}[Les quatre emplois du subjonctif II]
\textbf{Souhait irréel} : \all{Wäre der Markt doch sauber!} « Si seulement le marché était propre ! »\par
\textbf{Hypothèse} : \all{Wenn es keine Korruption gäbe, wären die Preise niedriger.} « S'il n'y avait pas de corruption, les prix seraient plus bas. »\par
\textbf{Conseil} : \all{Du solltest die Rechnung prüfen.} « Tu devrais vérifier la facture. »\par
\textbf{Politesse} : \all{Ich hätte eine Frage. Könnten Sie mir helfen?} « J'aurais une question. Pourriez-vous m'aider ? »
\end{exemplebox}

\begin{methode}[Choisir entre KI et KII : la check-list en 3 questions]
\begin{enumerate}
\item \textbf{Est-ce un discours rapporté (journalisme, témoignage) ?} $\rightarrow$ \textbf{KI} (présent ou passé selon le temps de l'énoncé original).
\item \textbf{Est-ce une hypothèse, un souhait, un conseil, une politesse ?} $\rightarrow$ \textbf{KII} (présent pour le présent/futur, passé pour le passé).
\item \textbf{Verbe faible au KII ?} $\rightarrow$ À l'écrit : utilisez le KI si c'est du discours rapporté ; à l'oral : utilisez \all{würde} + infinitif.
\end{enumerate}
\end{methode}

\courssub{Le duo KI + KII dans le texte journalistique}

Un bon article d'investigation alterne les deux modes : le KI pour rapporter les déclarations des acteurs, le KII pour formuler des hypothèses ou des critiques voilées. Modèle de phrase combinée : \all{Der Zollbeamte sagte, er sei bestochen worden; hätte er das Geld angenommen, hätte er es zugegeben.} « Le douanier a dit qu'il avait été soudoyé ; s'il avait accepté l'argent, il l'aurait admis. » Remarquez le passif au KI : \all{er sei bestochen worden} (il a été corrompu) — le participe \all{worden} sans ge- après \all{sei/wäre}.

\begin{attention}[Erreurs fréquentes des francophones]
\begin{itemize}
\item Ne traduisez pas le conditionnel français par un KI : « il serait » = \all{er wäre}, jamais \all{er sei}.
\item N'oubliez pas l'Umlaut au KII des verbes forts : \all{käme}, \all{gäbe}, \all{wüsste} — jamais \all{kame} ou \all{gabe}.
\item Au KII passé, l'auxiliaire suit les mêmes règles qu'au Perfekt : \all{er wäre gekommen} (mouvement), \all{er hätte gemacht} (action).
\item Avec deux verbes modaux au passé, l'infinitif double : \all{er hätte kommen müssen} « il aurait dû venir ».
\item Le français « j'aurais aimé » se dit \all{ich hätte gerne\ldots} (ex. \all{ich hätte gerne mitgemacht}) ou \all{das hätte ich gern gehabt} (pour un objet), mais jamais \all{ich hätte gern gehabt} seul sans complément.
\end{itemize}
\end{attention}

\courssub{Entraînement ciblé}

\begin{exoresolu}[Conjugaison à trous : KI ou KII ?]
Complétez avec la forme correcte du verbe entre parenthèses.\par
1. Der Anwalt sagte, sein Mandant \ldots\ unschuldig. (sein)\par
2. Wenn die Kontrollen strenger \ldots, gäbe es weniger Schmuggel. (sein)\par
3. Sie behauptete, sie \ldots\ das Geld nie gesehen. (haben)\par
4. \ldots\ ich an Ihrer Stelle, würde ich den Vertrag ablehnen. (sein)\par
5. Der Händler meinte, man \ldots\ die gefälschten Produkte sofort vernichten. (müssen, KI)\par
6. Ohne Bestechung \ldots\ alles schneller gehen. (können, KII)
\tcblower
\textbf{Corrigé détaillé.}\par
1. \all{sein Mandant sei unschuldig} — discours indirect : KI présent de \all{sein}.\par
2. \all{strenger wären} — hypothèse irréelle introduite par \all{wenn} : KII présent de \all{sein}, pluriel.\par
3. \all{sie habe das Geld nie gesehen} — discours indirect au passé : KI passé avec \all{habe} + Partizip II (verbe d'action).\par
4. \all{Wäre ich an Ihrer Stelle} — hypothèse avec inversion : KII de \all{sein} à la 1re personne.\par
5. \all{man müsse} — parole rapportée : KI de \all{müssen}.\par
6. \all{könnte alles schneller gehen} — conséquence irréelle : KII de \all{können} + infinitif.
\end{exoresolu}

\begin{exercice}[Transformations]
A. Mettez au discours indirect (KI) : 1. „Ich habe die Dokumente nicht gefälscht.“ (Der Kaufmann sagte, \ldots) 2. „Die Ware wurde am Hafen beschlagnahmt.“ (Die Zeitung berichtete, \ldots)\par
B. Mettez à l'irréel du passé (KII passé) : 3. Er hat das Geld nicht angenommen. 4. Wir haben die Rechnung nicht geprüft.\par
C. Choisissez le bon mode : 5. Der Minister erklärte, er (wäre / sei) bereit zu kooperieren. 6. Wenn ich Sie wäre, (würde / werde) ich den Anwalt rufen.
\end{exercice}

\begin{corrige}[Exercice : Transformations]
A. 1. \all{Der Kaufmann sagte, er habe die Dokumente nicht gefälscht.} 2. \all{Die Zeitung berichtete, die Ware sei am Hafen beschlagnahmt worden.}\par
B. 3. \all{Er hätte das Geld nicht angenommen.} 4. \all{Wir hätten die Rechnung nicht geprüft.}\par
C. 5. \all{sei} (discours rapporté : KI). 6. \all{würde} (conseil irréel : KII avec \all{würde} + infinitif).
\end{corrige}

\begin{aretenir}[L'essentiel de la section]
KI = radical de l'infinitif + -e (il rapporte : \all{er habe, er sei, er komme}) ; KII = radical du Präteritum + Umlaut (il imagine : \all{er hätte, er wäre, er käme}). Passé : KI avec \all{habe/sei} + Partizip II ; KII avec \all{hätte/wäre} + Partizip II. Verbes faibles : KII = Präteritum, donc à l'écrit on préfère le KI, à l'oral \all{würde} + infinitif. En français : KI $\rightarrow$ indicatif, KII $\rightarrow$ conditionnel.
\end{aretenir}

\coursec{Méthode du commentaire de texte et commentaire modèle}

Dans les sections précédentes, nous avons maîtrisé la conjugaison du subjonctif I et du subjonctif II. Ces formes ne sont pas de simples exercices grammaticaux : elles constituent l'outil indispensable du \cle{commentaire de texte} (\all{Textkommentar}), l'exercice phare de l'expression écrite au baccalauréat. Commenter un texte, c'est rapporter fidèlement les idées d'un auteur \emph{sans les faire siennes} — et c'est exactement la fonction du subjonctif I. Nous allons maintenant apprendre, pas à pas, à construire un commentaire complet, structuré et correct, en l'appliquant à notre texte sur les marchés publics et la criminalité économique.

\courssub{Qu'est-ce qu'un commentaire de texte ?}

\begin{definition}[Le commentaire de texte — der Textkommentar]
Le \cle{commentaire de texte} est une production écrite en trois parties qui présente un texte, en résume et en analyse le contenu, puis se termine par une opinion personnelle justifiée. Il ne s'agit ni d'une traduction, ni d'un simple résumé, ni d'une dissertation libre : c'est un exercice codifié, avec des phrases-types et une grammaire imposée (le discours indirect au subjonctif I).
\end{definition}

\begin{propriete}[Les trois parties obligatoires]
\begin{enumerate}
\item \textbf{\all{die Einleitung}} (l'introduction, 2 à 3 lignes) : on présente le texte — type de texte, source, thème — et on formule la \cle{problématique} (\all{die Problemstellung}) sous forme d'une question.
\item \textbf{\all{der Hauptteil}} (le développement, 6 à 9 lignes) : on résume le texte de façon structurée (\all{Zusammenfassung}) puis on analyse la position de l'auteur et ses arguments (\all{Analyse}), \textbf{impérativement au subjonctif I}.
\item \textbf{\all{der Schluss}} (la conclusion, 2 à 3 lignes) : on fait le bilan et on donne son \cle{opinion personnelle} (\all{die eigene Meinung}), nuancée et justifiée.
\end{enumerate}
\end{propriete}

\begin{popfigure}
\begin{tikzpicture}[
node distance=0.55cm,
box/.style={draw=popblue, thick, rounded corners, fill=popblueL, align=center, text width=4.6cm, inner sep=6pt},
arr/.style={-{Stealth[length=3mm]}, thick, popdark}]
\node[box, align=center] (ein){\textbf{Einleitung} (2--3 Zeilen)\\ Textsorte, Quelle, Thema\\ + Problemstellung};
\node[box, below=of ein, align=center] (haupt){\textbf{Hauptteil} (6--9 Zeilen)\\ Zusammenfassung + Analyse\\ (Konjunktiv I !)};
\node[box, below=of haupt, align=center] (schluss){\textbf{Schluss} (2--3 Zeilen)\\ Bilanz + eigene Meinung\\ (begründet)};
\draw[arr] (ein) -- (haupt);
\draw[arr] (haupt) -- (schluss);
\end{tikzpicture}
\legende{Organigramme du commentaire de texte : trois parties, avec le nombre de lignes indicatif.}
\end{popfigure}

\courssub{Les phrases-types du commentaire}

Un bon commentaire repose sur des formules fixes qu'il faut connaître par cœur. Les voici, classées par partie, avec leur traduction.

\textbf{1. Pour l'introduction (\all{Einleitung})}\par

\begin{center}
\begin{tabularx}{\linewidth}{|X|X|}
\hline
\rowcolor{popblueL} \textbf{Deutsch} & \textbf{Français} \\
\hline
\all{Der vorliegende Text ist ein Artikel aus einer Zeitung.} & Le texte présenté est un article de journal. \\
\hline
\all{Er behandelt das Problem der Korruption bei öffentlichen Aufträgen.} & Il traite le problème de la corruption dans les marchés publics. \\
\hline
\all{Es geht um die Frage: Wie kann man den Schmuggel bekämpfen?} & Il s'agit de la question : comment lutter contre la contrebande ? \\
\hline
\all{Der Text stammt aus dem Jahr 2024.} & Le texte date de 2024. \\
\hline
\end{tabularx}
\end{center}

\textbf{2. Pour le résumé et l'analyse (\all{Hauptteil})}\par

\begin{center}
\begin{tabularx}{\linewidth}{|X|X|}
\hline
\rowcolor{popblueL} \textbf{Deutsch} & \textbf{Français} \\
\hline
\all{Der Autor behauptet, dass die Korruption der Wirtschaft schade.} & L'auteur affirme que la corruption nuit à l'économie. \\
\hline
\all{Er vertritt die Meinung, die Lage sei ernst.} & Il soutient l'opinion que la situation est grave. \\
\hline
\all{Als Beispiel nennt er gefälschte Medikamente.} & Comme exemple, il cite les médicaments falsifiés. \\
\hline
\all{Seiner Ansicht nach seien strengere Kontrollen nötig.} & Selon lui, des contrôles plus stricts seraient nécessaires. \\
\hline
\all{Zuerst beschreibt er..., dann erklärt er..., schließlich fordert er...} & D'abord il décrit..., puis il explique..., enfin il exige... \\
\hline
\end{tabularx}
\end{center}

\textbf{3. Pour l'opinion personnelle (\all{Schluss})}\par

\begin{center}
\begin{tabularx}{\linewidth}{|X|X|}
\hline
\rowcolor{popblueL} \textbf{Deutsch} & \textbf{Français} \\
\hline
\all{Meiner Meinung nach muss der Staat strengere Kontrollen einführen.} & À mon avis, l'État doit instaurer des contrôles plus stricts. \\
\hline
\all{Ich bin der Ansicht, dass Aufklärung wichtiger ist als Strafen.} & Je suis d'avis que la sensibilisation est plus importante que les punitions. \\
\hline
\all{Man sollte die Bürger über die Gefahren informieren.} & On devrait informer les citoyens des dangers. \\
\hline
\all{Zusammenfassend lässt sich sagen, dass...} & En résumé, on peut dire que... \\
\hline
\end{tabularx}
\end{center}

\begin{exemplebox}[Deux exemples traduits à retenir]
\all{Der Autor vertritt die Meinung, die Korruption schade der Wirtschaft.} „ L'auteur soutient que la corruption nuit à l'économie. “ — Remarque : \all{schade} est le subjonctif I de \all{schaden} (3ᵉ pers. sg.) ; il signale que c'est l'opinion de l'auteur, pas un fait établi.

\all{Meiner Meinung nach muss der Staat strengere Kontrollen einführen.} „ À mon avis, l'État doit instaurer des contrôles plus stricts. “ — Ici, indicatif : c'est \emph{mon} opinion, je l'assume.
\end{exemplebox}

\courssub{Le subjonctif I : la grammaire du commentaire}

\begin{propriete}[Règle d'or du discours rapporté dans le commentaire]
Dans le \all{Hauptteil}, toutes les idées de l'auteur se rapportent au \cle{subjonctif I} :
\begin{itemize}
\item avec \all{dass} : \all{Der Autor behauptet, dass die Korruption der Wirtschaft schade.}
\item sans \all{dass} (verbe en deuxième position) : \all{Der Autor sagt, die Lage sei ernst.}
\end{itemize}
Si le subjonctif I est identique à l'indicatif (ex. \all{sie schaden} / \all{sie geben}), on le remplace par le subjonctif II : \all{Der Autor meint, sie schadeten der Wirtschaft.} / \all{Der Autor sagt, sie gäben den Auftrag.}
En revanche, dans le \all{Schluss}, ton opinion personnelle s'exprime à l'\cle{indicatif} (ou au subjonctif II pour un souhait ou une recommandation : \all{Man sollte handeln}).
\end{propriete}

\begin{attention}[L'erreur classique du francophone]
Au commentaire, ne jamais écrire \all{Der Autor sagt, dass die Lage ernst ist.} à l'indicatif pour rapporter une opinion : en allemand, cela sonne comme si tu présentais le propos comme un fait établi. Il faut le subjonctif I : \all{Der Autor sagt, die Lage sei ernst.} En français, on utilise simplement l'indicatif ou le conditionnel (« l'auteur affirme que la situation est grave ») : la nuance grammaticale allemande n'a pas d'équivalent direct, d'où l'oubli fréquent. Vérifie chaque verbe de ton \all{Hauptteil} !
\end{attention}

\begin{aretenir}[Lien avec la Grammaire 1 : le passif d'état (Zustandspassiv)]
Le texte support contient plusieurs \cle{passifs d'état} (\all{Zustandspassiv}) qui décrivent une situation résultante : \all{sind bestochen} (sont corrompus), \all{sind schlecht ausgeführt} (sont mal exécutés), \all{ist verschwendet} (est gaspillé), \all{sind benachteiligt} (sont désavantagés). Au \all{Hauptteil}, au discours indirect, ils deviennent : \all{seien bestochen}, \all{seien schlecht ausgeführt}, \all{sei verschwendet}, \all{seien benachteiligt} (subjonctif I de \all{sein} + participe II).
\end{aretenir}

\courssub{Le texte support (rappel)}

\begin{textebox}[Korruption bei öffentlichen Aufträgen — article original]
Wenn der Staat eine Schule oder eine Straße bauen will, schreibt er den Auftrag öffentlich aus: Mehrere Firmen machen Angebote, und die beste Firma soll den Auftrag bekommen. So will das Gesetz die Transparenz garantieren.

In der Praxis aber wird dieses System oft missbraucht. Manche Beamte sind bestochen: Sie geben den Auftrag nicht der besten, sondern der teuersten Firma — und bekommen dafür Geld unter dem Tisch. Die Folge: Die Arbeiten sind schlecht ausgeführt, das Geld des Staates ist verschwendet, und die ehrlichen Firmen sind benachteiligt.

Ein zweites Problem ist der Schmuggel. Waren werden heimlich über die Grenzen gebracht, ohne dass Zoll gezahlt wird. Noch gefährlicher ist die Fälschung von Produkten: Auf den Märkten werden gefälschte Medikamente verkauft, die krank machen oder sogar töten können.

Experten fordern deshalb strengere Kontrollen, unabhängige Kommissionen bei den Ausschreibungen und härtere Strafen für Betrüger. Außerdem müsse die Bevölkerung aufgeklärt werden, damit sie gefälschte Produkte erkennt und meldet. Nur so, meinen sie, könne das Vertrauen in die Wirtschaft wiederhergestellt werden.
\end{textebox}

\textbf{Glossaire :} \all{der Auftrag, die Aufträge} (le marché, la commande) ; \all{ausschreiben} (mettre en adjudication) ; \all{das Angebot, die Angebote} (l'offre) ; \all{die Transparenz} (la transparence) ; \all{missbrauchen} (détourner, abuser de) ; \all{bestechen / bestochen} (corrompre / corrompu) ; \all{unter dem Tisch} (sous la table, en secret) ; \all{benachteiligen} (désavantager) ; \all{der Zoll, die Zölle} (la douane) ; \all{die Fälschung, die Fälschungen} (la falsification) ; \all{die Strafe, die Strafen} (la sanction) ; \all{aufklären} (sensibiliser) ; \all{melden} (signaler) ; \all{das Vertrauen} (la confiance) ; \all{wiederherstellen} (rétablir).

\courssub{Commentaire modèle rédigé}

Voici un commentaire complet (environ 13 lignes), avec les trois parties signalées. Observe le subjonctif I (souligné mentalement par ta lecture attentive).

\begin{exemplebox}[Modèle de commentaire — Musterkommentar]
\textbf{[Einleitung]} \all{Der vorliegende Text ist ein Zeitungsartikel. Er behandelt das Problem der Korruption, des Schmuggels und der Fälschung von Produkten. Es geht um die Frage: Wie kann man die Wirtschaftskriminalität wirksam bekämpfen?}

\textbf{[Hauptteil]} \all{Zuerst erklärt der Autor, dass öffentliche Aufträge eigentlich transparent ausgeschrieben werden sollten. Dann behauptet er, dieses System werde oft missbraucht: Manche Beamte seien bestochen und gäben den Auftrag der teuersten Firma. Die Folge sei, dass das Geld des Staates verschwendet werde. Als zweites Problem nennt er den Schmuggel und die Fälschung von Produkten; gefälschte Medikamente seien besonders gefährlich, weil sie sogar töten könnten. Schließlich vertritt er die Meinung, strengere Kontrollen und härtere Strafen seien nötig, und die Bevölkerung müsse aufgeklärt werden.}

\textbf{[Schluss]} \all{Meiner Meinung nach hat der Autor recht: Korruption schadet allen Bürgern. Ich bin der Ansicht, dass man in Kamerun die Kontrollen an den Grenzen verstärken sollte, damit gefälschte Produkte nicht mehr auf unsere Märkte kämen.}
\end{exemplebox}

Observe la progression : présentation neutre, puis rapport fidèle des idées au subjonctif I (\all{werde missbraucht, seien bestochen, gäben, sei, werde, seien, könnten, seien, müsse}), et seulement à la fin l'opinion personnelle à l'indicatif (ou KII pour la recommandation : \all{sollte ... kämen}).

\courssub{Méthode et pièges à éviter}

\begin{methode}[Réussir le commentaire au Bac en 5 étapes]
\begin{enumerate}
\item \textbf{Lire le texte deux fois} : une première fois pour le sens global, une deuxième fois en soulignant les idées principales et les connecteurs (\all{zuerst, dann, schließlich}).
\item \textbf{Dégager la problématique en une seule question} : « De quoi s'agit-il et quel est le débat ? » (\all{Wie kann man die Korruption bekämpfen?}).
\item \textbf{Résumer sans citer de longs passages} : reformuler avec ses propres mots ; une courte citation est tolérée, jamais un paragraphe recopié.
\item \textbf{Analyser la position de l'auteur au subjonctif I} : chaque idée rapportée prend la forme du discours indirect (attention aux verbes identiques indicatif/K1 → KII de remplacement).
\item \textbf{Conclure avec une opinion nuancée et justifiée} : \all{Meiner Meinung nach... weil... / denn...} — une opinion sans justification ne vaut rien.
\end{enumerate}
\end{methode}

\begin{attention}[Les quatre erreurs qui coûtent des points]
\begin{itemize}
\item \textbf{Recopier le texte} : le correcteur reconnaît immédiatement les phrases du texte support ; il faut reformuler.
\item \textbf{Donner son opinion dès l'introduction} : l'opinion personnelle n'apparaît que dans le \all{Schluss}.
\item \textbf{Oublier la problématique} : sans question directrice, l'introduction est incomplète.
\item \textbf{Traduire mot à mot depuis le français} : « l'auteur pense que » ne se traduit pas par \all{der Autor denkt, dass} suivi d'un indicatif, mais par \all{der Autor ist der Meinung, ...} + subjonctif I. Attention aussi à \all{behaupten} (affirmer, prétendre — faux ami partiel : ce n'est pas « bêhauter »).
\end{itemize}
\end{attention}

\begin{exoresolu}[Transformer des phrases pour le Hauptteil]
Énoncé : Mets ces phrases au discours indirect (subjonctif I) pour un commentaire, en commençant par \all{Der Autor sagt, ...} :
a) \all{Die Korruption schadet der Wirtschaft.}
b) \all{Gefälschte Medikamente sind gefährlich.}
c) \all{Die Beamten werden bestochen.}
\tcblower
Solution détaillée :
a) \all{Der Autor sagt, die Korruption schade der Wirtschaft.} — \all{schade} : subjonctif I de \all{schaden} (3ᵉ pers. sg.).
b) \all{Der Autor sagt, gefälschte Medikamente seien gefährlich.} — \all{seien} : subjonctif I de \all{sein} (pluriel), distinct de l'indicatif \all{sind}.
c) \all{Der Autor sagt, die Beamten würden bestochen.} — le subjonctif I \all{werden bestochen} est identique à l'indicatif ; on recourt donc au subjonctif II de remplacement \all{würden bestochen} (passif au discours indirect).
\end{exoresolu}

\begin{exercice}[À ton tour]
Rédige un commentaire de 10 à 12 lignes sur le texte support, en respectant le plan \all{Einleitung -- Hauptteil -- Schluss} et en utilisant au moins cinq verbes au subjonctif I. Tu peux t'appuyer sur les phrases-types des tableaux.
\end{exercice}

\begin{corrige}[Exercice — éléments de réponse]
Introduction attendue : type de texte (\all{Zeitungsartikel}), thème (\all{Wirtschaftskriminalität}), problématique sous forme de question. Développement : les idées du texte rapportées au subjonctif I (\all{der Autor behauptet, das System werde missbraucht; er nennt als Beispiel gefälschte Medikamente; er fordert, die Bevölkerung müsse aufgeklärt werden}). Conclusion : opinion personnelle à l'indicatif (ou KII pour recommandation), justifiée par \all{weil} ou \all{denn}, par exemple : \all{Ich bin der Ansicht, dass ehrliche Ausschreibungen die Entwicklung Kameruns fördern, weil das Geld des Staates dann richtig benutzt wird.}
\end{corrige}

\coursec{Production guidée : écrire et parler sur la criminalité économique}

Dans la section précédente, nous avons appris à \cle{commenter} un texte de presse : le lire globalement, relever les idées principales, analyser le vocabulaire et les procédés grammaticaux (Zustandspassiv, discours indirect au subjonctif I). Nous allons maintenant passer de la \emph{réception} à la \emph{production} : répondre par écrit aux questions de compréhension, résumer le texte, rédiger un court paragraphe argumenté et débattre à l'oral. C'est l'aboutissement de la leçon : prouver que vous savez \textbf{réutiliser} le lexique et la grammaire étudiés dans vos propres phrases.

\courssub{Répondre par écrit aux questions de compréhension}

\begin{methode}[Répondre à une question de compréhension]
\begin{enumerate}
\item Lire la question et repérer le \cle{mot interrogatif} : \all{warum} (pourquoi), \all{was} (quoi), \all{wer} (qui), \all{wie} (comment).
\item Repérer dans le texte le passage qui contient la réponse.
\item Rédiger une \textbf{phrase complète} (sujet + verbe conjugué en 2\textsuperscript{e} position), en \textbf{reformulant} avec vos propres mots : ne recopiez jamais une phrase entière du texte.
\item Reprendre dans la réponse un élément de la question : \all{Warum ist die Korruption gefährlich?} $\rightarrow$ \all{Die Korruption ist gefährlich, weil...}
\item Vérifier : place du verbe, cas, majuscules aux noms.
\end{enumerate}
\end{methode}

\begin{exoresolu}[Questions de compréhension sur le texte de la leçon]
Répondez par des phrases complètes.
\begin{enumerate}
\item \all{Was ist ein öffentlicher Auftrag?}
\item \all{Warum ist die Ausschreibung wichtig?}
\item \all{Welche Folgen hat der Schmuggel für den Staat?}
\item \all{Warum sind gefälschte Medikamente besonders gefährlich?}
\item \all{Was fordern die Experten von der Regierung?}
\end{enumerate}
\tcblower
\textbf{Solutions (réponses modèles) :}
\begin{enumerate}
\item \all{Ein öffentlicher Auftrag ist ein Auftrag, den der Staat an eine Firma gibt, zum Beispiel für den Bau einer Straße oder einer Schule.} « Un marché public est un contrat que l'État attribue à une entreprise, par exemple pour la construction d'une route ou d'une école. »
\item \all{Die Ausschreibung ist wichtig, weil sie den Wettbewerb zwischen den Firmen garantiert und die Korruption verhindern soll.} « L'appel d'offres est important car il garantit la concurrence entre les entreprises et doit prévenir la corruption. »
\item \all{Der Schmuggel hat schwere Folgen für den Staat: Er verliert Steuern und Zölle, und die einheimischen Firmen werden geschädigt.} « La contrebande a de lourdes conséquences pour l'État : il perd des impôts et des droits de douane, et les entreprises locales sont lésées. »
\item \all{Gefälschte Medikamente sind besonders gefährlich, weil sie keine Wirkung haben oder sogar giftig sein können; die Patienten werden nicht geheilt.} « Les médicaments falsifiés sont particulièrement dangereux car ils n'ont aucun effet ou peuvent même être toxiques ; les patients ne sont pas soignés. »
\item \all{Die Experten fordern, dass die Regierung die Täter härter bestraft und die Kontrollen an den Grenzen und auf den Märkten verstärkt.} « Les experts exigent que le gouvernement punisse les coupables plus sévèrement et renforce les contrôles aux frontières et sur les marchés. »
\end{enumerate}
\end{exoresolu}

\begin{attention}[La reprise du mot interrogatif \all{warum}]
En allemand, la réponse à \all{warum} se construit avec \all{weil} et le verbe \textbf{à la fin} : \all{..., weil die Kontrollen schwach sind.} Erreur fréquente du francophone : calquer « parce que les contrôles sont faibles » en gardant l'ordre français : *\all{weil die Kontrollen sind schwach} est faux.
\end{attention}

\courssub{Résumé guidé du texte (5 à 6 phrases)}

Un résumé (\all{eine Zusammenfassung}) ne contient que les \textbf{idées essentielles}, sans exemples ni détails. Pour vous guider, les connecteurs sont \textbf{imposés} : \all{zuerst} (d'abord), \all{dann} (ensuite), \all{außerdem} (de plus), \all{jedoch} (cependant), \all{schließlich} (finalement).

\begin{center}
\begin{tabularx}{\linewidth}{|l|X|X|}
\hline
\rowcolor{popblueL} \textbf{Connecteur} & \textbf{Fonction} & \textbf{Place du verbe} \\
\hline
\all{zuerst} & première étape & verbe en 2\textsuperscript{e} position \\
\hline
\all{dann} & étape suivante & verbe en 2\textsuperscript{e} position \\
\hline
\all{außerdem} & ajout d'un argument & verbe en 2\textsuperscript{e} position \\
\hline
\all{jedoch} & opposition & verbe en 2\textsuperscript{e} position \\
\hline
\all{schließlich} & conclusion & verbe en 2\textsuperscript{e} position \\
\hline
\end{tabularx}
\end{center}

\begin{attention}[Jedoch et außerdem ne sont pas des conjonctions de subordination]
Après \all{jedoch} et \all{außerdem} en tête de phrase, le verbe reste en \textbf{deuxième position} : \all{Jedoch sind die Kontrollen schwach.} (Cependant les contrôles sont faibles.) et non *\all{Jedoch die Kontrollen sind schwach.} Ces mots sont des \emph{adverbes} (ils occupent la 1\textsuperscript{re} place), pas des conjonctions comme \all{weil} ou \all{dass}.
\end{attention}

\begin{exemplebox}[Résumé modèle du texte de la leçon]
\all{Zuerst wird erklärt, was öffentliche Aufträge sind und warum die Ausschreibung den Wettbewerb sichern soll. Dann zeigt der Text, dass Korruption dieses System zerstört: Manche Firmen erhalten Aufträge durch Bestechung. Außerdem werden der Schmuggel und die Fälschung von Produkten beschrieben, die dem Staat und den Bürgern schaden. Jedoch sind die Kontrollen oft schwach. Schließlich fordern Experten härtere Strafen und mehr Transparenz.}

« D'abord, on explique ce que sont les marchés publics et pourquoi l'appel d'offres doit garantir la concurrence. Ensuite, le texte montre que la corruption détruit ce système : certaines entreprises obtiennent des marchés par des pots-de-vin. De plus, la contrebande et la falsification des produits, qui nuisent à l'État et aux citoyens, sont décrites. Cependant, les contrôles sont souvent faibles. Finalement, les experts réclament des peines plus dures et plus de transparence. »
\end{exemplebox}

\courssub{Production écrite guidée : rédiger un paragraphe argumenté}

\textbf{Sujet :} \all{Welche Maßnahmen sollte Kamerun gegen Korruption und Fälschung ergreifen?} (Quelles mesures le Cameroun devrait-il prendre contre la corruption et la falsification ?) — 8 à 10 lignes.

\begin{methode}[Rédiger une production guidée en trois temps]
\begin{enumerate}
\item \textbf{Plan en mots-clés allemands} : notez d'abord vos idées en allemand, sans phrases : \all{Problem: gefälschte Medikamente / Ursache: schwache Kontrollen / Maßnahme: härtere Strafen...}
\item \textbf{Phrases simples d'abord} : écrivez des phrases courtes et correctes (sujet – verbe – complément). N'écrivez une phrase complexe (avec \all{weil}, \all{dass}, subordonnée relative) que si vous êtes sûr de la place du verbe.
\item \textbf{Vérification finale} : relisez en contrôlant trois points : les \textbf{accords} (adjectifs, articles), la \textbf{place du verbe} (2\textsuperscript{e} position, verbe à la fin dans les subordonnées) et les \textbf{modes} (subjonctif I pour rapporter une parole, subjonctif II pour un conseil : \all{man sollte...}).
\end{enumerate}
\end{methode}

\begin{popfigure}
\begin{tikzpicture}[node distance=0.55cm and 0.9cm,
box/.style={rectangle, rounded corners=4pt, draw=popblue, fill=popblueL, thick, align=center, minimum width=3.1cm, minimum height=1.15cm, font=\small},
arr/.style={-{Stealth[length=3mm]}, thick, popdark}]
\node[box, align=center] (a){\textbf{Feststellung}\\ \all{Das Problem}};
\node[box, right=of a, draw=poporange, fill=poporangeL, align=center] (b){\textbf{Ursachen}\\ \all{Warum?}};
\node[box, right=of b, draw=popgreen, fill=popgreenL, align=center] (c){\textbf{Maßnahmen}\\ \all{Was tun?}};
\node[box, right=of c, draw=poppurple, fill=poppurpleL, align=center] (d){\textbf{Meinung}\\ \all{Ich denke...}};
\draw[arr] (a) -- (b);
\draw[arr] (b) -- (c);
\draw[arr] (c) -- (d);
\end{tikzpicture}
\legende{Plan de la production guidée : constat $\rightarrow$ causes $\rightarrow$ mesures $\rightarrow$ opinion personnelle.}
\end{popfigure}

\textbf{Contraintes linguistiques obligatoires :}
\begin{itemize}
\item au moins \textbf{deux phrases au Zustandspassiv} (\all{sein} + participe II) : \all{Viele Produkte auf dem Markt sind gefälscht.} (Beaucoup de produits sur le marché sont falsifiés.)
\item au moins \textbf{deux phrases au subjonctif I} (discours indirect) : \all{Experten sagen, die Regierung müsse härter bestrafen.} (Les experts disent que le gouvernement devrait punir plus sévèrement.)
\item au moins \textbf{cinq mots du lexique thématique} : \all{der Markt, die Korruption, der Schmuggel, die Fälschung, die Ausschreibung, der Betrug, die Bestechung, die Kontrolle, die Strafe, die Transparenz}.
\end{itemize}

\begin{textebox}[Production modèle (réponse attendue)]
\all{In Kamerun sind viele Produkte auf den Märkten gefälscht, besonders Medikamente und Lebensmittel. Auch die Korruption ist ein großes Problem: Bei öffentlichen Aufträgen wird oft Geld unter dem Tisch gezahlt. Die Ursachen sind bekannt: Die Kontrollen sind schwach, und die Strafen sind zu mild. Experten sagen, der Staat müsse die Grenzen besser kontrollieren, und die Justiz sei zu langsam. Meiner Meinung nach sollte die Regierung zuerst die Ausschreibungen transparenter machen. Dann sollte sie die Zöllner besser bezahlen, damit sie keine Bestechung akzeptieren. Außerdem sollten gefälschte Produkte sofort vernichtet werden. Schließlich bin ich überzeugt: Wenn die Bürger informiert sind, wird der Kampf gegen Betrug und Fälschung erfolgreich sein.}

« Au Cameroun, beaucoup de produits sur les marchés sont falsifiés, surtout les médicaments et les denrées alimentaires. La corruption aussi est un grand problème : pour les marchés publics, de l'argent est souvent versé sous la table. Les causes sont connues : les contrôles sont faibles et les peines trop légères. Les experts disent que l'État devrait mieux contrôler les frontières et que la justice serait trop lente. À mon avis, le gouvernement devrait d'abord rendre les appels d'offres plus transparents. Ensuite, il devrait mieux payer les douaniers pour qu'ils n'acceptent pas de pots-de-vin. De plus, les produits falsifiés devraient être détruits immédiatement. Finalement, je suis convaincu : si les citoyens sont informés, la lutte contre la fraude et la falsification sera couronnée de succès. »

\textbf{Glossaire :} \all{die Lebensmittel} (pl.) : les denrées alimentaires ; \all{unter dem Tisch zahlen} : payer sous la table (corrompre) ; \all{mild} : léger, clément (pour une peine) ; \all{der Zöllner, -} : le douanier ; \all{vernichten} : détruire ; \all{erfolgreich} : couronné de succès.
\end{textebox}

\begin{exemplebox}[Repérage des contraintes dans le modèle]
\begin{itemize}
\item \textbf{Zustandspassiv} : \all{viele Produkte ... sind gefälscht} ; \all{wenn die Bürger informiert sind}.
\item \textbf{Subjonctif I} (discours indirect) : \all{der Staat müsse ... kontrollieren} ; \all{die Justiz sei zu langsam}.
\item \textbf{Lexique thématique} : \all{Markt, Korruption, öffentliche Aufträge, Bestechung, Ausschreibung, Betrug, Fälschung, Kontrollen, Strafen, Transparenz}.
\end{itemize}
\end{exemplebox}

\courssub{Production orale : débat en binômes}

\begin{experience}[Débat : \all{Ist die Korruption unvermeidbar?} — La corruption est-elle inévitable ?]
Travaillez en binômes. L'élève A défend la thèse « oui », l'élève B la thèse « non ». Chacun pioche trois cartes d'arguments et doit les développer en deux ou trois phrases, puis réagir à l'argument de son partenaire (\all{Da hast du recht, aber...} / \all{Das sehe ich anders: ...}).

\textbf{Cartes d'arguments :}
\begin{itemize}
\item \all{Menschen sind überall gierig.} (Les hommes sont partout avides.)
\item \all{Niedrige Gehälter fördern die Bestechung.} (Les bas salaires favorisent les pots-de-vin.)
\item \all{Andere Länder haben die Korruption besiegt.} (D'autres pays ont vaincu la corruption.)
\item \all{Erziehung und Information verändern das Verhalten.} (L'éducation et l'information changent les comportements.)
\item \all{Ohne Strafen gibt es keine Angst.} (Sans sanctions, il n'y a pas de peur / de dissuasion.)
\item \all{Die Tradition des Schenkens ist stark.} (La tradition des cadeaux est forte.)
\end{itemize}

\textbf{Expressions utiles pour débattre :} \all{Meiner Meinung nach...} (à mon avis) ; \all{Ich bin der Ansicht, dass...} (j'estime que) ; \all{Das stimmt, jedoch...} (c'est vrai, cependant) ; \all{Ich bin nicht einverstanden, weil...} (je ne suis pas d'accord parce que).
\end{experience}

\courssub{Grille d'auto-évaluation}

Après votre production écrite, évaluez-vous honnêtement avec la grille suivante (cochez oui / non) :

\begin{center}
\begin{tabularx}{\linewidth}{|l|X|l|}
\hline
\rowcolor{popblueL} \textbf{Critère} & \textbf{Question de contrôle} & \textbf{Oui / Non} \\
\hline
\cle{Lexique} & Ai-je utilisé au moins cinq mots du thème avec le bon article ? & \\
\hline
\cle{Zustandspassiv} & Ai-je écrit deux phrases avec \all{sein} + participe II ? & \\
\hline
\cle{Subjonctif I} & Ai-je rapporté deux paroles au subjonctif I (\all{müsse}, \all{sei}, \all{habe}) ? & \\
\hline
\cle{Structure} & Mon texte suit-il le plan : constat $\rightarrow$ causes $\rightarrow$ mesures $\rightarrow$ opinion ? & \\
\hline
\cle{Connecteurs} & Ai-je placé le verbe en 2\textsuperscript{e} position après \all{jedoch} et \all{außerdem} ? & \\
\hline
\cle{Orthographe} & Tous les noms ont-ils une majuscule ? & \\
\hline
\end{tabularx}
\end{center}

\begin{exercice}[À vous de produire]
Rédigez votre propre paragraphe de 8 à 10 lignes sur le sujet : \all{Welche Maßnahmen sollte Kamerun gegen Korruption und Fälschung ergreifen?} Respectez le plan en quatre étapes, les contraintes grammaticales (deux Zustandspassiv, deux subjonctifs I) et utilisez cinq mots du lexique thématique. Puis remplissez la grille d'auto-évaluation et préparez le débat oral avec votre voisin.
\end{exercice}

\begin{aretenir}[L'essentiel de la production guidée]
Une bonne production en allemand repose sur trois piliers : un \textbf{plan en mots-clés} avant de rédiger, des \textbf{phrases simples et correctes} plutôt que des phrases longues et fausses, et une \textbf{relecture systématique} (accords, place du verbe, modes). Les contraintes imposées — Zustandspassiv, subjonctif I, lexique du thème — ne sont pas des obstacles : ce sont les preuves que vous maîtrisez la leçon.
\end{aretenir}

\newpage

\coursec{Activité d'intégration}

\courssub{Objectifs de l'activité}

À la fin de cette activité d'intégration, tu seras capable de :

\begin{itemize}
\item mobiliser dans une situation de communication authentique le \cle{vocabulaire thématique} des marchés publics et de la criminalité économique (\all{die Ausschreibung}, \all{die Korruption}, \all{der Schmuggel}, \all{die Fälschung}, \all{der Betrug}, \all{beschlagnahmen}) ;
\item poser des questions orales précises et polies, comme un journaliste lors d'une conférence de presse ;
\item rapporter des déclarations au \cle{discours indirect} en employant correctement le \cle{subjonctif I} (\all{Der Minister sagte, das Verfahren sei korrekt gewesen.}) ;
\item décrire des faits et des résultats au \cle{passif d'état} (Zustandspassiv : \all{Die Waren sind beschlagnahmt.}) ;
\item rédiger un court article de presse de 10 à 12 lignes en allemand, clair, structuré et correct ;
\item prendre la parole en public avec assurance et écouter activement les interventions des autres.
\end{itemize}

\courssub{Situation}

Tu es journaliste à Yaoundé. Hier soir, la douane camerounaise a annoncé une vaste opération au port de Douala : des containers remplis de médicaments falsifiés ont été saisis. Or ces médicaments avaient été achetés dans le cadre d'un marché public de fourniture pour les hôpitaux de la région du Littoral. Des rumeurs de corruption entourent l'attribution du marché. Le ministère concerné organise aujourd'hui une conférence de presse. Voici la dépêche d'agence qui a déclenché le scandale :

\begin{textebox}[Dépêche d'agence — Douala, ce matin]
Große Schmugglerware im Hafen von Douala entdeckt

In der Nacht von Dienstag auf Mittwoch haben Zollbeamte im Hafen von Douala drei Container beschlagnahmt. Die Container waren als „medizinische Geräte" deklariert, enthielten aber gefälschte Medikamente: Antibiotika und Malariamittel ohne Wirkstoff. Der Wert der Ware wird auf mehrere Millionen Euro geschätzt.

Besonders brisant: Die gefälschten Medikamente waren für ein öffentliches Krankenhaus bestimmt. Der Auftrag war im Rahmen einer öffentlichen Ausschreibung an eine Firma vergeben worden, die erst vor sechs Monaten gegründet wurde. Oppositionelle Politiker sprechen von Korruption und Betrug und fordern eine unabhängige Untersuchung.

Der zuständige Minister und der Chef der Zollpolizei geben heute um 10 Uhr eine Pressekonferenz.
\end{textebox}

\textbf{Glossaire :} \all{der Zollbeamte, -n} (le douanier) ; \all{beschlagnahmen} (saisir, confisquer) ; \all{deklariert} (déclaré) ; \all{der Wirkstoff, -e} (le principe actif) ; \all{der Auftrag, „-e} (le marché, la commande) ; \all{die Ausschreibung, -en} (l'appel d'offres) ; \all{vergeben — vergab — vergeben} (attribuer) ; \all{brisant} (explosif, délicat) ; \all{die Untersuchung, -en} (l'enquête).

\courssub{Consignes}

\textbf{Étape 1 — Compréhension du document (10 min, travail individuel).}
Lis la dépêche attentivement, puis réponds par écrit en allemand, par des phrases complètes :
\begin{enumerate}
\item \all{Was wurde im Hafen von Douala entdeckt?}
\item \all{Warum ist dieser Skandal besonders schwerwiegend?}
\item \all{Wer wird verdächtigt, und warum?}
\end{enumerate}

\textbf{Étape 2 — Préparation des rôles (15 min, travail en groupes).}
La classe se divise en trois groupes :
\begin{enumerate}
\item \textbf{Groupe A — les journalistes.} Préparez au moins huit questions pour la conférence de presse. Utilisez des questions fermées avec \all{ist es wahr, dass ...?}, \all{können Sie bestätigen, dass ...?} et des questions ouvertes avec \all{warum}, \all{seit wann}, \all{wie}. Employez le vocabulaire du thème : \all{die Ausschreibung}, \all{die Bestechung}, \all{der Schmuggel}, \all{die Fälschung}.
\item \textbf{Groupe B — le ministre.} Préparez vos réponses : vous défendez la procédure, mais vous annoncez une enquête. Prévoyez des formulations au passif : \all{Der Auftrag wurde korrekt vergeben. Die Akten sind geprüft.}
\item \textbf{Groupe C — le chef de la police douanière.} Préparez un bref rapport factuel : que s'est-il passé, qu'est-ce qui est saisi, fermé, arrêté ? Utilisez le Zustandspassiv pour décrire l'état actuel : \all{Die Container sind geöffnet, die Waren sind beschlagnahmt, zwei Verdächtige sind festgenommen.}
\end{enumerate}

\textbf{Étape 3 — La conférence de presse (20 min, oral).}
Le ministre et le chef de la douane s'installent devant la classe. Les journalistes posent leurs questions ; chacun prend des notes sur les réponses. Le professeur joue le modérateur et veille à la correction de la langue.

\textbf{Étape 4 — Rédaction de l'article (25 min, travail individuel).}
Chaque « journaliste » rédige un article de 10 à 12 lignes qui :
\begin{enumerate}
\item rapporte au moins \textbf{trois déclarations au discours indirect} (subjonctif I) : \all{Der Minister sagte, das Verfahren sei korrekt gewesen. Er erklärte, man habe nichts gewusst.} ;
\item décrit au moins \textbf{deux faits au passif d'état} : \all{Die Waren sind beschlagnahmt. Die Firma ist geschlossen.} ;
\item contient au moins \textbf{cinq mots du lexique thématique} de la leçon ;
\item commence par un titre et une phrase d'introduction (qui ? quoi ? où ? quand ?).
\end{enumerate}

\textbf{Étape 5 — Publication (10 min).}
Les deux ou trois meilleurs articles sont lus à voix haute par leurs auteurs, corrigés collectivement, puis affichés au mur de la classe.

\courssub{Ressources à mobiliser}

\begin{aretenir}[Boîte à outils linguistique]
\textbf{Lexique :} \all{der Markt, die Märkte} ; \all{die Ausschreibung, -en} (l'appel d'offres) ; \all{der Auftrag, „-e} (le marché) ; \all{die Korruption} ; \all{die Bestechung, -en} (la corruption active, les pots-de-vin) ; \all{der Schmuggel} (la contrebande) ; \all{die Fälschung, -en} (la falsification) ; \all{der Betrug} (l'escroquerie) ; \all{beschlagnahmen} (saisir) ; \all{festnehmen} (arrêter) ; \all{der Verdächtige, -n} (le suspect).

\textbf{Zustandspassiv :} \emph{sein} + participe II. Il décrit un \textbf{état résultant} d'une action : \all{Die Container sind geöffnet.} (les containers sont ouverts — résultat). Ne pas confondre avec le Vorgangspassiv (\emph{werden} + participe II) qui décrit l'action : \all{Die Container werden geöffnet.}

\textbf{Discours indirect — subjonctif I :} \all{Der Minister sagte, das Verfahren sei korrekt gewesen.} Pour le passé au subjonctif I : \emph{habe/sei} + participe II. Si la forme du subjonctif I est identique à l'indicatif (\all{sie haben}), on la remplace par le subjonctif II : \all{Er sagte, sie hätten nichts gewusst.}
\end{aretenir}

\begin{attention}[Pièges à éviter]
\begin{itemize}
\item Ne dis pas \all{Der Minister sagte, dass das Verfahren ist korrekt.} — au discours indirect, le verbe passe au subjonctif I : \all{... dass das Verfahren korrekt sei.}
\item Attention à la place du verbe : après \all{dass}, le verbe conjugué va à la fin ; sans \all{dass}, il reste en deuxième position : \all{Er sagte, er habe nichts gewusst.}
\item \all{bekommen} ne signifie pas « devenir » mais « recevoir » : \all{Die Firma bekam den Auftrag.} (la firme a obtenu le marché).
\item N'oublie pas la majuscule à tous les noms : \all{die Korruption}, \all{der Schmuggel}, \all{die Pressekonferenz}.
\end{itemize}
\end{attention}

\courssub{Proposition de production et corrigé commenté}

\begin{exoresolu}[Article modèle : le scandale de Douala]
Rédige un article de presse de 10 à 12 lignes rapportant la conférence de presse, avec au moins trois déclarations au subjonctif I et deux faits au Zustandspassiv.
\tcblower
\textbf{Article modèle :}

\all{Skandal um gefälschte Medikamente: Minister verteidigt Ausschreibung}

\all{Auf einer Pressekonferenz in Yaoundé hat der Minister gestern die Vorwürfe der Korruption zurückgewiesen. Er sagte, die Ausschreibung sei nach allen Regeln gelaufen, und der Auftrag sei korrekt vergeben worden. Er habe von den gefälschten Medikamenten nichts gewusst. Der Chef der Zollpolizei erklärte, die Beamten hätten in der Nacht drei Container im Hafen von Douala geöffnet. Die Waren seien als medizinische Geräte deklariert gewesen, hätten aber Antibiotika ohne Wirkstoff enthalten. Heute sind die Medikamente beschlagnahmt, die Firma ist geschlossen und zwei Verdächtige sind festgenommen. Der Minister kündigte an, man werde eine unabhängige Untersuchung einleiten. Die Opposition bleibt skeptisch und spricht weiter von Betrug.}

\medskip
\textbf{Commentaire des choix de langue :}
\begin{enumerate}
\item \all{die Ausschreibung sei nach allen Regeln gelaufen} : discours indirect au passé — subjonctif I de \emph{sein} (\all{sei}) + participe II (\all{gelaufen}), car \all{laufen} forme son passé avec \emph{sein}.
\item \all{der Auftrag sei korrekt vergeben worden} : passif au passé rapporté au discours indirect — \all{worden} (et non \all{geworden}) est le participe II de \emph{werden} au passif.
\item \all{Er habe ... nichts gewusst} : subjonctif I de \emph{haben} + participe II ; la phrase sans \all{dass} garde le verbe en deuxième position.
\item \all{die Beamten hätten ... geöffnet} : ici le subjonctif I (\all{haben}) serait identique à l'indicatif à la 3\textsuperscript{e} personne du pluriel ; on utilise donc le subjonctif II (\all{hätten}) pour marquer le discours rapporté.
\item \all{sind beschlagnahmt}, \all{ist geschlossen}, \all{sind festgenommen} : Zustandspassiv — \emph{sein} au présent + participe II, pour décrire l'état actuel résultant des opérations.
\item Lexique thématique mobilisé : \all{die Ausschreibung}, \all{der Auftrag}, \all{die Korruption}, \all{der Betrug}, \all{beschlagnahmen}, \all{die Untersuchung}.
\end{enumerate}
\end{exoresolu}

\courssub{Bilan de l'activité}

Cette activité t'a permis de réinvestir l'ensemble des acquis de la leçon dans une situation de communication réaliste et motivante :

\begin{itemize}
\item \textbf{Compréhension :} tu as lu et exploité une dépêche de presse sur un scandale économique, en repérant les informations essentielles (qui, quoi, où, pourquoi).
\item \textbf{Lexique :} tu as employé activement le vocabulaire des marchés publics et de la criminalité économique, à l'oral comme à l'écrit.
\item \textbf{Grammaire :} tu as distingué et utilisé le Vorgangspassiv et le Zustandspassiv, et tu as rapporté des paroles au discours indirect avec le subjonctif I (et le subjonctif II de substitution).
\item \textbf{Expression orale :} tu as formulé des questions de journaliste et pris la parole en public.
\item \textbf{Expression écrite :} tu as rédigé un article de presse structuré, conforme aux conventions du genre (titre, introduction, citations rapportées).
\end{itemize}

\begin{methode}[Pour aller plus loin]
Pour t'entraîner davantage : choisis un fait divers économique dans un journal camerounais, souligne les verbes au passif, puis réécris trois phrases au discours indirect comme si tu rapportais les déclarations d'un responsable. Vérifie chaque fois : le verbe est-il au subjonctif I ? La forme est-elle distincte de l'indicatif ? Sinon, passe au subjonctif II.
\end{methode}

\newpage

\coursec{Méthodes et exercices résolus}

\courssub{Savoir-faire 1 : Comprendre un texte sur les marchés publics et la criminalité économique}

\begin{methode}[Comprendre un texte économique en trois lectures]
\begin{enumerate}
\item \textbf{Lecture globale (survol).} Lire le titre, la source et le premier paragraphe. Identifier le \cle{thème} (marchés publics, corruption, contrebande, falsification), le \cle{type de texte} (article de presse, reportage, interview) et la \cle{thèse} de l'auteur. Ne pas s'arrêter sur les mots inconnus.
\item \textbf{Lecture détaillée.} Souligner les mots-clés du champ lexical : \all{die Ausschreibung} (l'appel d'offres), \all{der Auftrag} (le marché, la commande), \all{die Bestechung} (la corruption active), \all{der Schmuggel} (la contrebande), \all{gefälschte Produkte} (produits falsifiés). Repérer les chiffres, les noms propres, les connecteurs (\all{weil, obwohl, damit, jedoch}).
\item \textbf{Lecture analytique.} Pour chaque question de compréhension, localiser le passage concerné, reformuler avec ses propres mots en allemand — \textbf{jamais recopier} une phrase entière du texte sans la citer.
\item \textbf{Réflexe de réponse.} Commencer par une phrase de reprise de la question : \all{Der Text erklärt, dass...} / \all{Nach dem Autor...} / \all{Laut dem Artikel...}. Utiliser le discours indirect (subjonctif I) pour rapporter les propos d'un témoin cité dans le texte.
\end{enumerate}
\end{methode}

\begin{exoresolu}[Compréhension de texte — type Bac]
\textbf{Texte (extrait) :} \all{In vielen afrikanischen Ländern werden öffentliche Aufträge nicht immer transparent vergeben. Ein Beamter aus Yaoundé erklärte, manche Firmen erhielten Aufträge, weil sie Bestechungsgelder gezahlt hätten. Außerdem werden an den Grenzen täglich gefälschte Medikamente beschlagnahmt, die für die Bevölkerung gefährlich sind.}

\textbf{Questions :} 1. Wie werden öffentliche Aufträge laut dem Text manchmal vergeben? 2. Was hat der Beamte aus Yaoundé erklärt? 3. Was geschieht an den Grenzen?
\tcblower
\textbf{Raisonnement.} Question 1 : le passage clé est \all{nicht immer transparent vergeben} — on reformule sans recopier. Question 2 : il s'agit de rapporter une parole, donc \textbf{discours indirect au subjonctif} : le texte donne déjà \all{erhielten} (Konj. II, forme de remplacement du Konj. I \all{erhalten}, identique à l'indicatif) et \all{hätten gezahlt} (plus-que-parfait du subjonctif II, remplaçant le Konj. I \all{hätten gezahlt}... identique à l'indicatif : le remplacement est donc obligatoire). Question 3 : repérer le passif d'action \all{werden ... beschlagnahmt}.

\textbf{Réponses rédigées :}
\begin{enumerate}
\item \all{Laut dem Text werden öffentliche Aufträge manchmal ohne Transparenz vergeben, das heißt nicht nach fairen Regeln.} (Selon le texte, les marchés publics sont parfois attribués sans transparence, c'est-à-dire pas selon des règles équitables.)
\item \all{Der Beamte erklärte, manche Firmen hätten Aufträge erhalten, weil sie Bestechungsgelder gezahlt hätten.} (Le fonctionnaire a déclaré que certaines entreprises avaient obtenu des marchés parce qu'elles avaient versé des pots-de-vin.) Justification : après \all{erklärte}, on rapporte au subjonctif ; \all{hätten erhalten} et \all{hätten gezahlt} sont les formes de remplacement (Konj. II) car le Konj. I serait identique à l'indicatif.
\item \all{An den Grenzen werden jeden Tag gefälschte Medikamente beschlagnahmt. Diese Produkte sind gefährlich für die Bevölkerung.} (Aux frontières, des médicaments falsifiés sont saisis chaque jour. Ces produits sont dangereux pour la population.)
\end{enumerate}
\textbf{Conclusion.} Une bonne réponse reformule, cite le texte entre guillemets si nécessaire et utilise le subjonctif pour rapporter les paroles.
\end{exoresolu}

\courssub{Savoir-faire 2 : Employer le vocabulaire des marchés publics et de la criminalité économique}

\begin{methode}[Construire et utiliser le champ lexical]
\begin{enumerate}
\item \textbf{Apprendre chaque nom avec son article et son pluriel} : \all{der Markt, die Märkte} ; \all{die Ausschreibung, -en} (appel d'offres) ; \all{der Auftrag, die Aufträge} (marché/commande) ; \all{die Korruption} (sans pluriel) ; \all{der Schmuggel} (sans pluriel) ; \all{die Fälschung, -en} ; \all{der Betrug} (escroquerie, en général sans pluriel) ; \all{das Bestechungsgeld, -er} (pot-de-vin).
\item \textbf{Mémoriser les familles de mots} : \all{fälschen} (verbe) $\rightarrow$ \all{die Fälschung} (nom) $\rightarrow$ \all{gefälscht} (participe/adjectif) ; \all{schmuggeln} $\rightarrow$ \all{der Schmuggel} $\rightarrow$ \all{der Schmuggler, -} (contrebandier) ; \all{bestechen} (verbe fort : \all{er besticht, er bestach, er hat bestochen}) $\rightarrow$ \all{die Bestechung} $\rightarrow$ \all{bestechlich} (corruptible).
\item \textbf{Associer verbes et compléments} : \all{einen Auftrag vergeben} (attribuer un marché), \all{an einer Ausschreibung teilnehmen} (participer à un appel d'offres), \all{Waren schmuggeln} (faire passer des marchandises en contrebande), \all{Produkte fälschen}, \all{Waren beschlagnahmen} (saisir des marchandises), \all{jemanden bestechen} (corrompre quelqu'un).
\item \textbf{Réutiliser le lexique dans sa production} : au moins cinq mots du thème dans toute rédaction, avec les bons articles et les bons cas.
\end{enumerate}
\end{methode}

\begin{exoresolu}[Vocabulaire — complétion et réemploi]
\textbf{Énoncé :} Complétez avec le mot correct (article + nom) : \all{der Schmuggel / die Ausschreibung / das Bestechungsgeld / die Fälschung / der Auftrag}.\\
1. \all{Die Regierung hat eine neue \_\_\_ für den Bau einer Schule in Douala veröffentlicht.} 2. \all{Die Firma hat den \_\_\_ bekommen, obwohl ihr Angebot zu teuer war.} 3. \all{Der Beamte hat \_\_\_ angenommen und wurde verhaftet.} 4. \all{An der Grenze wird der \_\_\_ von Medikamenten bekämpft.} 5. \all{Die \_\_\_ von Markenprodukten ist ein großes Problem.}
\tcblower
\textbf{Raisonnement.} On identifie d'abord le genre grammatical exigé par le contexte, puis le sens. 1. \all{eine neue ... veröffentlicht} : nom féminin, sens « appel d'offres » $\rightarrow$ \textbf{die Ausschreibung}. 2. \all{den \_\_\_ bekommen} : accusatif masculin, sens « marché » $\rightarrow$ \textbf{den Auftrag} (verbe \all{bekommen} + Akkusativ). 3. \all{hat ... angenommen} : accusatif neutre $\rightarrow$ \textbf{das Bestechungsgeld}. 4. \all{der \_\_\_ von Medikamenten} : nominatif masculin $\rightarrow$ \textbf{der Schmuggel}. 5. \all{Die \_\_\_ ... ist} : nominatif féminin $\rightarrow$ \textbf{die Fälschung}.

\textbf{Phrase de réemploi (production libre corrigée) :} \all{Nach der Ausschreibung hat eine Firma den Auftrag erhalten, weil sie den Beamten bestochen hatte; später wurde der Schmuggel gefälschter Produkte entdeckt.} (Après l'appel d'offres, une entreprise a obtenu le marché parce qu'elle avait corrompu le fonctionnaire ; plus tard, la contrebande de produits falsifiés a été découverte.) Justification : \all{bestechen} au plus-que-parfait (\all{hatte bestochen}) pour l'antériorité ; \all{gefälschter Produkte} : génitif pluriel après \all{der Schmuggel} ; attention à l'orthographe : \all{gefälschter} avec \textbf{ä}.
\end{exoresolu}

\courssub{Savoir-faire 3 : Appliquer la grammaire — le passif d'état (Zustandspassiv)}

\begin{methode}[Former et employer le Zustandspassiv]
\begin{enumerate}
\item \textbf{Former} : \cle{sein} (conjugé au temps voulu) + \cle{participe II} : \all{Die Grenze ist geschlossen.} (La frontière est fermée.) Präteritum : \all{war geschlossen} ; Perfekt : \all{ist geschlossen gewesen}.
\item \textbf{Distinguer du Vorgangspassiv} (passif d'action) : le passif d'action avec \all{werden} décrit \textbf{l'action en train de se faire} (\all{Die Ware wird beschlagnahmt.} — on est en train de saisir la marchandise) ; le passif d'état avec \all{sein} décrit \textbf{le résultat, l'état} (\all{Die Ware ist beschlagnahmt.} — la marchandise est saisie, c'est fait).
\item \textbf{Transformation} : pour passer du passif d'action au passif d'état, remplacer \all{werden} par \all{sein} au même temps et garder le participe II.
\item \textbf{Réflexe de traduction} : le français « est + participe passé » se traduit souvent par le Zustandspassiv quand il décrit un état : « le marché est attribué » $\rightarrow$ \all{Der Auftrag ist vergeben.} ; « les produits sont falsifiés » $\rightarrow$ \all{Die Produkte sind gefälscht.}
\end{enumerate}
\end{methode}

\begin{exoresolu}[Grammaire — transformation au passif d'état]
\textbf{Énoncé :} Transformez au Zustandspassiv (Präsens), puis traduisez en français.\\
1. \all{Die gefälschten Medikamente werden beschlagnahmt.} 2. \all{Der Auftrag wird an die beste Firma vergeben.} 3. \all{Die Grenze wird nachts kontrolliert.}
\tcblower
\textbf{Raisonnement.} Règle appliquée : \all{werden} $\rightarrow$ \all{sein} au Präsens (ist/sind), le participe II reste inchangé en fin de phrase.

\textbf{Solutions :}
\begin{enumerate}
\item \all{Die gefälschten Medikamente \textbf{sind} beschlagnahmt.} — « Les médicaments falsifiés sont saisis » (état : la saisie est accomplie).
\item \all{Der Auftrag \textbf{ist} an die beste Firma vergeben.} — « Le marché est attribué à la meilleure entreprise. » Justification : sujet singulier masculin $\rightarrow$ \all{ist} ; le participe \all{vergeben} garde son préfixe inséparable \all{ver-} sans \all{ge-}.
\item \all{Die Grenze \textbf{ist} nachts kontrolliert.} — « La frontière est contrôlée la nuit. »
\end{enumerate}
\textbf{Conclusion.} Le Zustandspassiv insiste sur le résultat : il est idéal pour décrire une situation constatée (marchandises saisies, dossiers falsifiés, frontières fermées). Ne jamais écrire \all{ist geworden} pour un état présent : c'est un contresens (\all{ist ... geworden} = « est devenu », passé d'une action).
\end{exoresolu}

\courssub{Savoir-faire 4 : Appliquer la grammaire — le discours indirect au subjonctif I}

\begin{methode}[Rapporter une parole au Konjunktiv I]
\begin{enumerate}
\item \textbf{Repérer le verbe introducteur} : \all{sagen, erklären, berichten, behaupten, meinen}. Après ces verbes, la parole rapportée passe au \cle{subjonctif I}.
\item \textbf{Former le subjonctif I} : radical de l'infinitif + désinences (-e, -est, -e, -en, -et, -en) : \all{er sage, sie komme, sie haben}. Seul \all{sein} est irrégulier : \all{ich sei, du seist, er sei, wir seien, ihr seiet, sie seien}.
\item \textbf{Forme de remplacement} : quand le subjonctif I est identique à l'indicatif (\all{sie haben, sie gehen}), on utilise le \cle{subjonctif II} : \all{sie hätten, sie gingen}.
\item \textbf{Gérer les temps} : présent $\rightarrow$ Konj. I présent (\all{er habe}) ; passé $\rightarrow$ Konj. I parfait (\all{er habe gezahlt}) ; futur $\rightarrow$ \all{er werde zahlen}.
\item \textbf{Adapter les pronoms et possessifs} : \all{ich} devient \all{er/sie} ; \all{mein} devient \all{sein/ihr}. La construction \all{..., dass + verbe en fin} est possible, mais la construction sans \all{dass} (verbe en position 2) est plus élégante.
\end{enumerate}
\end{methode}

\begin{exoresolu}[Version et discours indirect — type Bac]
\textbf{Énoncé :} a) Mettez au discours indirect : \all{Der Minister sagte: „Wir haben die korrupten Beamten entlassen und wir werden den Schmuggel an den Grenzen bekämpfen.“} b) Traduisez en allemand : « Le journaliste a déclaré que les produits falsifiés étaient dangereux pour les consommateurs. »
\tcblower
\textbf{Raisonnement a).} Deux temps à convertir : \all{haben entlassen} (Perfekt $\rightarrow$ parfait du subjonctif : \all{hätten entlassen} — ici le subjonctif II est obligatoire car le Konj. I \all{sie haben entlassen} serait identique à l'indicatif) ; \all{werden bekämpfen} (futur $\rightarrow$ Konj. I : \all{sie werden bekämpfen}, forme distincte de l'indicatif \all{werden}... attention : \all{sie werden} est identique à l'indicatif, donc on préfère la forme de remplacement \all{sie würden bekämpfen}). Pronom : \all{wir} $\rightarrow$ \all{sie}.

\textbf{Solution a) :} \all{Der Minister sagte, sie hätten die korrupten Beamten entlassen und sie würden den Schmuggel an den Grenzen bekämpfen.} (Le ministre a dit qu'ils avaient licencié les fonctionnaires corrompus et qu'ils combattraient la contrebande aux frontières.) Variante acceptée pour le futur : \all{..., und sie würden den Schmuggel an den Grenzen bekämpfen} est la forme la plus sûre à l'écrit.

\textbf{Raisonnement b).} « a déclaré » $\rightarrow$ \all{erklärte} ; « étaient dangereux » : imparfait rapporté $\rightarrow$ Konj. I de \all{sein} : \all{seien} ; « pour les consommateurs » $\rightarrow$ \all{für die Verbraucher} (accusatif après \all{für}).

\textbf{Solution b) :} \all{Der Journalist erklärte, die gefälschten Produkte seien gefährlich für die Verbraucher.} (ou avec \all{dass} : \all{Der Journalist erklärte, dass die gefälschten Produkte für die Verbraucher gefährlich seien.}) Justification : \all{seien} = Konj. I de \all{sein}, 3e personne du pluriel ; majuscules aux noms \all{Produkte, Verbraucher}.
\end{exoresolu}

\courssub{Savoir-faire 5 : Conjugaison — tableaux de synthèse du subjonctif I et II}

\begin{propriete}[Formation du Konjunktiv I et du Konjunktiv II]
\textbf{Konjunktiv I} (discours indirect) : radical de l'infinitif + \textbf{-e, -est, -e, -en, -et, -en}. Seule exception : \all{sein} (\all{sei, seist, sei, seien, seiet, seien}), sans -e à la 1re et 3e personne du singulier.

\textbf{Konjunktiv II} (irréel, remplacement) : radical du \textbf{Präteritum} + \textbf{-e, -est, -e, -en, -et, -en}, avec Umlaut pour les verbes forts à voyelle a/o/u : \all{ich käme, er gäbe, sie wären}. Pour la plupart des verbes, on préfère \all{würde} + infinitif : \all{ich würde kommen}.
\end{propriete}

\begin{center}
\begin{tabular}{|l|l|l|l|}
\hline
\rowcolor{popblueL} Personne & Konj. I : \all{haben} & Konj. II : \all{haben} & Konj. I : \all{sein} \\
\hline
ich & habe & hätte & sei \\
\hline
du & habest & hättest & seist \\
\hline
er/sie/es & habe & hätte & sei \\
\hline
wir & haben & hätten & seien \\
\hline
ihr & habet & hättet & seiet \\
\hline
sie/Sie & haben & hätten & seien \\
\hline
\end{tabular}
\end{center}

\begin{center}
\begin{tabular}{|l|l|l|l|}
\hline
\rowcolor{popblueL} Personne & Konj. I : \all{werden} & Konj. II : \all{kommen} & Konj. II : \all{geben} \\
\hline
ich & werde & käme / würde kommen & gäbe \\
\hline
du & werdest & kämest / würdest kommen & gäbest \\
\hline
er/sie/es & werde & käme / würde kommen & gäbe \\
\hline
wir & werden & kämen / würden kommen & gäben \\
\hline
ihr & werdet & kämet / würdet kommen & gäbet \\
\hline
sie/Sie & werden & kämen / würden kommen & gäben \\
\hline
\end{tabular}
\end{center}

\begin{aretenir}[Les trois temps du discours indirect]
\begin{itemize}
\item \textbf{Présent} : Konj. I présent — \all{Er sagt, er habe keine Zeit.} (Il dit qu'il n'a pas le temps.)
\item \textbf{Passé} (Perfekt, Präteritum ou Plusquamperfekt du discours direct) : Konj. I parfait = \all{habe/sei} + participe II — \all{Sie erklärte, sie habe das Geld nicht gesehen.} (Elle a déclaré qu'elle n'avait pas vu l'argent.)
\item \textbf{Futur} : \all{werde} + infinitif — \all{Er sagte, er werde die Wahrheit sagen.} (Il a dit qu'il dirait la vérité.)
\end{itemize}
Remplacement par le Konj. II dès que le Konj. I est identique à l'indicatif : \all{sie haben} $\rightarrow$ \all{sie hätten} ; \all{sie werden} $\rightarrow$ \all{sie würden}.
\end{aretenir}

\courssub{Savoir-faire 6 : Produire un court texte argumenté sur la criminalité économique}

\begin{methode}[Rédiger un paragraphe d'opinion (8 à 12 lignes)]
\begin{enumerate}
\item \textbf{Introduction} : une phrase de situation générale avec le vocabulaire du thème : \all{Die Korruption und die Fälschung von Produkten sind große Probleme in vielen Ländern.}
\item \textbf{Développement} : deux ou trois arguments reliés par des connecteurs : \all{zuerst, außerdem, schließlich, deshalb, obwohl}. Employer au moins une fois le \textbf{passif d'état} (\all{Viele Aufträge sind schon vergeben, bevor die Ausschreibung beginnt.}) et une fois le \textbf{discours indirect} (\all{Ein Experte meint, die Kontrollen seien zu schwach.}).
\item \textbf{Conclusion} : proposition de solution avec \all{man sollte...} ou le subjonctif II : \all{Man sollte die Beamten besser kontrollieren.} / \all{Es wäre gut, wenn die Regierung strenger wäre.}
\item \textbf{Relecture systématique} : verbe en position 2 (ou en fin dans les subordonnées), majuscules aux noms, articles corrects, Umlaute (\all{gefälscht, Fälschung}).
\end{enumerate}
\end{methode}

\begin{exoresolu}[Expression écrite guidée — type Bac]
\textbf{Énoncé :} Rédigez un court paragraphe (8-10 lignes) : \all{Wie kann man Korruption und Fälschung in Ihrem Land bekämpfen?}
\tcblower
\textbf{Plan suivi :} 1) constat ; 2) deux mesures concrètes ; 3) opinion personnelle avec subjonctif II.

\textbf{Production modèle :}\\
\all{Die Korruption und der Schmuggel gefälschter Produkte sind ernste Probleme in meinem Land. Zuerst sollte die Regierung die öffentlichen Aufträge transparent vergeben: jede Ausschreibung muss veröffentlicht werden. Außerdem sind viele gefälschte Medikamente schon auf dem Markt, deshalb müssen die Kontrollen an den Grenzen strenger sein. Ein Experte hat erklärt, die bestechlichen Beamten seien ein großes Hindernis für die Entwicklung. Meiner Meinung nach wäre es gut, wenn die Bürger die Korruption melden würden. Schließlich sollte man die jungen Leute in den Schulen über diese Gefahren informieren.}

(Traduction : « La corruption et la contrebande de produits falsifiés sont de graves problèmes dans mon pays. D'abord, le gouvernement devrait attribuer les marchés publics de manière transparente : chaque appel d'offres doit être publié. De plus, beaucoup de médicaments falsifiés sont déjà sur le marché, c'est pourquoi les contrôles aux frontières doivent être plus stricts. Un expert a déclaré que les fonctionnaires corruptibles étaient un grand obstacle au développement. À mon avis, il serait bon que les citoyens dénoncent la corruption. Enfin, il faudrait informer les jeunes dans les écoles sur ces dangers. »)

\textbf{Points grammaticaux vérifiés :} passif d'action au Konj. II potentiel \all{muss veröffentlicht werden} ; discours indirect \all{seien} ; subjonctif II \all{wäre, würden melden} ; verbe conjugué en fin de subordonnée dans \all{wenn die Bürger die Korruption melden würden} ; génitif \all{der Schmuggel gefälschter Produkte} ; datif pluriel \all{in den Schulen}.
\end{exoresolu}

\begin{attention}[Les 5 erreurs qui coûtent des points au Bac]
\begin{enumerate}
\item \textbf{Confondre passif d'action et passif d'état.} Écrire \all{Die Ware ist beschlagnahmt geworden} pour un état présent est faux : le passif d'état au présent se forme avec \all{sein} + participe II : \all{Die Ware ist beschlagnahmt.} \all{gewesen} n'apparaît qu'au Perfekt : \all{Die Ware ist beschlagnahmt gewesen.}
\item \textbf{Oublier le subjonctif dans le discours indirect.} Après \all{er sagte / erklärte}, l'indicatif simple est une faute de niveau Terminale : écrire \all{..., er habe gelogen} et non \all{er hat gelogen}. Penser au remplacement par le subjonctif II (\all{hätten, würden}) quand le Konj. I ressemble à l'indicatif.
\item \textbf{Négliger majuscules et Umlaute.} \all{falschung} sans majuscule, \all{gefalscht} sans Umlaut : fautes graves. Les formes correctes : \all{die Fälschung, gefälscht, die Märkte}. Tout nom allemand prend la majuscule : \all{der Betrug, die Korruption, der Schmuggel}.
\item \textbf{Faux amis et calques du français.} \all{bekommen} ne signifie pas « devenir » mais « recevoir, obtenir » (\all{Die Firma bekommt den Auftrag.} = l'entreprise obtient le marché) ; « devenir » = \all{werden}. « Contrôler » au sens de « vérifier » = \all{kontrollieren}, et « un contrôle » (policier, douanier) = \all{die Kontrolle}. Ne pas traduire « marché public » par \all{öffentlicher Markt} : on dit \all{der öffentliche Auftrag}.
\item \textbf{Place du verbe et cas.} Dans la subordonnée, le verbe conjugué va en fin : \all{..., weil die Firmen Bestechungsgelder gezahlt haben} (et non \all{haben gezahlt}). Après \all{für, gegen, ohne} : accusatif (\all{gegen die Korruption}) ; après \all{mit, bei, von} : datif (\all{bei der Ausschreibung}). Enfin, \all{jemanden bestechen} exige l'accusatif : \all{Er hat den Beamten bestochen.}
\end{enumerate}
\end{attention}

\newpage

\coursec{Exercices}

\courssub{Je vérifie mes connaissances}

\begin{exercice}[QCM : marchés publics et criminalité économique]
Entoure la bonne réponse.
\begin{enumerate}
\item \all{Die Ausschreibung} signifie :
\begin{enumerate}
\item la falsification \item l'appel d'offres \item la contrebande
\end{enumerate}
\item \all{Der Schmuggel} est :
\begin{enumerate}
\item le passage illégal de marchandises à travers une frontière \item la fabrication de faux billets \item un contrat public
\end{enumerate}
\item Quel article pour \all{Fälschung} ?
\begin{enumerate}
\item der \item die \item das
\end{enumerate}
\item \all{Der Betrug} correspond en français à :
\begin{enumerate}
\item le vol à main armée \item l'escroquerie / la fraude \item l'impôt
\end{enumerate}
\item Lequel de ces verbes introduit typiquement le discours indirect ?
\begin{enumerate}
\item \all{behaupten} \item \all{schlafen} \item \all{backen}
\end{enumerate}
\end{enumerate}
\end{exercice}

\begin{exercice}[Vrai ou faux — justifie ta réponse]
Réponds par \all{richtig} ou \all{falsch} et justifie en une phrase en français.
\begin{enumerate}
\item Le passif d'état se forme avec \all{werden} + participe passé.
\item \all{Die Tür ist geschlossen} décrit un état, pas une action.
\item Au discours indirect, on utilise normalement le subjonctif II en premier choix.
\item \all{Er sagte, er habe die Ware nicht gefälscht} : \all{habe} est bien au subjonctif I.
\item Le participe passé de \all{stehlen} est \all{gestohlen}.
\end{enumerate}
\end{exercice}

\begin{exercice}[Vocabulaire : article, pluriel, famille de mots]
Complète le tableau suivant (article, pluriel, traduction française, un mot de la même famille).

\begin{center}
\begin{tabularx}{\linewidth}{|l|l|l|X|X|}
\hline
\rowcolor{popblueL} Mot & Article & Pluriel & Traduction & Famille de mots \\
\hline
Markt & & & & \\
\hline
Korruption & & & & \\
\hline
Schmuggel & & & & \\
\hline
Fälschung & & & & \\
\hline
Ausschreibung & & & & \\
\hline
Betrug & & & & \\
\hline
Ware & & & & \\
\hline
Grenze & & & & \\
\hline
\end{tabularx}
\end{center}
\end{exercice}

\begin{exercice}[Conjugaison : subjonctif I à compléter]
Complète avec la forme correcte du subjonctif I du verbe entre parenthèses.
\begin{enumerate}
\item Der Minister erklärte, er \_\_\_\_\_\_\_\_ (sein) gegen jede Form der Korruption.
\item Die Firma behauptete, sie \_\_\_\_\_\_\_\_ (haben, au Perfekt) alle Steuern bezahlt.
\item Der Händler gab zu, dass die Medikamente gefälscht \_\_\_\_\_\_\_\_ (sein).
\item Die Zeugin sagte, sie \_\_\_\_\_\_\_\_ (kennen) den Schmuggler nicht.
\item Der Zöllner berichtete, die Ware \_\_\_\_\_\_\_\_ (werden, Passiv) an der Grenze beschlagnahmt.
\item Er meinte, man \_\_\_\_\_\_\_\_ (müssen) die Kontrollen verstärken.
\end{enumerate}
\end{exercice}

\courssub{Je m'entraîne}

\begin{exercice}[Thème : traduis en allemand]
Traduis les phrases suivantes en allemand en utilisant le discours indirect (subjonctif I) ou le passif d'état selon le cas.
\begin{enumerate}
\item Le ministre a déclaré que les produits étaient falsifiés.
\item Le commerçant a affirmé qu'il n'avait jamais vendu de marchandises de contrebande.
\item La porte du bureau des douanes est fermée depuis ce matin.
\item Le journaliste a écrit que l'appel d'offres avait été truqué.
\item Les médicaments sont détruits par les douaniers ; maintenant le stock est détruit.
\end{enumerate}
\end{exercice}

\begin{exercice}[Transformation : du discours direct au discours indirect]
Mets les phrases suivantes au discours indirect en commençant par le verbe introducteur indiqué. Utilise le subjonctif I.
\begin{enumerate}
\item \all{„Ich habe die gefälschten Produkte nicht verkauft.“} (Der Händler behauptet, ...)
\item \all{„Die Ausschreibung war nicht fair.“} (Der Unternehmer erklärt, ...)
\item \all{„Wir werden die Schmuggler verhaften.“} (Die Polizei kündigt an, ...)
\item \all{„Ich habe das Geld an der Grenze versteckt.“} (Der Angeklagte gibt zu, ...)
\item \all{„Die Kontrollen sind zu schwach.“} (Der Experte kritisiert, ...)
\item \all{„Man muss die Konsumenten besser informieren.“} (Die Ministerin fordert, ...)
\end{enumerate}
\end{exercice}

\begin{exercice}[Texte à trous : Vorgangspassiv ou Zustandspassiv ?]
Complète le paragraphe avec la forme correcte de \all{sein} ou \all{werden} + le participe passé du verbe entre parenthèses. Attention : choisis le passif d'action ou le passif d'état selon le sens.

\begin{textebox}[An der Grenze]
Gestern Abend \_\_\_\_\_\_\_\_ ein Lastwagen an der Grenze \_\_\_\_\_\_\_\_ (kontrollieren). Dabei \_\_\_\_\_\_\_\_ gefälschte Medikamente \_\_\_\_\_\_\_\_ (finden). Der Fahrer \_\_\_\_\_\_\_\_ sofort \_\_\_\_\_\_\_\_ (verhaften). Heute Morgen ist der Lkw immer noch nicht \_\_\_\_\_\_\_\_ (abfertigen) : die Papiere sind noch nicht \_\_\_\_\_\_\_\_ (prüfen). Die Ware bleibt \_\_\_\_\_\_\_\_ (beschlagnahmen), bis die Untersuchung \_\_\_\_\_\_\_\_ (abschließen) ist.
\end{textebox}
\end{exercice}

\begin{exercice}[Appariements : mots et définitions]
Relie chaque mot allemand (1--8) à sa définition (a--h).

\begin{center}
\begin{tabularx}{\linewidth}{|l|X|}
\hline
\rowcolor{popblueL} Mot & Définition \\
\hline
1. die Korruption & a. Waren illegal über die Grenze bringen \\
\hline
2. der Schmuggel & b. ein öffentlicher Wettbewerb für einen Auftrag \\
\hline
3. die Fälschung & c. ein Beamter nimmt Geld, um eine Entscheidung zu ändern \\
\hline
4. die Ausschreibung & d. ein Produkt, das nicht echt ist \\
\hline
5. der Betrug & e. jemanden bewusst täuschen, um Geld zu bekommen \\
\hline
6. der Zoll & f. die Kontrolle der Waren an der Grenze \\
\hline
7. beschlagnahmen & g. der Staat nimmt illegale Waren weg \\
\hline
8. der Konsument & h. die Person, die ein Produkt kauft und benutzt \\
\hline
\end{tabularx}
\end{center}
\end{exercice}

\courssub{Je me prépare au Bac}

\begin{exercice}[Type Bac — Compréhension écrite]
Lis le texte suivant, puis réponds aux questions en allemand, avec des phrases complètes.

\begin{textebox}[Gefälschte Medikamente : ein Skandal erschüttert die Region]
In einer großen Stadt Zentralafrikas haben Zollbeamte am Dienstag mehrere Tonnen gefälschter Medikamente beschlagnahmt. Die Ware war in Containern versteckt, die offiziell „Landmaschinen“ transportieren sollten. Der Leiter der Zollbehörde erklärte, die Fälschungen seien so gut gemacht gewesen, dass selbst Apotheker sie kaum erkannt hätten. Ein Händler, der anonym bleiben will, gab zu, er habe die Produkte seit Monaten auf dem Markt verkauft. Er behauptete jedoch, er habe nicht gewusst, dass die Medikamente gefälscht gewesen seien. Die Gesundheitsministerin kündigte an, man werde die Kontrollen an allen Grenzen verstärken und die Verantwortlichen vor Gericht bringen. Experten warnen : gefälschte Medikamente seien besonders gefährlich, weil sie oft keine Wirkstoffe enthielten oder sogar giftig seien. Die Bevölkerung wird aufgefordert, Medikamente nur in offiziellen Apotheken zu kaufen.
\end{textebox}

\textbf{Questions :}
\begin{enumerate}
\item Richtig oder falsch ? Justifie mit einer Textstelle :
\begin{enumerate}
\item Die Medikamente waren in Containern versteckt.
\item Die Fälschungen waren leicht zu erkennen.
\item Der Händler wusste von Anfang an alles.
\end{enumerate}
\item Was sollte in den Containern offiziell transportiert werden ?
\item Was hat die Gesundheitsministerin angekündigt ? (2 Punkte)
\item Warum sind gefälschte Medikamente laut den Experten gefährlich ? (2 Gründe)
\item Was empfiehlt man der Bevölkerung ?
\item Reformuliere mit deinen eigenen Worten : \all{„Die Ware wurde beschlagnahmt.“}
\end{enumerate}
\end{exercice}

\begin{exercice}[Type Bac — Grammaire : discours indirect et passif]
\textbf{A.} Mets les phrases suivantes au discours indirect (subjonctif I) avec le verbe introducteur indiqué.
\begin{enumerate}
\item \all{„Die Produkte sind falsch etikettiert.“} (Der Inspektor erklärt, ...)
\item \all{„Ich habe Bestechungsgeld angenommen.“} (Der Beamte gibt zu, ...)
\item \all{„Wir haben den Auftrag ehrlich bekommen.“} (Die Firma behauptet, ...)
\item \all{„Der Schmuggel wird härter bestraft.“} (Der Minister verspricht, ...)
\item \all{„Man hat die Grenze nachts überquert.“} (Der Zeuge berichtet, ...)
\item \all{„Die Kontrollen waren nicht streng genug.“} (Der Bericht stellt fest, ...)
\end{enumerate}

\textbf{B.} Complète avec \all{sein} ou \all{werden} (forme correcte) + participe passé, et justifie ton choix (action ou état).
\begin{enumerate}
\item Die gefälschten Waren \_\_\_\_\_\_\_\_ gestern \_\_\_\_\_\_\_\_ (zerstören).
\item Jetzt \_\_\_\_\_\_\_\_ der Markt endlich \_\_\_\_\_\_\_\_ (säubern) : es gibt keine Fälschungen mehr.
\item Der Verdächtige \_\_\_\_\_\_\_\_ noch nicht \_\_\_\_\_\_\_\_ (verurteilen) ; der Prozess dauert an.
\end{enumerate}
\end{exercice}

\begin{exercice}[Type Bac — Thème, version et expression écrite]
\textbf{A. Thème.} Traduis en allemand (attention au subjonctif I et au passif d'état).
\begin{enumerate}
\item Le directeur a affirmé que l'appel d'offres avait été transparent.
\item Les douaniers ont déclaré que la marchandise était falsifiée.
\item Depuis hier, le bureau est fermé et les dossiers sont bloqués.
\item Le ministre a promis que les coupables seraient punis.
\item On dit que la contrebande détruit l'économie du pays.
\end{enumerate}

\textbf{B. Version.} Traduis en français l'extrait suivant.

\begin{textebox}[Aus einem Zeitungsartikel]
Die Staatsanwaltschaft teilte mit, dass drei Verdächtige festgenommen worden seien. Sie hätten über zwei Jahre gefälschte Medikamente in die Region geschmuggelt. Ein Verdächtiger gab zu, er habe Beamte bestochen, damit die Container nicht kontrolliert würden. Die Ermittler erklärten, der Schaden für die öffentliche Gesundheit sei enorm. Die Regierung kündigte an, sie werde ein neues Gesetz gegen Produktfälschung vorschlagen.
\end{textebox}

\textbf{C. Expression écrite.} Rédige un commentaire d'environ 15 lignes (environ 120--150 mots) sur le sujet suivant :

\begin{center}
\all{„Wie kann man die Fälschung von Medikamenten in Afrika bekämpfen ?“}
\end{center}

Respecte le plan imposé :
\begin{enumerate}
\item \textbf{Einleitung} : présente le problème (2--3 lignes).
\item \textbf{Hauptteil} : propose trois solutions concrètes (contrôles, justice, information des consommateurs) avec des exemples.
\item \textbf{Schluss} : donne ton avis personnel.
\end{enumerate}
Utilise au moins deux fois le discours indirect et une fois le passif d'état.
\end{exercice}

\newpage

\coursec{Corrigés des exercices}

\begin{corrige}[Exercice 1 — QCM : marchés publics et criminalité économique]
\begin{enumerate}
\item \textbf{b} : \all{die Ausschreibung} = l'appel d'offres (procédure par laquelle l'État demande des offres pour un marché public). La falsification = \all{die Fälschung} ; la contrebande = \all{der Schmuggel}.
\item \textbf{a} : \all{der Schmuggel} = le passage illégal de marchandises à travers une frontière. La fabrication de faux billets relève de \all{die Fälschung} (faux-monnayage : \all{das Falschgeld}).
\item \textbf{b} : \all{die Fälschung}. Tous les noms en \all{-ung} sont féminins (comme \all{-ung} $\rightarrow$ \all{die}) : \all{die Ausschreibung, die Korruption, die Untersuchung}.
\item \textbf{b} : \all{der Betrug} = l'escroquerie, la fraude. Le vol à main armée = \all{der bewaffnete Raub} ; l'impôt = \all{die Steuer}.
\item \textbf{a} : \all{behaupten} (affirmer, prétendre) est un verbe de déclaration qui introduit typiquement le discours indirect au subjonctif I : \all{Er behauptet, er sei unschuldig.} \all{schlafen} et \all{backen} n'introduisent pas de parole rapportée.
\end{enumerate}
\end{corrige}

\begin{corrige}[Exercice 2 — Vrai ou faux]
\begin{enumerate}
\item \all{falsch}. Le passif d'état (\all{Zustandspassiv}) se forme avec \all{sein} + participe passé (\all{Die Tür ist geschlossen}). C'est le passif d'action (\all{Vorgangspassiv}) qui se forme avec \all{werden} + participe passé (\all{Die Tür wird geschlossen}).
\item \all{richtig}. \all{Die Tür ist geschlossen} décrit le résultat d'une action achevée, c'est-à-dire un état (la porte est dans l'état « fermée »), et non l'action de fermer en train de se dérouler.
\item \all{falsch}. Au discours indirect, le premier choix est le \textbf{subjonctif I} (\all{er habe, sie sei}). On ne recourt au subjonctif II que lorsque la forme du subjonctif I est identique à l'indicatif (ex. \all{sie haben} $\rightarrow$ \all{sie hätten}).
\item \all{richtig}. \all{habe} (3\textsuperscript{e} pers. sg. de \all{haben}) est la forme de subjonctif I ; l'indicatif serait \all{hat}. \all{Er sagte, er habe die Ware nicht gefälscht} = « Il a dit qu'il n'avait pas falsifié la marchandise » (Perfekt au discours indirect : \all{habe} + participe).
\item \all{richtig}. \all{stehlen -- stahl -- hat gestohlen} : verbe fort à participe en \all{ge-...-en} avec changement de voyelle (\all{ie} $\rightarrow$ \all{o}). Le participe passé est bien \all{gestohlen}.
\end{enumerate}
\end{corrige}

\begin{corrige}[Exercice 3 — Vocabulaire : article, pluriel, famille de mots]
\begin{center}
\begin{tabularx}{\linewidth}{|l|l|l|X|X|}
\hline
\rowcolor{popblueL} Mot & Article & Pluriel & Traduction & Famille de mots \\
\hline
Markt & der & die Märkte & le marché & der Marktplatz, marktfähig \\
\hline
Korruption & die & (die Korruptionen, rare) & la corruption & korrupt, korrumpieren \\
\hline
Schmuggel & der & (sans pluriel) & la contrebande & schmuggeln, der Schmuggler \\
\hline
Fälschung & die & die Fälschungen & la falsification & fälschen, falsch, der Fälscher \\
\hline
Ausschreibung & die & die Ausschreibungen & l'appel d'offres & ausschreiben, die Schreibung \\
\hline
Betrug & der & (die Betrüge, rare) & l'escroquerie, la fraude & betrügen, der Betrüger \\
\hline
Ware & die & die Waren & la marchandise & das Warenhaus, der Warenverkehr \\
\hline
Grenze & die & die Grenzen & la frontière & der Grenzverkehr, grenzenlos \\
\hline
\end{tabularx}
\end{center}
\textbf{Points de vigilance :} \all{der Markt} prend l'Umlaut au pluriel (\all{die Märkte}) ; \all{der Schmuggel} et \all{die Korruption} sont des noms abstraits, généralement sans pluriel ; \all{falsch} (faux, adjectif) est un faux ami partiel : il ne signifie pas « falsifié » au sens de contrefait, on préfère \all{gefälscht}.
\end{corrige}

\begin{corrige}[Exercice 4 — Conjugaison : subjonctif I]
\begin{enumerate}
\item Der Minister erklärte, er \textbf{sei} gegen jede Form der Korruption. (subj. I de \all{sein}, 3\textsuperscript{e} pers. sg. : \all{sei} — forme irrégulière à connaître par cœur.)
\item Die Firma behauptete, sie \textbf{habe} alle Steuern \textbf{bezahlt}. (Perfekt au discours indirect : subj. I de \all{haben} + participe passé ; \all{bezahlt} sans \all{ge-} car verbe en \all{-ieren} ? Non : \all{bezahlen} prend \all{be-} préfixe inséparable, donc pas de \all{ge-}.)
\item Der Händler gab zu, dass die Medikamente gefälscht \textbf{seien}. (subj. I de \all{sein}, 3\textsuperscript{e} pers. pl. : \all{seien} ; passif d'état rapporté.)
\item Die Zeugin sagte, sie \textbf{kenne} den Schmuggler nicht. (subj. I de \all{kennen}, 3\textsuperscript{e} pers. sg. : radical + \all{-e}.)
\item Der Zöllner berichtete, die Ware \textbf{werde} an der Grenze beschlagnahmt. (passif d'action au discours indirect : subj. I de \all{werden} + participe passé.)
\item Er meinte, man \textbf{müsse} die Kontrollen verstärken. (subj. I de \all{müssen}, 3\textsuperscript{e} pers. sg. : \all{müsse}, avec Umlaut conservé ; ne pas confondre avec le subj. II \all{müsste}.)
\end{enumerate}
\textbf{Rappel de formation :} subjonctif I = radical de l'infinitif + \all{-e, -est, -e, -en, -et, -en} ; seul \all{sein} est irrégulier : \all{ich sei, du seiest/seist, er sei, wir seien, ihr seiet, sie seien}.
\end{corrige}

\begin{corrige}[Exercice 5 — Thème : traduis en allemand]
\begin{enumerate}
\item \all{Der Minister erklärte, die Produkte seien gefälscht.} (ou : \all{..., dass die Produkte gefälscht seien.}) — Passif d'état rapporté : \all{seien} + participe passé.
\item \all{Der Händler behauptete, er habe nie Schmuggelware verkauft.} — Perfekt au discours indirect : \all{habe} + \all{verkauft} ; « marchandises de contrebande » = \all{die Schmuggelware}.
\item \all{Die Tür des Zollamtes ist seit heute Morgen geschlossen.} — Passif d'état au présent : \all{ist} + \all{geschlossen} ; « depuis ce matin » = \all{seit heute Morgen} (+ datif).
\item \all{Der Journalist schrieb, die Ausschreibung sei manipuliert worden.} (ou : \all{... sei nicht fair gewesen.}) — Plus-que-parfait du passif au discours indirect : \all{sei} + participe + \all{worden}.
\item \all{Die Medikamente werden von den Zöllnern zerstört ; jetzt ist der Bestand zerstört.} — Première phrase : passif d'action (\all{werden zerstört}) ; deuxième : passif d'état (\all{ist zerstört}), le résultat. « Le stock » = \all{der Bestand}.
\end{enumerate}
\textbf{Points de vigilance :} ne jamais écrire \all{*ist zerstört worden} pour un état présent ; \all{worden} marque le passé du passif d'action. « Les douaniers » = \all{die Zöllner} (ou \all{die Zollbeamten}).
\end{corrige}

\begin{corrige}[Exercice 6 — Transformation : du discours direct au discours indirect]
\begin{enumerate}
\item Der Händler behauptet, er \textbf{habe} die gefälschten Produkte nicht \textbf{verkauft}. (Perfekt $\rightarrow$ \all{habe} + participe.)
\item Der Unternehmer erklärt, die Ausschreibung \textbf{sei} nicht fair \textbf{gewesen}. (Präteritum \all{war} $\rightarrow$ Perfekt au discours indirect : \all{sei gewesen}.)
\item Die Polizei kündigt an, sie \textbf{werde} die Schmuggler verhaften. (Futur \all{werden} $\rightarrow$ subj. I : \all{werde} + infinitif.)
\item Der Angeklagte gibt zu, er \textbf{habe} das Geld an der Grenze \textbf{versteckt}. (Perfekt $\rightarrow$ \all{habe versteckt}.)
\item Der Experte kritisiert, die Kontrollen \textbf{seien} zu schwach. (présent \all{sind} $\rightarrow$ subj. I pl. \all{seien}.)
\item Die Ministerin fordert, man \textbf{müsse} die Konsumenten besser informieren. (présent \all{muss} $\rightarrow$ subj. I \all{müsse}.)
\end{enumerate}
\textbf{Méthode :} 1) repérer le temps de la parole directe ; 2) présent $\rightarrow$ subj. I du présent ; passé (Perfekt/Präteritum) $\rightarrow$ subj. I de \all{haben/sein} + participe ; futur $\rightarrow$ \all{werde} + infinitif ; 3) adapter les pronoms (\all{ich} $\rightarrow$ \all{er/sie}, \all{wir} $\rightarrow$ \all{sie}).
\end{corrige}

\begin{corrige}[Exercice 7 — Texte à trous : Vorgangspassiv ou Zustandspassiv ?]
Gestern Abend \textbf{wurde} ein Lastwagen an der Grenze \textbf{kontrolliert}. Dabei \textbf{wurden} gefälschte Medikamente \textbf{gefunden}. Der Fahrer \textbf{wurde} sofort \textbf{verhaftet}. Heute Morgen ist der Lkw immer noch nicht \textbf{abgefertigt} : die Papiere sind noch nicht \textbf{geprüft}. Die Ware bleibt \textbf{beschlagnahmt}, bis die Untersuchung \textbf{abgeschlossen} ist.

\textbf{Justifications :}
\begin{itemize}
\item \all{wurde kontrolliert / wurden gefunden / wurde verhaftet} : passif d'action au Präteritum (\all{werden} $\rightarrow$ \all{wurde(n)}), car il s'agit d'événements ponctuels (« gestern Abend »). Attention à l'accord : \all{Medikamente} pluriel $\rightarrow$ \all{wurden}.
\item \all{ist abgefertigt, sind geprüft, bleibt beschlagnahmt, ist abgeschlossen} : passif d'état au présent (\all{sein} + participe), car on décrit le résultat actuel (« heute Morgen », « immer noch nicht »). Après \all{bleiben}, le participe passé a aussi valeur d'état.
\end{itemize}
\textbf{Astuce :} un marqueur de moment ponctuel (\all{gestern, plötzlich}) appelle \all{werden} ; un marqueur de durée/résultat (\all{seit, immer noch, jetzt}) appelle \all{sein}.
\end{corrige}

\begin{corrige}[Exercice 8 — Appariements : mots et définitions]
\begin{center}
\begin{tabular}{|c|c|l|}
\hline
\rowcolor{popblueL} Mot & Définition & Traduction de la définition \\
\hline
1. die Korruption & c & un fonctionnaire prend de l'argent pour modifier une décision \\
\hline
2. der Schmuggel & a & faire passer illégalement des marchandises à la frontière \\
\hline
3. die Fälschung & d & un produit qui n'est pas authentique \\
\hline
4. die Ausschreibung & b & une mise en concurrence publique pour un marché \\
\hline
5. der Betrug & e & tromper volontairement quelqu'un pour obtenir de l'argent \\
\hline
6. der Zoll & f & le contrôle des marchandises à la frontière \\
\hline
7. beschlagnahmen & g & l'État confisque les marchandises illégales \\
\hline
8. der Konsument & h & la personne qui achète et utilise un produit \\
\hline
\end{tabular}
\end{center}
\textbf{Correspondances :} 1--c, 2--a, 3--d, 4--b, 5--e, 6--f, 7--g, 8--h.
\end{corrige}

\begin{corrige}[Exercice 9 — Type Bac : compréhension écrite]
\textbf{Barème indicatif (10 points) :} question 1 : 3 $\times$ 1 pt (réponse + justification) ; questions 2 à 5 : 1 pt chacune ; question 6 : 2 pts ; qualité de la langue : 1 pt.

\textbf{Réponses modèles :}
\begin{enumerate}
\item \begin{enumerate}
\item \all{Richtig. Im Text steht : „Die Ware war in Containern versteckt.“}
\item \all{Falsch. Der Leiter der Zollbehörde erklärte, die Fälschungen seien so gut gemacht gewesen, dass selbst Apotheker sie kaum erkannt hätten.} (Elles étaient donc très difficiles à reconnaître.)
\item \all{Falsch. Er behauptete, er habe nicht gewusst, dass die Medikamente gefälscht gewesen seien.}
\end{enumerate}
\item \all{Offiziell sollten in den Containern Landmaschinen transportiert werden.} (Des machines agricoles.)
\item \all{Sie kündigte an, man werde die Kontrollen an allen Grenzen verstärken und die Verantwortlichen vor Gericht bringen.} (Renforcer les contrôles + traduire les responsables en justice.)
\item \all{Gefälschte Medikamente sind gefährlich, weil sie oft keine Wirkstoffe enthalten oder sogar giftig sind.} (Pas de principe actif ; parfois toxiques.)
\item \all{Man empfiehlt der Bevölkerung, Medikamente nur in offiziellen Apotheken zu kaufen.}
\item Reformulation possible : \all{Die Zollbeamten haben die Ware weggenommen / konfisziert.} ou \all{Der Staat hat die illegalen Waren eingezogen.}
\end{enumerate}
\textbf{Points de vigilance :} répondre par des phrases complètes ; reprendre le discours indirect du texte sans le transformer en indicatif ; \all{beschlagnahmen} = confisquer/saisir (verbe faible, participe \all{beschlagnahmt}).
\end{corrige}

\begin{corrige}[Exercice 10 — Type Bac : grammaire]
\textbf{Barème indicatif :} A : 6 $\times$ 1 pt ; B : 3 $\times$ 2 pts (forme + justification).

\textbf{A. Discours indirect (subjonctif I) :}
\begin{enumerate}
\item Der Inspektor erklärt, die Produkte \textbf{seien} falsch etikettiert. (\all{sind} $\rightarrow$ \all{seien}.)
\item Der Beamte gibt zu, er \textbf{habe} Bestechungsgeld \textbf{angenommen}. (Perfekt : \all{habe} + participe ; \all{annehmen} $\rightarrow$ \all{angenommen}, verbe fort à particule séparable.)
\item Die Firma behauptet, sie \textbf{habe} den Auftrag ehrlich \textbf{bekommen}. (\all{bekommen} : faux ami = « recevoir, obtenir », pas « devenir » !)
\item Der Minister verspricht, der Schmuggel \textbf{werde} härter bestraft. (Futur du passif rapporté : \all{werde} + participe.)
\item Der Zeuge berichtet, man \textbf{habe} die Grenze nachts \textbf{überquert}. (Perfekt avec \all{haben} car \all{überqueren} est transitif ici.)
\item Der Bericht stellt fest, die Kontrollen \textbf{seien} nicht streng genug \textbf{gewesen}. (Präteritum \all{waren} $\rightarrow$ \all{seien gewesen}.)
\end{enumerate}

\textbf{B. Passif d'action ou d'état :}
\begin{enumerate}
\item Die gefälschten Waren \textbf{wurden} gestern \textbf{zerstört}. — Action ponctuelle (« gestern ») $\rightarrow$ Vorgangspassiv au Präteritum.
\item Jetzt \textbf{ist} der Markt endlich \textbf{gesäubert} : es gibt keine Fälschungen mehr. — Résultat présent (« jetzt », « endlich ») $\rightarrow$ Zustandspassiv au présent.
\item Der Verdächtige \textbf{ist} noch nicht \textbf{verurteilt} ; der Prozess dauert an. — État actuel (« noch nicht », le procès continue) $\rightarrow$ Zustandspassiv. (\all{wird noch nicht verurteilt} serait possible mais décrirait l'action ; le sens demandé est l'état.)
\end{enumerate}
\end{corrige}

\begin{corrige}[Exercice 11 — Type Bac : thème, version et expression écrite]
\textbf{Barème indicatif :} A : 5 $\times$ 2 pts ; B : 5 pts ; C : 10 pts (contenu 4, langue 4, respect des contraintes 2).

\textbf{A. Thème — propositions de traduction :}
\begin{enumerate}
\item \all{Der Direktor behauptete, die Ausschreibung sei transparent gewesen.} (plus-que-passé $\rightarrow$ \all{sei} + participe + \all{gewesen}.)
\item \all{Die Zöllner erklärten, die Ware sei gefälscht.} (passif d'état rapporté : \all{sei} + participe.)
\item \all{Seit gestern ist das Büro geschlossen und die Akten sind blockiert.} (passif d'état au présent ; « les dossiers » = \all{die Akten}.)
\item \all{Der Minister versprach, die Schuldigen würden bestraft.} (futur rapporté : la forme de subj. I \all{werden} étant identique à l'indicatif au pluriel, on utilise le subj. II \all{würden} — cas obligatoire de substitution.)
\item \all{Man sagt, der Schmuggel zerstöre die Wirtschaft des Landes.} (ou : \all{..., dass der Schmuggel die Wirtschaft des Landes zerstöre.})
\end{enumerate}

\textbf{B. Version — traduction modèle :}

« Le parquet a fait savoir que trois suspects avaient été arrêtés. Ils auraient, pendant deux ans, introduit en contrebande des médicaments falsifiés dans la région. Un suspect a avoué qu'il avait soudoyé des fonctionnaires pour que les conteneurs ne soient pas contrôlés. Les enquêteurs ont déclaré que le préjudice pour la santé publique était énorme. Le gouvernement a annoncé qu'il proposerait une nouvelle loi contre la contrefaçon de produits. »

\textbf{Points de vigilance (version) :} \all{die Staatsanwaltschaft} = le parquet (pas « l'avocat de l'État ») ; \all{festnehmen} = arrêter ; \all{bestechen -- bestach -- bestochen} = soudoyer/corrompre ; \all{der Schaden} = le préjudice/dommage ; rendre le subjonctif I par le conditionnel journalistique (« aurait introduit ») ou l'indicatif selon le degré de certitude.

\textbf{C. Expression écrite — production modèle :}

\begin{textebox}[Modèle de commentaire]
Gefälschte Medikamente sind ein großes Problem in vielen afrikanischen Ländern. Sie werden auf den Märkten und sogar in manchen Apotheken verkauft, und viele Patienten sterben, weil die Produkte keine Wirkstoffe enthalten. Wie kann man diese Kriminalität bekämpfen ?

Zuerst muss der Staat die Kontrollen an den Grenzen verstärken. Experten sagen, die Zollbeamten seien oft zu wenig ausgebildet und schlecht bezahlt. Wenn ein Beamter korrupt ist, bleibt die Grenze praktisch offen. Zweitens müssen die Gerichte härter arbeiten : ein Schmuggler, der nie bestraft wird, fängt immer wieder an. Die Minister erklären oft, die Verantwortlichen würden vor Gericht gebracht, aber das muss auch wirklich passieren. Drittens muss man die Konsumenten informieren. Kampagnen im Radio und in den Schulen können erklären, dass man Medikamente nur in offiziellen Apotheken kaufen soll.

Meiner Meinung nach ist der Kampf gegen die Fälschung möglich, aber nur, wenn Staat, Justiz und Bevölkerung zusammenarbeiten. Ein falsifiziertes Medikament ist kein kleines Verbrechen : es ist eine Waffe gegen die Gesundheit aller.
\end{textebox}

\textbf{Contraintes respectées :} discours indirect : \all{Experten sagen, die Zollbeamten seien...} ; \all{Die Minister erklären oft, die Verantwortlichen würden...} ; passif d'état : \all{bleibt die Grenze praktisch offen} et \all{ein Schmuggler, der nie bestraft wird} (passif d'action) — on peut aussi citer \all{ist... verkauft} en début de texte.

\textbf{Points de vigilance :} majuscules à tous les noms ; verbe en deuxième position ; \all{die Wirkstoffe} (principes actifs) ; éviter le calque du français « *informieren über » sans préposition correcte : \all{jemanden über etwas informieren}.
\end{corrige}

\newpage

\coursec{Fiche bilan}

\begin{popfigure}
\begin{tikzpicture}[mindmap, grow cyclic, every node/.style={concept, align=center, font=\small},
  root concept/.append style={concept color=popblue, font=\small\bfseries, text=white},
  level 1 concept/.append style={level distance=3.4cm, sibling angle=60, font=\footnotesize},
  level 2 concept/.append style={level distance=2.2cm, font=\scriptsize}]
\node[root concept, align=center]{Die öffentlichen Märkte\\et la criminalité économique}
  child[concept color=poporange]{ node[align=center]{Lexique\\thématique}
    child[concept color=poporangeL]{ node[align=center]{der Markt\\die Ausschreibung} }
    child[concept color=poporangeL]{ node[align=center]{die Korruption\\der Betrug} }
  }
  child[concept color=popgreen]{ node[align=center]{Zustandspassiv\\(passif d'état)}
    child[concept color=popgreenL]{ node[align=center]{sein + Partizip II\\Das Tor ist geöffnet.} }
  }
  child[concept color=poppurple]{ node[align=center]{Discours indirect\\Konjunktiv I}
    child[concept color=poppurpleL]{ node {Er sagte, er habe\ldots} }
    child[concept color=poppurpleL]{ node[align=center]{KI $\rightarrow$ KII\\si identique} }
  }
  child[concept color=poppink]{ node[align=center]{Conjugaison\\KI et KII}
    child[concept color=poppinkL]{ node[align=center]{habe, sei, werde\\hätte, wäre, würde} }
  }
  child[concept color=popgold]{ node[align=center]{Méthode\\commentaire}
    child[concept color=popgoldL]{ node[align=center]{lire -- relever\\analyser -- juger} }
  }
  child[concept color=popblue!60]{ node[align=center]{Production\\guidée}
    child[concept color=popblueL]{ node[align=center]{résumé, débat\\court paragraphe} }
  };
\end{tikzpicture}
\legende{Carte mentale de la leçon : marchés publics et criminalité économique}
\end{popfigure}

\begin{bilan}
\textbf{1. Lexique clé (à connaître avec article et pluriel)}
\begin{itemize}
  \item \all{der Markt, die Märkte} : le marché ; \all{der öffentliche Auftrag, die Aufträge} : le marché public ; \all{die Ausschreibung, -en} : l'appel d'offres ; \all{der Auftragnehmer, -} : l'attributaire.
  \item \all{die Korruption} (sans pluriel) : la corruption ; \all{die Bestechung, -en} : la corruption/le pot-de-vin ; \all{bestechen, besticht, bestach, hat bestochen} : corrompre ; \all{korrupt} : corrompu.
  \item \all{der Schmuggel} : la contrebande ; \all{schmuggeln, hat geschmuggelt} : faire passer en contrebande ; \all{der Schmuggler, -} : le contrebandier.
  \item \all{die Fälschung, -en} : la falsification/le faux ; \all{fälschen, hat gefälscht} : falsifier ; \all{gefälschte Medikamente} : des médicaments falsifiés.
  \item \all{der Betrug} : l'escroquerie ; \all{betrügen, betrügt, betrog, hat betrogen} : tromper, escroquer ; \all{der Betrüger, -} : l'escroc.
  \item \all{bekämpfen, hat bekämpft} : combattre ; \all{der Zoll, die Zölle} : la douane ; \all{beschlagnahmen, hat beschlagnahmt} : saisir ; \all{die Strafe, -n} : la peine, la sanction.
\end{itemize}

\textbf{2. Grammaire à connaître par cœur}
\begin{center}
\begin{tabularx}{\linewidth}{|l|X|X|}
\hline
\rowcolor{popblueL} \textbf{Point} & \textbf{Règle} & \textbf{Exemple} \\
\hline
Zustandspassiv & \cle{sein} + Partizip II ; décrit un \textbf{état résultant} d'une action (pas l'action elle-même) & \all{Die Waren sind beschlagnahmt.} « Les marchandises sont saisies (état). » \\
\hline
Vorgangspassiv & \cle{werden} + Partizip II ; décrit l'\textbf{action} & \all{Die Waren werden beschlagnahmt.} « On saisit les marchandises. » \\
\hline
Konjunktiv I & formé sur l'\textbf{infinitif} + désinences (-e, -est, -e, -en, -et, -en) ; sert au \textbf{discours indirect} & \all{Er sagte, die Firma habe betrogen.} « Il a dit que l'entreprise avait escroqué. » \\
\hline
Konjunktiv II & formé sur le \textbf{Präteritum} (souvent avec Umlaut) ; hypothèse, irréel ; remplace le KI identique à l'indicatif & \all{Er sagte, sie hätten geschmuggelt.} (au lieu de \all{sie haben}) \\
\hline
sein au KI & ich \all{sei}, du \all{sei(e)st}, er \all{sei}, wir \all{seien}, ihr \all{sei(e)t}, sie \all{seien} & \all{Er behauptet, er sei unschuldig.} \\
\hline
\end{tabularx}
\end{center}

\textbf{3. Méthodes express}
\begin{itemize}
  \item \textbf{Comprendre le texte} : 1) lire le titre et la source ; 2) lecture globale (qui ? quoi ? où ?) ; 3) lecture détaillée avec le glossaire ; 4) repérer les chiffres et les exemples ; 5) reformuler en une phrase.
  \item \textbf{Discours indirect} : 1) repérer le verbe introducteur (\all{sagen, behaupten, erklären, berichten}) ; 2) passer le verbe au KI ; 3) si la forme KI = indicatif, utiliser le KII ; 4) adapter les pronoms et les possessifs.
  \item \textbf{Passif d'état} : demander « action ou résultat ? » ; résultat $\rightarrow$ \all{sein} + Partizip II ; action $\rightarrow$ \all{werden} + Partizip II.
  \item \textbf{Production} : 4 à 6 phrases, connecteurs (\all{zuerst, außerdem, deshalb, schließlich}), un exemple camerounais (marché de Yaoundé, frontière, douane de Douala), une conclusion personnelle.
\end{itemize}
\end{bilan}

\begin{aretenir}[Checklist avant le Bac]
\begin{itemize}
  \item $\square$ Je connais l'article et le pluriel de \all{der Markt, die Ausschreibung, die Fälschung, der Betrug, der Zoll}.
  \item $\square$ Je sais conjuguer \all{bestechen, betrügen, schmuggeln, fälschen} au Präsens et au Perfekt.
  \item $\square$ Je distingue \all{werden} + Partizip II (action) et \all{sein} + Partizip II (état).
  \item $\square$ Je sais former le Konjunktiv I à partir de l'infinitif (\all{er habe, er sei, er komme}).
  \item $\square$ Je sais former le Konjunktiv II à partir du Präteritum (\all{er hätte, er wäre, er käme}).
  \item $\square$ Je remplace le KI par le KII quand la forme est identique à l'indicatif (\all{sie haben} $\rightarrow$ \all{sie hätten}).
  \item $\square$ Je sais rapporter une parole au discours indirect en adaptant les pronoms.
  \item $\square$ Je mets la majuscule à tous les noms allemands et j'utilise les guillemets „\ldots“.
  \item $\square$ Je place le verbe conjugué en deuxième position et le participe en fin de phrase.
  \item $\square$ Je sais rédiger un court paragraphe argumenté sur la corruption ou la contrebande avec des connecteurs.
  \item $\square$ Je connais les faux amis : \all{der Betrug} = escroquerie (pas « bétruque »), \all{die Strafe} = peine/sanction.
  \item $\square$ Je relis ma production : accords, cas, place du verbe, Umlaute (ä, ö, ü) et ß.
\end{itemize}
\end{aretenir}

\findecours

\end{document}
